% ============================================================
%  Harald Høffding, "Begrebet Analogi"
%  "The Concept of Analogy"
%  Det Kgl. Danske Videnskabernes Selskab,
%  Filosofiske Meddelelser I, 4.
%  København: Bianco Lunos Bogtrykkeri, 1923.
%
%  TRANSLATION (English).  TRANSLATION-ONLY PROJECT: no Danish
%  transcription is produced or published.
%
%  ---------------------------------------------------------------
%  SOURCES, AND A WARNING ABOUT THE SCAN
%  ---------------------------------------------------------------
%  Base text: the DRM-free Danish e-book (SAGA / Lindhardt og
%  Ringhof, ISBN 9788726363067), an OCR of the 1923 print.  A copy
%  sits in this directory, gitignored by the blanket *.epub rule.
%
%  Page-accurate second witness: the Royal Danish Academy's own
%  digitisation of the 1923 printing,
%  <http://publ.royalacademy.dk/books/109/740>, on disk as
%  ~/bibliotek/Høffding, Harald/begrebet-analogi.pdf (142 pp.,
%  ABBYY FineReader Server, 2019).
%
%  ** THAT FILE IS STRUCTURALLY DAMAGED, AND THE DAMAGE IS AT THE
%  SOURCE ** -- the Academy's own download is byte-for-byte the same
%  file.  Its page tree is malformed from object 685 onward.  qpdf
%  cannot reconstruct it, MuPDF refuses it outright, Ghostscript
%  aborts at page 4 of 142, and poppler's pdftotext returns one page
%  of text out of 142.  So NOTHING in the Linux sandbox can read it,
%  and an agent that tries will conclude, wrongly, that there is no
%  second witness.
%
%  It IS readable on the host, through the PDF tooling there
%  (read_pdf_pages).  Every page checked so far comes back clean.
%  Use that route, and keep the print open alongside the e-book, as
%  on "Filosofiske Problemer" and "Totalitet som Kategori".
%
%  ---------------------------------------------------------------
%  THE PAGE MAP -- FROM THE BOOK'S OWN INDHOLD
%  ---------------------------------------------------------------
%  Like "Totalitet som Kategori" and unlike "Filosofiske Problemer",
%  the 79 endnotes carry no printed "(S. N)" page references; all 79
%  were checked.  The map below therefore comes from the print
%  itself -- the INDHOLD on p. 138 -- not from a reconstruction.
%
%    PRINTED PAGE = PDF PAGE - 2.  (PDF 5 = p. 3, PDF 82 = p. 80,
%    PDF 140 = p. 138.)  Text runs pp. 3-137; INDHOLD on p. 138.
%
%    Indledning                                    3
%    I.   Uvilkaarlige Analogier                   7
%    II.  Analogi og Logik                        37
%    III. Analogi mellem Erkendelsesfunktioner    61
%    IV.  Analogi mellem Erkendelsesomraader      79
%           A. Logik og Aritmetik                 80
%           B. Aritmetik og Geometri              85
%           C. Formal og real Videnskab           89
%           D. Naturvidenskab og Aandsvidenskab  109
%           E. Etik og teoretisk Videnskab       119
%           F. Verdensanskuelse og Videnskab     124
%    V.   Poetisk og religiøs Symbolik           129
%
%  No page number in this file was ever inferred: every section
%  opening below was read off the print as its batch was reached.
%  The complete verified table, of all twenty-two divisions:
%
%    §1  p. 3     §7  p. 41    §13 p. 66    §19 p. 109
%    §2  p. 7     §8  p. 45    §14 p. 70    §20 p. 119
%    §3  p. 16    §9  p. 48    §15 p. 79    §21 p. 124
%    §4  p. 26    §10 p. 52    §16 p. 80    §22 p. 129
%    §5  p. 32    §11 p. 55    §17 p. 85
%    §6  p. 37    §12 p. 61    §18 p. 89
%
%  ("Tillæg til § 19" opens on p. 117; §18's lettered runs f and g
%  open on pp. 99 and 102.)  Chapters overlap by a page at three
%  breaks -- ch. II ends and ch. III's head stands on p. 61, and
%  ch. III ends and ch. IV's head stands on p. 79 -- which the
%  INDHOLD's single figures conceal.
%
%  ---------------------------------------------------------------
%  THE CONTENTS PAGE DISAGREES WITH THE CHAPTER HEADS
%  ---------------------------------------------------------------
%  Recorded, not normalised, per the repo's editorial stance.  The
%  INDHOLD on p. 138 letters chapter IV's six divisions A.--F.; the
%  heads at the top of the text itself number them 16.--21., in the
%  continuous 1--22 sequence that runs through the whole book.
%  VERIFIED on p. 80, whose head reads "16. Logik og Aritmetik."
%  The chapter heads are followed here, being the reading at the
%  head of the text and the one consistent with the run-in numbering
%  of every other division in the book.
%
%  ---------------------------------------------------------------
%  STRUCTURE
%  ---------------------------------------------------------------
%  An Indledning and five chapters, with TWENTY-TWO numbered
%  divisions running STRAIGHT THROUGH the book (1--22) rather than
%  restarting at each chapter.  The print sets them two different
%  ways, and this translation mirrors that:
%
%    * Indledning and chapters I--III and V: RUN-IN numbers, set here
%      with \hnum (= \noindent\textbf{N.}~).
%    * Chapter IV: REAL HEADINGS for 16--21, set here with
%      \subsection*.  Section 15 is NOT in chapter III -- it is the
%      untitled introduction standing at the head of chapter IV,
%      before the first heading, and is set run-in like the rest.
%      (This resolves the one structural question left open by the
%      survey: chapter III carries 12--14 only.)
%
%    ch.   title                                  pp.      §§       notes
%    ---   -----------------------------------    -------  ------   -----
%    (--)  Indledning                             3--6     1        --
%    I     Uvilkaarlige Analogier                 7--36    2--5     1--17
%    II    Analogi og Logik                       37--60   6--11    18--29
%    III   Analogi mellem Erkendelsesfunktioner   61--79   12--14   30--40
%    IV    Analogi mellem Erkendelsesomraader     79--128  15--21   41--77
%            16. Logik og Aritmetik
%            17. Aritmetik og Geometri
%            18. Formal og real Videnskab
%            19. Naturvidenskab og Aandsvidenskab
%              (with a "Tillæg til § 19" on Hobbes)
%            20. Etik og teoretisk Videnskab
%            21. Verdensanskuelse og Videnskab
%    V     Poetisk og religiøs Symbolik           129--137 22       78--79
%
%  Lettered run-in subdivisions (a., b., c., ...) occur inside §§ 3,
%  4, 14, 18 and 19.  The e-book gives them no styling and none is
%  added; they are left as plain text at the head of the paragraph,
%  unlike the section numbers, which the user's specification sets
%  bold.
%
%  ---------------------------------------------------------------
%  THE TERMINOLOGY ANCHOR IS HØFFDING'S OWN ENGLISH
%  ---------------------------------------------------------------
%  Høffding published the English counterpart of the 1905 essay that
%  this book expands: "On Analogy and its Philosophical Importance",
%  Mind 14, no. 54 (April 1905), pp. 199--209 -- a paper read to the
%  Jowett Society in Oxford, 26 November 1904, and referred to by
%  Høffding himself in ch. IV § 18 i.  Openly readable at
%  <https://archive.org/details/sim_mind_1905-04_14_54>.
%
%  It is authoritative in a way Fisher's 1905 translation of
%  "Filosofiske Problemer" was not: it is the author's own English.
%  Its central sentence fixes this book's title concept --
%
%      "Analogy is likeness of the relations of different objects,
%       not likeness of single qualities."   (Mind, p. 200)
%
%  -- confirming the rendering of "Forholdslighed" as LIKENESS OF
%  RELATION that this repo's translations of "Filosofiske Problemer"
%  and "Totalitet som Kategori" had already settled on other grounds.
%
%  Where his 1905 English and the repo's settled vocabulary disagree,
%  the repo is followed and the disagreement recorded, since four
%  Høffding volumes here already share a vocabulary that must stay
%  consistent.  The three disagreements found, in full, are:
%     "Grund og Følge"  -- his English: "reason and consequence";
%                          used here: "ground and consequent".
%     "uvilkaarlig"     -- his English: "spontaneous" (Mind p. 208,
%                          "the spontaneously developed world-views");
%                          used here: "involuntary".
%     "Bestaaen" (of value) -- his English: "persistence";
%                          used here: "subsistence".
%  All three are one-command reversals if a later editor prefers the
%  author's own word; see RESUME-NOTES.md.
%
%  Terms his 1905 English CONFIRMS, and which are used here:
%  likeness of relation; analogies of experience (Kant); scheme;
%  symbol; image; totality; whole; multiplicity; unity; receptivity
%  and activity; world-view; the fate of values.  His own footnote on
%  Mind p. 205 gives the symbol/scheme distinction that ch. I § 5
%  turns on, referred to Kant, Kritik der Urtheilskraft § 59: "the
%  first is a concrete image, the second an abstract sign".
%
%  ---------------------------------------------------------------
%  SMALL CAPS
%  ---------------------------------------------------------------
%  The print sets cited authors' names in small caps.  All are
%  carried as \textsc{} -- 218 spans in the finished file -- per the
%  convention established in "Filosofiske Problemer".  The e-book
%  swallows the space that follows every such span; those are
%  restored silently and the class of defect is logged once.
%
%  ---------------------------------------------------------------
%  EDITORIAL INTERVENTIONS -- the complete list
%  ---------------------------------------------------------------
%  Six places where this translation departs from what the 1923
%  print says, rather than merely from what the e-book says.  Each
%  is argued in RESUME-NOTES.md; they are gathered here because this
%  header is the only part an outside reader gets.
%
%   1. "Prolegomma" -> Prolegomena (note 52).  Kant's title.
%   2. "Erdtmann" -> Erdmann (note 29).  Leibniz's editor; the book
%      itself has "Erdmann" at note 36.
%   3. "Brunschwicg" / "Brunschwieg" -> Brunschvicg (notes 32, 37,
%      42 and ch. IV § 18 i).  The print spells it two ways, neither
%      correct.
%   4. "Phänomenologie" and other third-language titles quoted in
%      running text are given their correct forms, as are Kant's
%      "die wir kennen" (Kritik der Urtheilskraft § 65) and the
%      German and French titles in the notes.
%   5. "(Smlgn. ovenfor p. 32)" in ch. V -> "(Cf. above, § 5.)".
%      THE ONLY PAGE CROSS-REFERENCE IN THE BOOK.  This edition
%      carries no printed pagination, so the reference is converted
%      to Høffding's own §-idiom, which he uses for every other
%      internal cross-reference (§§ 4 b, 5, 7, 10).  Printed p. 32
%      is where § 5 opens.
%   6. Names set as literal capitals or lower case inside a
%      small-caps span (AVENARIUS, KANT, FR. VISCHER, BECHTEREW,
%      "émile meyerson") are set \textsc{} with normal capitals.
%
%  Left AS PRINTED and logged, not emended: "Kritik der
%  Urtheilskraft § 40" for Kant on genius (ch. I § 4 b);
%  "Mysticism and Logic, p. 987" (note 76); "Yrjo Hirn" without the
%  umlaut; "Michel Angelo"; "Gleichniss" at ch. I § 5 against
%  "Gleichnis" at ch. IV § 17, the same line of Goethe spelled and
%  punctuated two ways; "matematisk Proposition" where the sense
%  wants "Proportion", which Høffding also wrote as "proposition" in
%  his own English in Mind 1905; "indbyr-Forhold" for "indbyrdes";
%  and "Arts- og Slægtskabsbegreber" where ch. IV § 16 has the clean
%  species/genus pair.
% ============================================================

\documentclass[12pt]{book}
\usepackage[utf8]{inputenc}
\usepackage[T1]{fontenc}
\usepackage[english]{babel}
\usepackage[protrusion=true,expansion=true]{microtype}
\usepackage[margin=1.3in]{geometry}
\usepackage{setspace}
\usepackage{amsmath}
\usepackage{enumitem}
\usepackage{libertinus}
\usepackage{libertinust1math}
\usepackage{textalpha}
\usepackage{fancyhdr}
\usepackage[colorlinks=true,linkcolor=black,urlcolor=black,
  pdftitle={The Concept of Analogy},
  pdfauthor={Harald Høffding}]{hyperref}
\renewcommand{\thechapter}{\Roman{chapter}}
\setlength{\headheight}{15pt}
\pagestyle{fancy}
\fancyhf{}
\fancyhead[LE,RO]{\thepage}
\fancyhead[RE]{\textit{The Concept of Analogy}}
\fancyhead[LO]{\textit{\leftmark}}
\setlength{\parindent}{1.5em}
\setstretch{1.3}

% run-in division number (Høffding's numbering, continuous 1--22)
\newcommand{\hnum}[1]{\noindent\textbf{#1.}~}

\title{The Concept of Analogy}
\author{Harald Høffding}
\date{Translated from the Danish\\[0.6em]
  \small \textit{Det Kgl.\ Danske Videnskabernes Selskab},\\
  \textit{Filosofiske Meddelelser} I, 4.\\
  Copenhagen: Bianco Lunos Bogtrykkeri, 1923}

\begin{document}
\maketitle
\tableofcontents


% ============================================================
%  INDLEDNING -- pp. 3--6 -- § 1 -- no notes
% ============================================================

\chapter*{Introduction}
\addcontentsline{toc}{chapter}{Introduction}
\markboth{Introduction}{Introduction}
\label{ch:intro}

\hnum{1} In this treatise I have set myself the task of continuing the closer
investigation of some of the fundamental concepts which, in my book
\textit{Den menneskelige Tanke} [The Human Thought] (1910), I attempted to
exhibit as the most characteristic for human cognition, such as we know it
after its development hitherto. I have thus wished to give a companion piece to
the treatise on \textit{Totalitet som Kategori} (Videnskabernes Selskabs
Skrifter. 1917) and to the treatise on \textit{Relation som Kategori}
(Videnskabernes Selskabs filosofiske Meddelelser. 1920), and it is the concept
of analogy that I shall here investigate more closely.

Provisionally, analogy can be defined as likeness of relation between two
subject-matters, that is to say, a likeness which does not hold of single
properties of those subject-matters or of single parts of them, but of the
relation of properties to one another or of parts to one another. There may be
a very great difference between two subject-matters' properties or parts, each
taken by itself, and yet there may be a likeness in the relation between the
properties or between the parts. Analogy is thus a peculiar kind of likeness.
It stands in a certain sense further from identity (absolute likeness) than
qualitative likeness, composite likeness and coincidence do, which concern only
single properties or parts, and which more easily form transitional forms to
identity (cf.\ \textit{Den menneskelige Tanke} § 21). Yet it must be remarked
that there is not only an identity of subject-matter, but also an identity of
relation, and already Aristotle distinguishes between quantitative analogy
(ἰσότης λόγων), proportionality in the strict sense, and qualitative analogy,
which may hold, e.g., between geometrical figures or between organs in
different living beings. And there is besides an identity of difference, which
lies at the base of the formation of an identically varying difference-series
(\textit{Den menneskelige Tanke} § 65) and especially of the formation of the
identical quality-series that are fundamental for exact science (number, time,
degree, place) (\textit{Den menneskelige Tanke} §§ 76--82).

Both totality and relation I have given a place among the categories, that is
to say, among the fundamental concepts, the forms of thought, which at
science's present stage play an essential role in thought-work, determining its
tasks and its whole direction (\textit{Den menneskelige Tanke} §§ 51--52; 57).
Neither Aristotle nor Kant has given analogy a place among the categories, and
in \textit{Den menneskelige Tanke} I followed them in this. Yet in that work I
was led to lay much weight on the fact that, in contrast to purely formal
thinking (logic and arithmetic), analogy plays an essential role in all
thinking about real subject-matters. The opposition between identity (identity
of subject-matter in logic, identity of difference and identity of relation in
arithmetic) and analogy is therefore of great importance for the right
apprehension of the nature and the validity of our cognition. But when analogy
really plays so great a role, while identity is properly an unattainable ideal
outside the formal world of abstractions, it might be warranted to give it a
place in the series of the categories; and in the table of categories to be
found on the last page of the treatise \textit{Totalitet som Kategori} this has
in fact been done. It is entered there among the formal categories, immediately
after the concept of identity, because it lies at the base of that reduction of
relations of quality which --- with the help of points of view of identity ---
takes place in the stricter development of the concepts of number, time, degree
and place. Such series form intermediate members between absolute
identity-series and qualitatively different series of subject-matters. They are
formed by analogizing abstraction from the latter. It is here that analogy
especially shows its importance. It should now be the task of the present
treatise to follow the concept of analogy further in its more special
importance for human cognition.

I dwell first on the involuntary analogies which make themselves felt in
primitive men and in children, but also, at the high points of the life of the
spirit, in the genius. They do not come forward as analogies for the
consciousness in question itself; that is possible only for the critical
consciousness, which has learned to distinguish between analogy and identity
and to see that every grounding and proving transition of thought takes place
in virtue of identity. Thereafter I go on to investigate analogy from a formal
logical point of view, and in the following sections I seek to exhibit
analogies between cognition's different functions and between its different
domains.

Provisionally it may be remarked that analogy's importance for cognition shows
itself in three different respects.

Its greatest importance it certainly has by working in a scouting,
anticipating and spurring manner. Categories in general lay themselves open to
view in the questions that are put (\textit{Den menneskelige Tanke} § 52).
Analogy may be the actual road to true cognition, even if the analogy's
validity is afterwards rejected. One can, after all, draw a correct conclusion
from false premisses. Naturally one cannot prove that the conclusion is correct
except by finding the right premisses; but it often happens that an assumption
stands for us as true, and that precisely thereby it sets us the task of
finding out why it is at bottom true. The premisses (the right premisses) are
often the last thing to be found; but one would not seek after them if an
analogy had not first led us to an assumption and thereby set us a task. The
analytic method rests on this, that a task has been set us, and that one then
seeks the conditions for solving that task. If the right premisses are then
found, one perhaps rejects the analogy that showed the way --- not only if it
is false, but also if it is unnecessary. The ground for an analogy's being set
up will lie in this, that there is a break in a continuous connection which has
formed itself involuntarily for us, and which an involuntary need leads us to
hold fast as long as possible. Analogy may here lead to an intellectual
venture, which is sometimes crowned with success. One may apply to analogy what
Pascal says of the imagination in general, that it deceives all the more easily
because it does not always deceive (\textit{d'autant plus fourbe qu'elle ne
l'est pas toujours}). If it does not deceive, or if it shows itself to be
dispensable, it is replaced by a determination of relation which is finally
grounded in identity.

But analogy is not merely preparatory. It may be the last bond of union between
different subject-matters or series of subject-matters which cannot be reduced
in such a way that a single subject-matter or a single series of
subject-matters is seen to lie at the base, or in such a way that they are all
seen to have their ground in more deeply lying subject-matters or series of
subject-matters. A certain rounding-off and survey becomes possible, in the
face of such subject-matters, only by the help of analogy. We do not thereby
obtain an explanation of why there are different subject-matters or series of
subject-matters, but we establish the actual state of the case, namely that at
certain points we get no further than to a connection, a synthesis, which
builds only on the remotest degree of likeness --- which analogy is. It will
appear that both in the face of cognition's different functions (see ch.\ III)
and in the face of its different domains (see ch.\ IV) we are finally placed in
this way. We end then in open questions, since it remains undecided whether a
further reduction, a more deeply lying continuity, might be found, or whether
there are here, in the nature of things, differences which are at any rate
irreducible for us. Questions can thus be not only the beginning of wisdom, but
also its end.

It will have its interest, finally, to dwell somewhat on analogies whose
importance lies in their being expressions of poetic and religious moods ---
expressions which may recall now the analogies of involuntary mental life, now
the seeking and finding of scientific thinking. Especially will intimations or
hopes of a unity and a continuity lying deeper than experience can exhibit find
expression in this kind of analogy.


% ============================================================
%  I. UVILKAARLIGE ANALOGIER -- pp. 7--36 -- §§ 2--5 -- notes 1--17
% ============================================================

\chapter{Involuntary Analogies}
\label{ch:1}

\hnum{2} As already remarked, analogy does not always presuppose that a
distinction has expressly been drawn between different subject-matters, or
respectively between their properties or parts, and that afterwards, in a
certain opposition to the differences, a connection is brought about between
the subject-matters on the basis of a likeness of relation. In involuntary
mental phenomena, such as sense-intuition, memory and imagination, it cannot
always be shown that the elements are first present as separated and only
afterwards enter into combination. The importance of these involuntary mental
processes lies precisely in this, that they bring about wholes, combinations,
which can be analysed only afterwards. Reflection constantly needs a material,
but this material need not, as has often been believed, consist in a multitude
of independent elements; it may be wholes, which it is then reflection's task
to resolve into elements that may not without more ado be regarded as identical
with the elements to which the involuntarily formed wholes are due. For newly
emerging elements may at once and involuntarily be woven together with
consciousness's given content, so that perhaps no alteration at all is noticed.
All the same, there is a tendency in reflection to regard the elements which
its analysis finds as the whole's original, elementary foundation. There is
here an analogy present, which at most can seek its justification in the
continuity with which the transition can take place from an involuntarily
formed mental continuum to the differentiation to which the very arising of
reflection already testifies. Properly, when we speak of the most primitive and
involuntary mental processes, all designations fall short, and this is no
wonder, since language's psychological designations are fetched from the
clearly differentiating, critically reflecting mental life. The very expression
``analogy'' we here use in an analogizing manner. It is here as when one speaks
of ``unconscious inferences''. These are not inferences in the proper sense,
but processes, transitions, which can be interpreted in analogy with the proper
inferences. The same holds further (as I have already shown in my
\textit{Psykologi} and later in \textit{Den menneskelige Tanke} § 5) when we
set up the concept of synthesis as an expression for the nature of all
conscious life. In epistemology, where we investigate the possibilities of
scientific cognition, we must also discuss the possibility that the given might
consist of a heap of equally great and mutually irreducible differences. Only
an approximation to such a chaotic difference-series can be exhibited (as also
to its counterpart, an absolute identity-series); there would in any case be
needed a searching and strenuous reflection to establish that such a series
really lay before us in a determinate given case.

In the manner in which a single lived experience, from the very moment of its
conception, is connected with other simultaneous or past lived experiences, or
with traditions into which the individual has lived himself so that these too
can be regarded as lived experiences, there lies already an interpretation. The
further back we come towards what we call the primitive man's stage, the closer
lived experience and interpretation move to one another, and finally there is
no opposition between them to be found. When we nevertheless attempt to hold
them apart, we make use of an analogy with what first makes itself felt for
testing reflection. And, as has been said, the very concept of analogy is used
analogically when we say that primitive men interpret every new lived
experience in analogy with what they know from the traditions into which they
have lived themselves. --- In my treatise \textit{Oplevelse og Tydning} I have
sought to describe different ways in which the relation between lived
experience and interpretation comes forward historically, as regards ecstatic
experiences. ---

In his important work \textit{Les fonctions mentales dans les sociétés
inférieures} (1910), \textsc{Lévy-Bruhl} has criticized the formerly common
conception that primitive man apprehends everything on the likeness of himself,
anthropomorphizes or analogizes. It is, on his conception, not comparison or
reflection that is at work here, but an entirely immediate identifying. In
primitive man there always lie traditions at the base, with which the new that
is lived through involuntarily melts together. The representations handed down
from the tribe make themselves felt with such strength of feeling that no need
for reflection arises. What deviates in high degree from what has been handed
down is provisionally not noticed; only later is an interpretation perhaps
thought out that agrees with the tradition. When the savage sees an image of an
object he knows, the image as image brings him nothing new, for the image is
the object itself, however he may afterwards arrange this relation for himself.
And with all the logical consistency which already distinguishes this stage ---
which Lévy-Bruhl has wrongly called the ``prelogical'' --- the savage therefore
fears to become dependent on one who has an image of him, and he fears that he
will come to suffer want of game when a European has drawn or photographed many
wild animals in the district and has departed with the pictures. And as it goes
with images, so it goes also with names and with ceremonies: they are what they
mean. Images, names and ceremonies are indeed only a part of a whole which has
other properties besides those that are depicted or named; but this part
\textit{is} the whole. Lévy-Bruhl finds in this identifying of part and whole
--- what he calls ``loi de participation''\footnote{To this corresponds what
\textsc{Grønbech} (\textit{Primitiv Religion,} p.\ 8) calls
``Helhedsrealisme'' [realism of the whole], and which is characteristic for the
primitive mental life.} --- the most characteristic peculiarity of the
primitive mental life. It is certain that the interpretation which primitive
man institutes under the influence of a tradition embraced with heartfelt
adherence and reverence cannot follow the principles of formal logic. These
make themselves felt only when doubt has been awakened; but he is accustomed to
regard the tradition as a self-evident foundation for his apprehension of
everything that happens. The basis of interpretation is not the work of
reflection. But it may lead to misunderstanding when Lévy-Bruhl declares
(p.\ 425) that the collective representations governed by the ``law of
participation'' are ``indifferent to the logical law of contradiction''. A
distinction must be drawn between the given basis of interpretation (here: the
tradition held fast in reverence, which is once for all stamped by the ``law of
participation'') and the interpretation of single phenomena that takes place on
the ground of this basis. Reflection works at all times and on all levels of
culture on the ground of determinate presuppositions which cannot themselves at
once be made an object for reflection. On the basis which primitive man once
for all has, it is entirely consistent that he does not understand why the
European distinguishes between the image and the thing. When Lévy-Bruhl calls
``the primitive mentality'' ``prelogical'', this fits the basis of
interpretation, not the interpretation made on the ground of that
basis.\footnote{Dr.\ \textsc{Albert Schweitzer}, who works as a physician in
the primeval forests of equatorial Africa, writes, after several years' stay
among the negroes: ``The natural man thinks more deeply than one commonly
believes. Even if he can neither read nor write, he yet institutes reflections
to a far wider extent than we suspect.'' (\textit{Mellem Floder og Urskov.}
Danish transl.\ 1922. p.\ 134.)} As there will later be occasion to mention,
Lévy-Bruhl has himself in another connection maintained, and rightly, that
human consciousness, as regards the ultimate basis for judgments and
inferences, properly never quite loses its ``prelogical'' character. There will
always be presuppositions present which we do not and cannot bring forward,
even in our most logical reflection. --- As regards especially the close
connection between name and thing, one can find examples of it even today.
There is a tendency to regard the terminology one uses oneself as the only
``right'' one, and the language one speaks oneself seems to be one with the
thing, while foreign languages' designations often seem foreign to the things.
It is told of a German who visited Paris that he reported home that there they
called bread ``pain'', and added, ``Wir nennen es Brot, und es ist doch auch
Brot!''

Tradition comes forward for the savage as a closed whole, within which he has
no ground for making any distinction. There provisionally rules an immediate
belief in the validity of the representations handed down, just as in the
sense-perceptions that immediately appear. And the expression ``belief in'' is
even too reflected, for no distinction at all is drawn between the belief and
that which is believed in. The very concept of reality or validity is first
formed when doubt has made itself felt. The fundamental proposition in
accordance with which thinking provisionally goes (though it is of course not
formulated) can be expressed by the question ``Why not?'' Why not rely on
everything that comes forward for us with living force as sensation or as
tradition? Only later does one learn to ask ``Why?'' --- and reflection can
then even reach the point of directing this question against the very thing
that has hitherto formed the undisputed basis. For that it is required that new
lived experiences should appear so numerous, and with such a peculiarity
deviating from everything customary, that they can gain power over what one has
already lived oneself into. \textsc{Albert Schweitzer} declares, as the result
of several years' living together with negroes in Central Africa: ``One must
not believe that one has got to the bottom of the negro's world of thought if
one has got no further than to look at his superstitious representations and
his inherited rules of law. Innermost in him there lives a dark intimation that
the apprehension of what is good must come by way of reflection'' (\textit{Mellem
Floder og Urskove} p.\ 136). The self-sacrificing physician, who through
historical criticism had learned to distinguish between the Christian dogmas
and the ethical idealism that comes forward in the Sermon on the Mount and in
the parables, experienced this especially when he sought to explain his
religious apprehension to the negroes. Yet he experienced also that, alongside
the ``morality of the heart'' that could be called forth in them, there long
maintained itself a desire to lie and to steal, and a need for revenge which
according to their own traditions was their right and perhaps even their duty.

I must here interpolate the remark that when in my \textit{Psykologiske
Undersøgelser} (p.\ 24 f.) I say that in the relation between new lived
experiences and practised, customary representations the chief weight lies on
the new, which ``sets the whole process going'', this presupposes another stage
of development than the undisturbed primitive mentality, where the new so
easily glances off against the customary and is able to gain power in the mind
only when it comes forward in a particularly insistent manner. ---

If one asks after the origin of the system of collective representations that
makes itself felt in the savage and determines his interpretation of what meets
him, Lévy-Bruhl answers that it ``is lost in the night of time''. And since
every human group hitherto met with comes forward with certain collective
representations, we cannot exhibit any absolutely first beginning. We must
here, as with the question of the origin of language,\footnote{Cf.\
\textsc{Otto Jespersen:} \textit{Language, its Nature, Development and Origin.}
(1922). p.\ 416 f.\ on the method in the discussion of the question of the
origin of language. One must (as already, in their way, Madvig and Whitney had
maintained) start out from the manner in which language actually develops, that
is, from its history, and from there infer back to a more primitive language
than any of those we know.} content ourselves with analogies fetched from the
levels of culture that lie nearest to us, and these analogies must of course be
used with great caution. We see that on the levels of culture we know, the
traditions that lie at the base of each new generation's spiritual development
undergo constant alterations, additions and displacements, which for an
essential part are due precisely to the new generation's own experiences and
reflection. The collective is constantly altered under the influence of the
individual. We cannot, as was often attempted in earlier times, think of a time
when there was nothing collective, or when what collective there was was due
exclusively to single individuals' experience and reflection. But just as
little can we think of a time when minds were governed exclusively by common
traditions, without individual experiences and thoughts exercising any
influence whatever. It must be assumed that on all levels of development there
has taken place a reciprocal action between the common and the individual, and
the difference of the levels can consist only in this, that on earlier levels
the common, on later levels the individual, has the greatest influence --- yet
so that there is again and again a swinging to and fro between the different
factors' predominance. Even on the most primitive levels we know, individual
differences can make themselves felt and deviations from the tribe's belief
come forward.\footnote{See my treatise \textit{Sociologi og Religionsfilosofi}
(in the book \textit{Oplevelse og Tydning}. 1918) p.\ 131--138.} And the
examples of such things would undoubtedly be still more numerous if we had more
knowledge of life within primitive societies than we have. We should have found
remarkable deviations and initiatives which now are drowned for us in the
ruling traditions.

Primitive conditions of life present great uni\-for\-mity, and therewith great
com\-mu\-nity, in the experiences which the single individuals can make. And on
account of the constant living together in group and family, the single man's
special experiences and inventions will very quickly become common property,
and his name will be forgotten. Added to this, the practical demand constantly
makes itself felt; action is required, and life is for a great part lived as
reflex-life and instinct-life. Every representation gets its essential
importance by releasing movement. But for that it is not necessary to enter
into the whole connection to which the representation's content belongs. A
single part or a single property is enough to release action, and one need not
distinguish between part and whole; the part stands involuntarily as a sign of
the whole. The image, the name or the sign works as if the object itself were
given in it. The instinct has enough in an intimation to begin to stir. In the
fact that a part can thus have the same influence as the whole, and that it may
even be advantageous that time and force are not expended on investigating
whether the whole is really represented by the single part or property, there
might lie an access to an understanding of the ``law of participation'' or the
``realism of the whole''. And when a part thus shows itself again and again, in
the race's experience, to have the same practical importance as the whole, one
sees the possibility that there may be formed a circle of collective
representations which obey the law of participation, and which it may be
difficult for new lived experiences to break.

Purely logically, this identifying of part and whole (\textit{pars pro toto})
can be formulated as an inference in the second figure: the part releases a
certain movement, which the whole would release; therefore part and whole are
the same! And this is precisely one of the ways in which an inference from
analogy can be formulated. When an instinctive action can lead astray, this
hangs together with the fact that the ``inference'' in question is wrong. Only
within quite determinate limits, which the consciousness governed by an
unconditional impulse to action does not give itself time to establish, can it
do to identify part and whole. Neither instinct nor prelogical mentality can
therefore maintain itself longer than the inference from analogy practically
suffices. Doubt will be awakened when this is no longer the case. A closer
psychological investigation will be able to elucidate this relation more nearly
and thereby cast light on the involuntary analogies.

\hnum{3} Just as the adult European is too inclined to read his own developed
and reflecting mental life into the mental life of the so-called primitive men,
so there is in general a tendency to read functions that presuppose a certain
development of reflection into the most immediate and involuntary mental
processes. One then speaks of association of representations and of judging,
where in reality it is more elementary processes that are going on. This comes
in part of the fact that language's expressions for mental life are essentially
fetched from the developed, differentiated mental life, and one is then not
clear that it can be warranted only by way of analogy to use such expressions of
the most elementary psychical processes. Some years ago (\textit{Psykologiske
Undersøgelser}. Vid.\ Selsk.\ Skr.\ 1889) I had occasion to draw attention to
this, and I now bring forward a few single points from that investigation,
especially in order to show that what Lévy-Bruhl has exhibited for the mental
life of primitive men is found again in mental processes that go on in all of us
to this day.

a. It has not been sufficiently noticed that there are several different kinds
or forms of recognition. In the treatise just mentioned (p.\ 35 f.) I thought I
had to distinguish at least ten kinds of recognition. It leads to confusion when
one speaks, as is commonly done, of recognition quite in general and yet in fact
keeps to a single form --- and that a form which presupposes a certain
differentiated development of the life of representation. There is then a
failure to recognize the peculiarity of the elementary, involuntary mental life.

The most elementary form of recognition that can be exhibited is as immediate
and as involuntary as a sense-sensation can be. A subject-matter here stands for
us at one stroke as familiar, with a quality of familiarity, just as a sensation
has its sense-quality, without any association, comparison or judgment being
demonstrable. A name, a facial expression, a shade of colour can appear with a
quality of familiarity without our being able to remember that we have observed
the same subject-matter before, and without our being able to refer it to any
determinate connection --- that is, without knowing who bears the name, and so
on. Perhaps the quality of familiarity gives us occasion to form the judgment
that we have heard the name before, and so on. But that judgment then comes
precisely afterwards. And if one will conceive such immediate recognition as
association, one must formulate it as a limiting case and think of the matter
thus: that when a sensation of sight or of sound is repeated, it can, in the
very moment of its birth, fuse with the after-effect of earlier sensations of
the same character (however one may care to conceive such an after-effect).
There takes place, then, a momentary reproduction, and what is reproduced fuses
at once with the sensation that is coming into being. The quality of familiarity
can then theoretically be conceived as the product of two elements.
Physiologically regarded, we have here an example of practice. All practice of a
function is due to a fusion of a new release with the after-effects of earlier
releases.

This exhibiting of immediate recognition as an elementary phenomenon, lying on
the boundary between sensation and memory proper (I have called it bound
memory), was not acknowledged or attended to when it first appeared. It was
rejected by Danish psychologists, and writers such as Wundt, Meinong and Bergson
assumed without more ado that association, comparison or judgment is always the
condition of recognition. It has interested me the more to see that G.\ E.\
\textsc{Müller}, who is an authority as regards the doctrine of representations,
in his psychological lectures printed as manuscript (1913) § 26 describes the
same phenomenon that I call immediate recognition, under the name of
indeterminate recognition. ``Recognition of an object of observation is,'' he
says, ``indeterminate when, in observing the object, we merely have the
impression that we must somewhere have observed it before. When the object not
only comes forward for us as familiar, but in addition awakens the recollection
of its name, of the circumstances under which we observed it earlier, and so on,
then we have a determinate recognition.'' I have against this description only
this to object, that the quality of familiarity need not be bound up with a
consciousness of (``the impression of'') our having apprehended the object
before. It can be without consequence. It is already an inference when we take
it that we have apprehended the object before --- an inference drawn from the
difference between sensations with and without the quality of familiarity, an
immediate inference of course. When immediate recognition at once releases a
reflex or an instinctive action, no such inference is drawn, and ``the
impression'' will not make itself felt.

\textsc{Müller} finds (in a handwritten addition to the printed lecture) the
explanation of indeterminate recognition in this, that after-effects of earlier
observations, and perhaps indistinct reproductions of these, are released by the
new sensation; or that subject-matters we have apprehended earlier are, on
repetition, apprehended more easily and more quickly than before. In quite the
same way I had also referred to the law of practice, which finds application
both to the psychological and to the physiological side of the phenomenon. It
can of course often be difficult to decide whether, and in what degree, an
indistinct association of representations takes place. I have therefore
designated immediate recognition as a limiting case, and I have in my
\textit{Psykologi} compared the relation between immediate recognition and
association of representations with the relation between the tangent and the
secant: the nearer the secant's two points of intersection with the circle come
to lie to one another, the more does the secant approach the tangent. Repeated
and exact observation has shown me many degrees here, and there is every
occasion for closer investigation.

Even before \textsc{Müller}'s lectures were printed, \textsc{Betz} (Archiv für
die gesammte Psychologie. 1910, p.\ 267) had declared: ``Erst die vielzitierte
Abhandlung von Høffding über Wiedererkennen zeigte, dass die Dinge bei weitem
nicht so einfach liegen, als es den Assozia\-tions\-theo\-re\-tikern schien''. --- The
Neo\-pla\-tonist Olym\-pio\-dorus,\footnote{\textit{Olympiodori in Phædonem
Commentaria,} ed.\ W.\ \textsc{Norvin}. 1913. p.\ 68.} moreover, had already
mentioned the phenomenon. In his commentary on Plato's \textit{Phaedo} he
supplements his master's doctrine of recollection by drawing attention to the
fact that there is not only a transition from one element to another on account
of likeness [association by likeness], or from one element to another element
different from it with which it has earlier occurred together [association by
contiguity], but that also ``the same can be connected with the same''. When
someone recognizes Socrates, for instance, after having forgotten him for a
time, then it is, says Olympiodorus, just as much a sensation as the first time
he saw Socrates; in such recognition there takes place no transition of
cognition (οὐ γέγονε μετάβασις γνώσεως). That Plato does not mention this kind
of recognition, Olympiodorus explains by the fact that in the \textit{Phaedo} he
had use only for the free (``determinate'') recognition, in which, according to
him, ``the Ideas sprang forth in the soul''.

b. When there lies before us a composite subject-matter whose elements (parts or
properties) we have often apprehended in succession, we need not always go
through all the elements anew for recognition to take place. Successive
recognition can be replaced by partial recognition. If the subject-matter's
elements are \textit{a b c d}, it can be enough that \textit{a} is recognized.
We have here quite the same phenomenon that \textsc{Lévy-Bruhl} has designated
as participation. Partial recognition has great practical importance, since it
would often take too long a time if the whole series of elements had to be
recognized before we were sure of having the subject-matter before us, so that
we could proceed to act towards it. The subject-matter thus contains more than
we notice, but we confine ourselves to the strictly necessary. We recognize,
e.g., a man by his walk or by some peculiarity of his dress, a chemical
substance by its smell, and so on. As with immediate, indeterminate recognition,
it is the need for action that shortens the process. The necessity of
intervening, if the recognition holds good, weighs practically more than the
certainty that the apprehension is correct. When spontaneous or instinctive
action leads into misfortune, it is because the single mark does not hold good.

Partial recognition can become sense-illusion when the immediately recognized
element has in different cases been connected with different other elements.
When both the series \textit{a b c d} and the series \textit{a f g h} have
occurred, and I, when \textit{a} comes again, at once rely on \textit{b c d}
being there too, I can be surprised by \textit{f g h} presenting themselves
instead. I use an example from my \textit{Psykologi}. On a walk along the
Jammerbugt on Jutland's west coast it happened to me, a whole day through, that
I took pieces of wreckage sticking up out of the sand for men. There took place
here an involuntary supplementing of what was really seen. The image that
involuntarily formed itself for me contained more than proved to be given. The
supplementing here takes place just as involuntarily as in the well-known
experience with the blind spot, or as when, after a part of the lower lip has
been locally anaesthetized, one feels a hole or a depression in the rim of the
glass one drinks from. Such a supplementing lies, seen from our modern point of
view, at the base of primitive man's participation. The ``prelogical'', which
only by a questionable analogy can be described as having arisen through
association or through judging, we thus know from daily experience. Purely
logically, sense-illusion can be formulated as a fallacy; but it is all the same
unwarranted on that account to read such an inference, as an actually occurring
psychical process, into the elementary phenomena we here have to do with. We
simply have no right to identify the elements which analysing reflection finds
with those that co-operate in the original lived experience. Here the analogy
between psychological analysis and chemical analysis fails.

When one formulates the sense-illusion, or partial recognition in general, as an
inference, it is the second or the third figure that we get. The property
\textit{B} is found both in the subject-matter \textit{A} and in the
subject-matter \textit{C}; from this is inferred partial or absolute identity of
\textit{A} and \textit{C}. (Second figure.) Or: the subject-matter \textit{B}
has both the property \textit{A} and the property \textit{C}; from this is
inferred that there is partial or absolute identity between \textit{A} and
\textit{C}. (Third figure.) As we shall see, we have here two of the ways in
which an inference from analogy can logically be formulated. What is
questionable in every inference from analogy already comes forward clearly
here. ---

In all participation, that is, in immediate, partial and illusory recognition,
an involuntary need for whole-formation makes itself felt. Even if only a single
element is actually given, one at once makes it into a whole. \textsc{Grønbech}
has therefore rightly designated what Lévy-Bruhl calls participation as realism
of the whole. Such an involuntary impulse toward the whole is especially
characteristic of the first stages of development of conscious life, which bear
the stamp of naivety; it does not --- fortunately --- disappear later either,
though it becomes an object of analysis and criticism. It lies at the base of
what I have called concrete and practical intuition. By concrete intuition I
understand sense-intuition, memory or imagination. In all these three forms,
which do not from the first stand as different in principle from one another, a
whole is built up on the occasion of a single element or of a few elements.
Since it will often be a ruling feeling that determines the centre about which
the whole-formation takes place, concrete intuition is more or less akin to
practical intuition, in which all the mind's forces can work along, without
there needing to be any consciousness of the sources of the whole-image or of
the energy with which it is formed and held fast. All practical faith, indeed
all personal concentration upon a thought or a task, is of this kind. Here too
we are reminded of the ``prelogical'' consciousness, which according to
Lévy-Bruhl has a purely emotional and mystical character, since it is the
tribe's sacred traditions that decide how the individual interprets, or rather
integrates, what he lives through.\footnote{Cf.\ my treatise on \textit{Begrebet
Intuition} (Vid.\ Selsk.\ Oversigt 1914) p.\ 75--77 and my work \textit{Henri
Bergson's Filosofi} p.\ 22 f. --- In contrast to the two kinds of intuition here
mentioned I set analytic and synthetic intuition; the former expresses
reflection's testing, the latter its gathering function.}

c. In all recognition, likeness to what has earlier been lived through plays a
role, and to that extent recognition can be regarded as a limiting case of
association by likeness. As regards association by likeness in general, there
has been dispute whether any such thing exists at all. My position on this
question has, ever since the first edition of my \textit{Psykologi}, been that I
do not assume any pure and unmixed association by likeness, since in all
demonstrable examples of likeness there will always be certain differences
present. No one has been clearer about this than Plato --- it is of course for
that reason that he makes perfect likeness an ``Idea''. And the elements of
difference will be able to release association by contiguity, which plays a
greater or a smaller role in the transition to the following representation.
Every actual association will therefore be a mixture of association by likeness
and association by contiguity, and the discussion must in each single case turn
on the degree to which moments of likeness co-operate. My formula for
association by likeness is therefore \textit{a}\textsubscript{1} +
\textit{a}\textsubscript{2}, and the indices designate the differences which
always make themselves felt, even where the likeness is preponderant. Since it
is peculiar to such transitions in the life of representation as are essentially
due to likeness that they take place involuntarily, and especially when the mind
is strongly occupied and moved, it will be difficult to exhibit such transitions
experimentally; it is very difficult in psychological experiments to avoid
preparation and expectation, whereby association by contiguity is at once
favoured.\footnote{See, besides my \textit{Psykologi}: \textit{Psykologiske
Undersøgelser} (1889), p.\ 43--59.}

While many psychologists (especially Danish ones) took a quite dismissive
attitude towards the assumption of an association by likeness, even with the
limitations stated, G.\ E.\ \textsc{Müller}, who has quite especially
investigated connections of representations, expressed himself more cautiously.
He did not find it probable (in the lectures mentioned above) that there was
pure association by likeness; when a portrait calls forth the representation of
the person himself, it may happen by its first calling forth the name, and this
then the person. There is of course no doubt that it can go this way. I have,
however, not been able to convince myself that it fits all cases; I can well
recognize a portrait even if I have forgotten the name, which perhaps even after
the recognition I have trouble in recalling.

In the same year in which Müller had his lectures printed as manuscript, he took
part as a subject in some experiments conducted by his pupil Louise Schlüter,
which were to test the relation between two methods of teaching a foreign
language, of which the one builds on a direct connection between an object and
the foreign word (the intuition method), while the other first names the
domestic name of the object and lets that recall the foreign word. It showed
itself here that when, the day after an object had been shown and the foreign
word learned, a similar object was shown, the foreign word came up without any
support from the domestic word. This presupposes that the object seen on the
preceding day was reproduced from the object now seen, although it was not quite
the same. This Louise Schlüter regards as proof that there can be a real
association by likeness,\footnote{\textit{Beiträge zur Prüfung der Anschauungs-
und der Übersetzungsmethoden} (Zeitschrift für Psychologie. 1913) p.\ 73; 103.}
and Müller has declared himself in agreement.\footnote{In handwritten additions
to a copy of the lectures, which Professor Axel Herrlin has kindly placed at my
disposal.}

I shall only remark on this, that in observations and experiments on this
relation a distinction must be drawn between the different kinds of likeness. In
Louise Schlüter's example it was, so far as I can see, qualitative likeness that
was in question: the object shown on the second day was not the same as the one
seen on the first day, but ``ein ähnliches Exemplar derselben Art''. It was thus
not what I call coincidence that made itself felt here. But since the relation
between examples of the same species is for the most part likeness of relation
--- inasmuch as it need not be a likeness with respect to a single part or
property that is decisive for the species-concept --- it would be worth
investigating how, in cases like Schlüter's, qualitative likeness and likeness
of relation stand to one another.

At bottom, Schlüter's case is not different from the cases on which
\textsc{Müller} in his lectures (§ 20) has built the form of association which
he calls active substitution, and which is an association by likeness based on
an association by contiguity. When \textit{a} is associated with the
representation \textit{b} different from it, then a representation \textit{α}
which resembles \textit{a} can also call forth \textit{b}, whether the relation
between \textit{a} and \textit{α} is qualitative likeness or analogy.

Such active substitution is of great interest for our investigation of analogy
in its epistemological respect. For when it shows itself that \textit{a} and
\textit{α} practically play the same role, whatever kind of likeness it may be,
we see here a psychological possibility for that confusion of identity and
analogy which proves so fateful in the history of thought. There will from the
first be no doubt about the validity of the transition from \textit{α} to
\textit{b}. Practically regarded, there is a tendency to treat all likeness as
identity. This is seen already from the much-discussed phenomenon that a burnt
child dreads the fire. According to Stuart Mill, we have here merely an
association by contiguity. But the child provisionally shuns all objects that
resemble the one on which it once burnt itself.

\hnum{4} In the child's play and language, in artistic intuiting and in the
genius's working, involuntary analogies make themselves felt which are as little
as those discussed in the foregoing due to express comparison or inference, and
are not due either to an association in which every single member is apprehended
by itself and only thereafter connected with other members.

a. The child's play is often due to imitation of the grown-ups' doings, but the
child can immerse itself in it with so complete a living-into that the
difference between what the play ``is'' and what it is to ``mean'' does not make
itself felt at all. If the play is interrupted from outside, this can awaken
pain and indignation. A child was playing housewife and had laid its play-table
with various imagined dishes. When it now became the house's dinner-time, and
the child was wanted to sit down at the real table with the rest of the family,
it became violently incensed and burst into tears, ``for a mother doesn't just
run away like that from the table she has to look after''. The indignation shows
how deep the living-into the play had been. In a similar way it can go with the
spectator at a drama, when he is suddenly disturbed in his full merging into
what he sees and hears. It has wrongly been attempted to explain the attraction
of a tragic drama by supposing that the reader or spectator is conscious of
standing outside what is presented. Such a knowledge would make the full
absorption impossible. The absorption presupposes that we for a while identify
our own fate with the one presented. It is an illusion, a ``participation'',
which is perhaps lifted when the curtain falls or the book ends, but whose
resonance or ``perseveration'' can make itself felt for a long time, since we
may have a tendency to live again and again in the same mood as when we were
immediately under the drama's influence.

Outside play, too, children show a capacity for analogical lived experiences. A
state or a situation that has brought them satisfaction they will have
continued, either in such a way that they themselves are still part of it, or,
more disinterestedly, so that the state is merely continued after the same
scheme as before. In both cases it is equivalents that are sought.
\textsc{Avenarius} has, in his interesting principal work \textit{Kritik der
reinen Erfahrung}, sought to show that here --- in the need to continue a state,
even if one thereby goes beyond the old forms --- a universally human tendency
makes itself felt. He adduces the following amusing example from child
psychology (which decidedly does not belong among the disinterested ones). A
little boy is sitting comfortably enough in an armchair and letting his guest, a
little girl, stand. His mother reproaches him for it, and he vacates the seat,
whereupon the girl settles herself in the chair with her toy on her lap. The boy
looks at her a little and then proposes that she should lay the toy on the floor
and let him sit on her lap.\footnote{\textit{Kritik der reinen
Erfahrung}\textsuperscript{2} II p.\ 493, cf.\ I, p.\ 136. --- \textsc{Avenarius}
finds the same tendency in such cases as those where men, through sudden blows
of fate or through the enfeeblement of age, can continue their spiritual life
only by the help of forms of faith which in their vigorous years they had
abandoned. \textit{Kritik der reinen Erfahrung}\textsuperscript{2} II, p.\ 295.}
--- To the disinterested continuation of a situation, on the other hand, belongs
the following example. A little boy was riding on his grandfather's knee, but
had to break off the situation when he was to go to bed. A little after, he came
rushing back with great eagerness, bringing his toy bear, which he placed on the
grandfather's knee. In this way a new member was set into the relation, which
otherwise would have had to cease. Perhaps there was also sympathy with the old
man, whose lap was now empty. ---

In the child's acquisition of language, not only recollection and imitation of
the grown-ups' speech, but also involuntary analogizing, plays a role. The child
has, e.g., heard the expression ``Jacob's hat'' together with the corresponding
visual image, and now forms on its own account the expression ``Peter's hat'' on
a given occasion, without the general representation ``hat'' being a necessary
intermediate member. Children moreover often form regular verbal forms on their
own account, when they have not learned the irregular ones. Afterwards they can
then practise using the irregular forms which usage requires, instead of the
regular ones they had themselves formed in involuntary
analogizing.\footnote{\textsc{Otto Jespersen:} \textit{Language,} p.\ 130 f.}

b. The fact that creation by genius was comparatively late acknowledged as
something other and more than the product of a mechanical joining-together of
particulars in a conscious purpose\footnote{Cf., e.g., the way in which
\textsc{Holberg} himself explains the coming into being of his various works.
See my \textit{Mindre Arbejder} II, p.\ 150--154.} shows how hard it has been to
grasp involuntary psychical whole-formation. All mental processes were regarded
from the standpoint of analysing reflection. After Rousseau (especially in the
article ``Génie'' in his \textit{Dictionnaire de Musique}) had described genius
as the capacity for involuntary creation, and had entreated ``l'homme vulgaire''
not to misuse this sublime word, it was \textsc{Kant} who fixed the concept in
the form in which it is now commonly used, in apprehending genius as that rare
originality of natural endowment which makes it possible to bring forth works of
exemplary character --- works which cannot be imitated, but which can awaken
originality in other natural endowments (Kritik der Urtheilskraft § 40). Genius
is, like instinct, an immediate union of need and capacity. The involuntary and
the exemplary together make up the essence of genius. \textsc{Goethe} rejoiced
(Dichtung und Wahrheit. 19. Buch) that the best word for involuntary and
model-worthy creation was thus rescued. --- A fate like that of the word genius
threatens the word imagination, when one uses it of all free representation,
even where this is merely recollection, and does not restrict it to mean the
capacity for the new formation of concrete individual representations.
\textsc{Meumann} is surely right when he holds\footnote{\textit{Intelligenz und
Wille}\textsuperscript{2}, p.\ 136.} that the narrower meaning of the word
imagination first arose under the influence of the modern age's more penetrating
study of artistic production. Artistic imagination leads to works which to an
observer can seem to be the fruit of a reflection that with full consciousness
lays in the single features in such a way that they co-operate with other
features to the formation of a whole. On the occasion of a feature which Solger
had brought out in his analysis of the characters in
``Wahlverwandtschaften'', \textsc{Goethe} said: ``Ich habe selber nicht daran
gedacht, aber es liegt darin''. (Gespräche mit Eckermann. 21.\ Januar 1827.)

The artist thus comes involuntarily to say more than he is conscious of. That he
analogizes, and thereby comes to connect elements that were otherwise separated,
he does not notice until perhaps later, if testing reflection awakes. By a
living-into --- or, as the Germans say, a feeling-into (Einfühlung) --- of the
subject-matter, which involuntarily leads him to seek and to find expression for
its value, he can awaken, if not an independent creation, then at any rate an
independent capacity for a living-into like his own. In this latter living-into
it is not only the subject-matter's content, but also the form under which it
comes forward in the work, that exercises a decisive effect --- just as the
artist's living-into the subject-matter leads him from the first to a forming
activity which might let the subject-matter come forward in its whole fullness.
In the definition which \textsc{Julius Lange} (Om Kunstværdi. 1876) gave of
artistic value, as the value the subject-matter has for the artist and which it
acquires through his work for the beholder, the form of presentation did not
come into its own (as I sought to show in a review of Lange's fine book in
``Nær og Fjern'' 1876); this Lange regarded as something purely external and
mechanical. In the artistic genius, need (feeling of value) and capacity
(forming power) are joined from the first --- and joined in the most intimate
way, since there is a need to use the capacity, and not merely to feel the
value, and since the feeling of value presupposes a capacity for opening oneself
to what can call forth that feeling. \textsc{Friedrich Vischer} has rightly
maintained that the feeling-into must determine both the work's form and its
content. Lange, by emphasizing the living-into of the subject-matter, pointed to
what is involuntary in the imagination and in the genius; but to this belongs
also the involuntarily forming activity. On Lange's definition it might look as
if the work of art were merely an external, inessential means by which others
than the artist himself could experience the subject-matter's value. Need and
capacity are not bound together in so outward a way, even though the carrying
through of the forming activity can at many points make possible, indeed
require, conscious deliberation and choice.

The concept ``Einfühlung'' was known neither to Lange nor to me when we
conducted our discussion. It had been set up some years earlier by the German
aesthetician \textsc{Robert Vischer}, but became more generally known only
through \textsc{Volkelt}'s and \textsc{Friedrich Vischer}'s
writings.\footnote{\textsc{Johannes Volkelt:} \textit{Der Symbolbegriff in der
neueren Ästhetik.} 1876, p.\ 55. --- \textsc{Friedrich Vischer:} \textit{Das
Symbol} (Festskrift til Zeller. 1887) p.\ 170--187. --- \textsc{Rousseau} knew
feeling-into at first hand. His joy in nature springs from it. Cf.\
\textit{Confessions} (éd.\ Bever) I, p.\ 263 f.: Je dispose en maître de la
nature entière; mon c\oe ur, errant d'objet en objet, \textit{s'unit,
s'identifie} à ceux qui le flattent. --- But in Rousseau the need was greater
than the capacity. A work of art he could not form (apart from the work of art
his own Confessions are). But his heartfelt feeling gave the impulse to modern
poetry.}

While for the primitive consciousness the image stood as the bearer of a magical
power which made its possessor master over the subject-matter, in art it is the
bearer of a value that was born in the artist's soul and through his work is
born again in the beholder's soul. In both places there is an analogy which
works in an involuntary manner, and which is seen as analogy only by testing
reflection. That the involuntary mental life's manner of working has not been
understood led to a failure to recognize the nature both of the genius and of
primitive man. But the understanding that can be won of the involuntary mental
life comes about itself by the help of analogy, since the understanding starts
out from the differentiated life of consciousness, which can make the primitive
syntheses clear to itself only by dissolving them.

Answering to the image's magical significance there stands, in art, its symbolic
significance; but in the history of religion magic and symbolism are broken
wholly apart only in the most recent time.

\hnum{5} The symbol rests on participation, in that a part stands as the
expression of the whole and works as such. The small can become a symbol of the
great when it presents in itself relations that have a likeness to relations
within the great. For the great humour, everything single, however limited it
may be, can become a symbol by being seen as part of a great whole and as
presenting the same type as the whole. ``The humblest thing defends its place in
life's garland.'' But life's garland is never wholly bound, and the possibility
is always present that new flowers are to be woven in. To this the great humour
will not be blind; it looks humorously on its own garlands,\footnote{\textit{Den
store Humor.} Chap.\ 6 (Forstaaelse og Humor).} and it agrees with
\textsc{Goethe} (Gespräche mit Eckermann. 5.\ Juli 1827) that in images used as
expressions for life there always remains something incommensurable for the
understanding. The image's value rests precisely, on one essential side, on its
giving as a whole what for testing reflection stands as an unclosed connection
of elements --- perhaps of constantly changing elements. Seen from this side,
the image gives less than the always advancing experience and reflection. As a
whole-image it gives more, since as such it has elements not to be found in the
subject-matter from which it is fetched. When God is called Father, the element
of omnipotence comes in with it, which unfortunately is lacking in every father
given in experience, who must often struggle hard both for himself and for his
own. Or (to use \textsc{Fr.\ Vischer}'s example) when the lion is a symbol of
magnanimity, it has a property the real lion lacks. Pseudo-hallucinations are
therefore often better suited as a foundation for symbols than hallucinations
proper, which as a rule take in all the elements of reality.

Religious symbolism is especially characterized by the fact that its foundation
of feeling has formed itself under the struggle for what, for the human group in
question, is the highest thing of value. The images in which expression is found
for the experiences made in the struggle for existence are fetched from the
great oppositions of nature and of life: light and darkness, life and death, joy
and sorrow, power and powerlessness, good and evil; and, as the most elementary
nature-religions give place to higher forms of religion, a need makes itself
felt to see these oppositions represented and working through more or less
individualized beings.

But like the concept of analogy in general, the concept of symbol too is first
formed at the standpoint of testing reflection. The difference between part and
whole, between the sign and the signified, does not make itself felt from the
first. What for the artist and the artistic beholder holds only in moments of
enthralment is, for the positively religious consciousness, so long as it
subsists in its full peculiarity, the constant state. Every distinction between
part and whole, image and subject-matter, is heretical. Hegel wrongly held that
in all symbolic apprehension a dissatisfaction must be felt on account of the
ever-present disproportion between form and content. Since Hegel presupposed
such a dissatisfaction already at primitive stages, and in general within the
positively religious consciousness, it became comparatively easy for him to show
that it was a liberation for the religious consciousness to pass over to forms
that could subsist for testing reflection; and the difference between religion
and philosophy became for him in the end only a difference of form, while the
spiritual content was to be the same.\footnote{\textit{Den nyere Filosofis
Historie.}\textsuperscript{3} III, p.\ 182--186.} Hegel reads the cleavage to
which reflection can lead into the religious consciousness itself, for which, on
the contrary, the relation between part and whole, sign and subject-matter is a
relation of identity. It is denied, e.g., from the side of positive religion,
that the expression Father used of God is figurative. So far is it from being
experiences of the significance of the human father-relation that are to lie at
the base when God is called Father, that the case is precisely the reverse: it
is God who is the ground and the model for every father-relation in the world
(the Epistle to the Ephesians III, 15). By virtue of this identity it is
possible that actions which for reflection stand as symbolic --- such as prayer,
sacrifice and sacrament --- can be assumed to exercise an influence on the
Godhead. This is consistent, for by virtue of the relation of identity there is
a real connection between sign and object. Added to this, the positively
religious symbolism is the result not only of the individual's but of a whole
human group's experiences of life, and thereby comes forward for new generations
as the given and self-evident expression of what can be experienced about human
life and its conditions.

Concepts such as analogy and symbol are correlative concepts; they express a
reciprocal relation. As we have already seen with the father-symbol, the
relation therefore turns quite naturally round. From the world of experience in
which human life is led, symbols were fetched for a higher world; but in the end
the very world from which the symbols were originally fetched becomes a symbol.
In the mystics of the Middle Ages this turn comes forward in a characteristic
way. For Hugo a St.\ Victore the whole bodily world is a symbol of the spiritual
world. And when Goethe uses medieval expressions to denote the soaring above
everything finite, he declares

\begin{verse}
Alles Vergängliche ist nur ein Gleichniss.
\end{verse}

Medieval thinking moved, as it has aptly been said, as often from symbol to
symbol as from experience to symbol. There had once for all been awakened a
fundamental mood which constantly sought new expressions, and in relation to
which different experiences of life could mean the same. Here the emotional
foundation is at work which is in general determinative for all
``participation'', all ``prelogical mentality''. Only critical reflection makes
that demand on religious representations which Schleiermacher maintained, that
each single symbol shall be formed by itself on the foundation of a particular
lived experience. In the doctrine of faith, he says, the one proposition must
not be derived from the other, but is to be set up immediately with reference to
a single determinate lived experience, and only then afterwards compared with
other, previously established propositions. (Der christliche Glaube § 17.) So
long, on the other hand, as the once developed fundamental mood maintains
itself, new symbols will spring again and again from it, the one developed in
analogy with the other. Thus \textsc{Yrjo Hirn} (Det heliga Skrinet. Studier i
den katolska Kyrkans Poesi och Konst. 1909) has shown in a fine way how the
heartfelt worship of the Madonna leads to a series of images which express the
fundamental mood from different sides: the dwelling-place of the Host in a
shrine on the altar, the place where the divine comes forward in bodily form,
became a symbol for the Madonna, in whose womb God became man. But the Madonna
is at the same time the morning star that heralds the sun's coming, and she is
the spring, where light and warmth relieve the darkness and cold of winter. The
mighty expansion of feeling constantly puts out new shoots.

The involuntary formation of images is the presupposition of all dogmatics. Not
all symbols become dogmas. There is talk of dogmas only when certain symbols
stand as norms which individuals ought to follow in their attempts to interpret
their experiences of life. A choice is then made, and what is chosen stands fast
in virtue of the religious community's authority over its members. The dogma
rests on a mixture of lived experience, reflection and authority. From what has
been dogmatically fixed a struggle is then carried on against freely developed
symbols, as well as against the attempts to apprehend the dogmas themselves as
mere symbols. Dogmatics' struggle against symbolism is a struggle between ``it
is'' and ``it means''.

A free formation of symbols within a religious frame is presented by Jesus'
parables, in which through narratives or features seized out of life he makes
intuitable the relations and the demands of the promised kingdom of God. An
ethical-religious wisdom of life is here uttered in poetic form. The figurative
is chosen with deliberate intent, and often (especially in the parables of the
Gospel of John) passes over from symbolism to allegory.

The symbol becomes allegory the more expressly a distinction is drawn between
the two members in the symbolizing. In allegory the relation between the two
members is more external than in the symbol. Poetic symbolism springs, like
religious symbolism, from the enthralment a subject-matter has awakened; but it
does not exclude, as religious formation of images in the classical ages of the
religions does, that the enthralment may alternate rhythmically with a clear
consciousness of the formation of images that has taken place.

In intuitive natures, where intuition and imagination have the preponderance
over reflection, there will be a tendency to live in symbols as long as
possible. Still stronger will this tendency be in ecclesiastical natures, for
whom it is a need to hold fast to inherited symbols --- not only on account of
their intuitability, but because they are the inherited expression of the
fathers' faith, and thus on account of a feeling of solidarity with the past.
They hear a centuries-old tone in the church's hymns and find again the race's
dearly bought experiences in the church's dogmas. Yet it must be noted that it
need not be intellectual criticism alone that bursts the inherited forms; it may
also be a new and strong movement of feeling that makes it impossible to find
one's own experiences of life expressed in these forms.\footnote{Cf.\ my
\textit{Religionsfilosofi} §§ 95--96. --- \textit{Oplevelse og Tydning} p.\
102--107.}


% ============================================================
%  II. ANALOGI OG LOGIK -- pp. 37--60 -- §§ 6--11 -- notes 18--29
% ============================================================

\chapter{Analogy and Logic}
\label{ch:2}

\hnum{6} In the involuntary analogies, likenesses and differences are from the
first so woven into one another that it can be difficult to find the boundary
where the likeness ends and the difference begins. The definite attempt to find
such a boundary, through a sharp determination of the different degrees of
likeness (and of difference), marks the beginning of scientific thinking. Plato
stands in the history of thought as the first who drew attention to there being
many steps upward from the likenesses that come forward for immediate
apprehension to pure identity, the ``Idea of likeness'', which alone makes
strict science possible. What Parmenides had proclaimed, under the name of ``the
One'', as the sole condition of true knowledge is determined by Plato more
exactly as identity. --- In my treatise on the Platonic dialogue
\textit{Parmenides} (Vid.\ Selsk.\ filosofiske Meddelelser Nr.\ 3) I express
myself perhaps too strongly when I say (p.\ 28) that Plato ``expressly
declares'' unity and identity to be the same; more cautiously it is said on
p.\ 30 that it is ``let out'' that unity and identity are the same. The fact is
that Plato in the strange second part of the dialogue apparently denies this;
but since he grounds the denial only on the point that \textit{several} things
(τὰ πολλά) can be One (i.e.\ can come under the same concept) without
coinciding, and adduces no proof at all in the case of ``the One'', which is
just what is in question, I find here an involuntary concession. So much is
certain, that Parmenides' ``the One'' is obscure and mystical when it does not
mean identity --- that is, when it does not express at once the fundamental law
of all thinking and the highest ideal of all thinking.

For Plato, now, the Idea of likeness was not merely fundamental law and ideal,
but the expression of true being, and the degrees of likeness were so many lower
and higher steps of existence. All the same, in the setting up of this Idea a
direction was given for the dissolution of the involuntary analogies. When that
dissolution is carried out with a hard hand, absolute likeness and absolute
difference come to stand in opposition to one another, identity in opposition to
chaos, and it becomes cognition's task to find a road between these oppositions.
It is characteristic of the sharp Greek thought that not only the Idea of
absolute identity but also the Idea of absolute difference is brought forward,
albeit with a polemical purpose (by Zeno). And towards the close of his career
Plato saw that the principle of identity stood in the end as powerless in the
face of the manifold differences that existence presents, and therefore refers
to analogizing as a kind of ``spurious thinking'' (λογισμὸς νόθος. Tim.\ p.\ 52
B). Analogy stands here, then, as a last resort, where no higher degrees of
likeness can be exhibited. Whereas in mythological thinking it was everything,
it now comes to denote the lowest stage of cognition. The programme of European
thinking was herewith given. And the task then becomes to exhibit the different
intermediate forms between chaotic difference and absolute identity, for it is
in these intermediate forms that human thinking actually moves. The order of
rank of these intermediate forms (of which I have described six in
\textit{Den menneskelige Tanke}) will depend on how far the subject-matters
before us can be ordered in such a relation to one another that inference
becomes possible. It will then show itself that the more inference becomes
possible, the more we approach pure identity. It shows itself at the same time
that identity no longer stands as the highest form of existence, but as the form
of thought which makes grounded cognition possible without calling in the help
of observation at every moment. The principle of identity dawns most clearly in
such series as those where each member stands to its predecessor as that one
does to its own predecessor. If, e.g., \textit{A} is subject to \textit{B} and
\textit{B} to \textit{C}, or \textit{A} is smaller than \textit{B} and
\textit{B} smaller than \textit{C}, then \textit{B} can be eliminated and we get
a new judgment. The presupposition is \textit{B}'s identity with itself, whether
it be set in relation to \textit{A} or to \textit{C}. But such a presupposition
holds for every member in every series we set up. To that extent the principle
of identity lies at the base of all series-formation; but cognition consists not
merely in the formation of single concepts, but in the exhibiting of
intelligible relations between concepts, in rational series-formations.

The expression ``participation'', so happily chosen to denote involuntary
analogy, can well be retained after logic has begun to have a word to say
concerning the presuppositions of the course of thought. It denotes a partial
likeness which in the involuntary analogizing was made into an absolute
likeness, in spite of the difficulties there are in grounding the right to
assume, from a likeness at a single point, an agreement in general. It is
precisely these difficulties in ``participation'' that set scientific thinking
its tasks.

If one seeks a presupposition which for all that lies at the base both of
prescientific thinking's involuntary analo\-giz\-ing and of the sci\-en\-tific search
after grounded tran\-si\-tions, and which lies at the base of all series-formations
from the chaotic difference-series to the absolute identity-series, then such a
presupposition lies in the law of synthesis, which holds both in empirical
psychology and in rational epistemology. It lies in the concept of a series ---
whether the members are held together by mere memory or by grounded
understanding of the transitions --- that different members are gathered
together with one another, either as different or as alike. If thought
constantly lost the preceding members when it went over to the following ones,
no series-formation would take place. The more it becomes possible to ground the
transition from member to member, the smaller the role that mere memory ---
which is sole ruler in the chaotic difference-series --- needs to play; the
transitions then take place in virtue of a law.

The fundamental law of synthesis rules at all stages of the con\-nec\-tion of
rep\-re\-sen\-ta\-tions, and thereby it becomes possible to compare these different
stages and to exhibit the relation between them. If the different sciences,
which find themselves at different stages between chaos and identity, did not
have such a common fundamental form, we could not set up the very concept of
science as common to them all. We can now elucidate not only one science through
analogy with other sciences, and one part of a science through analogy with
other parts, but also prescientific thinking through analogy with the
scientific; and we can within psychology elucidate the elementary mental
functions by analogy with the more developed ones. Such analogies are not mere
fictions, as \textsc{Vaihinger} held in his \textit{Philosophie des Als
Ob}.\footnote{An interesting presentation and criticism of \textit{Philosophie
des Als Ob} has been given by \textsc{Jørgen Fr.\ Jørgensen} in the journal
``Vor Tid''. 1914--15. p.\ 227--249.} In virtue of the continuity which through
the fundamental form of synthesis makes itself felt in all life of the spirit,
and especially within the intellectual life, analogies such as those mentioned
become possible; naturally their logical value is highly various. It is, e.g.,
no fiction when a tangent is called a secant whose two points of intersection
with a circle approach one another in the highest degree; for the synthesis can
bring its elements together in a great many degrees. According to a modern
tendency in geometry the tangent even always touches the circle in more than one
point, so that the difference between secant and tangent falls away in the end.

\hnum{7} Even between the two outermost limits of our cognition, the chaotic
difference-series and the absolute identity-series, analogy is present. The only
operation of thought they make possible is --- apart from useless substitutions
--- counting. One can count how many absolute differences, or how many absolute
identities, one knows. As we shall later see, arithmetic takes hold precisely
where pure logic ends. Another analogy between them lies in this, that neither
of the two series can be immediately given. Only through a strictly critical
testing, which could not be shown to be exhaustive, could it be shown that there
were no points of likeness at all, or no differences at all, within a series of
subject-matters. The two series mark limits of our cognition which cannot be
proved to be present in any single case whatever; only through arbitrary
assertions could their validity be maintained. But these series serve indirectly
to characterize human cognition, by exhibiting what would make it impossible
from the first, or would be a conclusion admitting of no possible continuation.

One might say that epistemology has only the task of discussing the principles
that make cognition possible, and that it must therefore first of all set up the
principle which lies at the base of all formation of concepts and judgments and
inferences whatever, namely the principle of identity. Cognition consists of
judgments, one might say, and every judgment asserts a relation of identity
(partial or absolute). The concept of identity would therefore seem to have to
open the series of the categories. To this it must be answered that epistemology
ought to investigate all possibilities. Every actual cognition is always only an
approximation to identity, and all degrees of approximation must be
investigated. Epistemology's proper subject-matter is the degrees of clash
between identity and difference. Identity itself has only this significance,
that it sets the advancing reduction of the differences as the great task.
Identity is, after all, itself a relation. To strive after finding identity is
therefore to seek to put one relation in the place of another. The concept of
identity thus presupposes the concept of relation and the concept correlative
with it, synthesis, and these become the first categories. Besides, identity can
be defined only as a limiting case of likeness or of continuity, and these two
concepts therefore go --- when we proceed
analytically\footnote{Cf.\ \textit{Den menneskelige Tanke.} p.\ 177 f.} ---
before the concept of identity.

When it is said that cognition consists of judgments, it must not be forgotten
that the judgment's subject is always taken out of a whole connection, and that
this taking-out becomes possible only through the difference there is between it
and other subject-matters that are members of the same connection. Only by such
a distinguishing does a subject-matter come to exist for us with such
independence that closer determinations of it can be sought and found. As it has
aptly been said, every concept has fringes or rents which show that it is
properly a fragment. Only when this separating-out has been carried out can a
judgment come into being. If every subject-matter drowned in, or were wholly
fused with, the medium in which it exists from the first, no judgment would ever
be passed. A relation of difference is both the logical and the psychological
presupposition of a judgment's coming into being.

Naturally the principle of identity is presupposed in all logical thinking, also
in the setting up of a doctrine of categories, and especially in the setting up
of the categories synthesis and relation. Epistemology must of course build on
the principles all science builds on; and since it is precisely these principles
that epistemology wishes to bring forward and elucidate, it can to that extent
be said to build on itself, or to move in a circle. Occasion has sometimes been
taken from this to demonstrate the impossibility of an epistemology. So
Schelling, Lotze and, in the most recent time, \textsc{Leonard Nelson} (partly
in his work \textit{Über die sogenannte Erkenntnistheorie} 1908, whose first
section is entitled \textit{Die Unmöglichkeit der Erkenntnistheorie}, partly in
a lecture at the congress of philosophers at Bologna in 1911). Nelson's train of
thought is this: the criterion of cognition's truth must either be cognition or
not be cognition. If it is cognition, then this cognition must itself require a
criterion. And if it is not cognition, it must all the same be something that
can be an object for cognition, and then we require a criterion for that
cognition's being true. --- This train of thought we know already from
antiquity. It was brought forward on the sceptics' side against the Academics'
attempt to set up a criterion of reality, when they asked by what criterion we
can decide that the criterion of reality so set up is the right
one!\footnote{See more on this in \textit{Relation som Kategori.}} At the base of
this scepticism, as of the modern objections to the possibility of epistemology,
lies a confusion of the concepts of circularity and incompletability.
Epistemology's task must always be taken up anew. It must always be tested anew
whether the principles which epistemology sets up as conditioning the validity
of cognition coincide with the principles that lie at the base of the very
attempt to set up the epistemology. If there is a difference here, the
investigation must begin over again. A conclusion of these investigations is not
to be thought. It may very well be that a new investigation of human thought-work
can lead to the necessity of a new epistemology. Before Carneades too, and in
modern times before Hume and Kant, men had an epistemology; but it showed itself
that it had not been carried through, and that its fundamental significance had
not been demonstrated.

Epistemology comes into being through the investigation of thought-work itself,
that is to say, of the alteration which the application of certain principles
brings about in respect of clarity and connection in our representations. If we
find advancing clarity and connection, we say that epistemologically a step
forward has been made, even if we perhaps come to strike out certain assumptions
on which we formerly laid weight. The more numerous and the deeper the
differences between the subject-matters that can be brought out, and the more we
nevertheless succeed in finding grounded connection between them by the
application of the principle of identity, the more we become convinced of that
principle's validity. As we have seen, the principle of identity can be said to
dawn already in the prescientific analogies, as well as in those
series-formations which do not yet form a foundation for inferences. When, as in
a regularly varying difference-series (e.g.\ the colour-series), each member
lies fast, so that no interchange can take place, the principle of identity has
come more into force than where the order is either quite arbitrary (as in the
chaotic difference-series), or where the order is due to a wholly external point
of view, e.g.\ a scheme for the memory (as in the indeterminately varying
difference-series). And the history of the sciences shows us many stages in this
respect before full clarity and certainty are attained.

\hnum{8} In the analogies that play so great a role in prescientific
thought-life there already stirs the energy which can later lead to scientific
thinking. They ground the expectation with which men come to meet the
subject-matters that present themselves to them, and which makes them not stand
passive in the face of what observation offers. With a sanguine confidence in
all emerging representations, wherever they may come from, the life of thought
begins its course of development; indeed, even before free representations
appear, an analogy to this confidence declares itself in the involuntary
propensity to immediate motor reaction, to which reflexes and instinctive
actions testify. Life begins with a surplus of energy and has a tendency to
spread and unfold itself. Before memory in the narrower sense of the word can
form itself (with a clear consciousness that what is represented belongs to the
past), expectation can be at work; and representations which are not bound up
with any expectation, but whose content stands definitely as past, are held fast
only when a higher development has been reached. Hume rightly laid weight on the
need for expansion, for paying out the line, for continuing a train of thought
one has once entered upon. But when he would derive from this tendency the
belief in rational and causal principles, he overlooked that these principles
come forward most clearly when expectations are disappointed or rejected, and
when we are then compelled to seek an order and connection of another kind than
the one our involuntary analogies and expectations had announced. Rationality
and causality are norms which limit and dam in our involuntary course of
thought, which in itself has a tendency to pursue certain ways of looking at
things beyond all limits. The hypotheses which science on its way has shaped,
but has then let fall again, testify that it is for a great part through
disappointments that we work our way forward to truth. If disappointments had
not come, we should at many points not have got beyond the prescientific stage
of development; but the disappointments would not have come if the surplus
energy had not fostered expectations which were involuntarily tested by our
acting on them. Children's need to grasp and work upon all things leads them to
make important experiences. They experiment involuntarily. The need for movement
enters by degrees into the service of curiosity. In his autobiography the German
surgeon Schleich relates that at five years old he had already been given a
watch, but smashed it by knocking it with a stone; and when he was then scolded
for it, he said: ``Jungens müssen doch wissen, wa da'inn is!'' This need to know
what is ``inside'' stirs from the childish experimenting all the way up to the
highest speculations, without there always stirring any intimation of how fine
and how complete is that connection of things into which we seek to penetrate.
That need is by degrees dammed in, as it is experienced that determinate laws
rule for things, and that it is through the understanding of these that we may
perhaps penetrate into what we call nature's interior.

In a spirited work (De la contingence des lois de la nature. 1875)
\textsc{Émile Boutroux} developed the conception that the laws of nature signify
essentially forms or channels for the unfolding of free forces --- the river-bed
for a stream that works its way forward. A step further went \textsc{Henri
Bergson}, in demanding of philosophy that it should penetrate behind the
mechanism of the laws of nature and make itself one with the force, the
unfolding of life, which is constantly hampered and limited by the order of
nature; and he found a foundation for this in an immediate intuition of our own
inner unfolding of force (\textit{élan vital}), in analogy with which we must
then apprehend all other forms of existence than our own (Introduction à la
Métaphysique. 1903). --- Such attempts of thought are in reality due to a
translation into metaphysical form of what psychology and epistemology teach.
They are good examples of what we shall in a later connection go into somewhat
more closely, that all metaphysics (cosmology) builds on analogy --- here on
analogy with the way in which we see, in the history of science, the human
impulse to expansion and human impatience clash with the fixed order of things
which science works to find. Of special interest here is a chapter in
\textsc{Ernst Mach}'s work \textit{Erkenntnis und Irrtum} (1905). The chapter is
headed ``Sinn und Wert der Naturgesetze'' and seeks to show that the laws of
nature come into being for us as limitations which, under the guidance of
experience, we set to our expectations, and that through such advancing
limitation there takes place an ever closer adaptation of thought to facts. Even
daily experience shows us a constant damming in of the involuntary life of
thought.

\hnum{9} Just as the cognition of laws is a limitation of the freedom to form
connections of representations, so the formation of concepts is a limitation of
the freedom to shape the single representations. We form concepts in order to be
sure that a representation which emerges anew in our train of thought has the
same content as when it appeared formerly; only thereby can it enter as a fixed
member in a series of thoughts and stand in clear and determinate relations to
the other members of the series. If it is merely a matter of holding fast what
is peculiar to a subject-matter in a quite determinate situation, this is
attained by what one might call a concrete individual concept. Thus ``Napoleon
at Waterloo'' is a concrete individual concept, when we have noted Napoleon's
whole conduct on that occasion. If the concept is to hold for all sides of a
subject-matter's essence and for its appearance in all different situations ---
e.g.\ for Napoleon according to his whole character and course of life --- we get
a typical individual concept. Each single situation or each single action then
serves only as an example of this concept. And in a similar way it is with
general concepts, which are to hold for a whole group of subject-matters. The
``common'' properties will here, just as little as with the typical individual
concepts, be sufficient to be the content of a concept by itself. The ``common''
cannot be thought by itself alone. More or less distinctly, on the other hand, a
single subject-matter stands as an example of the whole series; or one makes use
of a scheme or a word which, if it be needed, can be replaced by an example.

Between the different examples themselves there will not be complete identity in
all respects. It is some point of likeness (perhaps of identity) that leads one
to set subject-matters together as examples of the same concept. An Indian who
saw a book open for the first time called it an oyster. That was already an
attempt at a formation of a concept. When men early spoke of feet both in
animals and in men, and later extended the concept to tables, chairs, mountains
and so on, that too was an attempt at the formation of a concept. But not even
in scientific formation of concepts do the single examples completely cover one
another. It is by the help of analogy that one goes from example to example
within a concept's extension. Under the concept of the conic section belong
circle, ellipse, parabola, hyperbola, and the relation between these
subject-matters is not identity but analogy. Only from a determinate point of
view, which the definition of the conic section states in words, can they be
regarded as identical, in so far as any one of them can in principle be used as
an example just as well as the others. The very definition of the conic section
(as a figure that arises by a plane section through a circular cone) requires a
construction which leads to one or another of the examples adduced. It is the
relation between properties or parts in a series of subject-matters that makes
these subject-matters examples of the same concept. With the concept of colour
there is provisionally only qualitative likeness between the single examples, and
a concept proper is got only by keeping (as with the conic sections) to the way
in which they arise. Each single colour-quality stands to a certain degree of
refraction of light as the other colours stand to certain other degrees of
refraction. For brevity's sake (when one has not time to bring forward examples)
words or signs (schemes) can express the concept. Thus letters in arithmetic are
a scheme for the concept of magnitude, whose examples are had in the members of
the number-series. ``Mathematicians'', says \textsc{Pascal} (De l'esprit
géométrique), ``and all who proceed methodically, give things names only in order
to shorten the course of thought''. The word or the scheme is to give directions
for the construction of an example. Pascal requires that for safety's sake one
should try to put what is defined in the place of the definition. But it must be
remarked on this that, since ``what is defined'' comes forward in different
shapes, this substitution can never become absolute. No example contains merely
the features the definition states. Each single example requires other examples,
and it is the analogy between them that is the concept's proper content. The
concept's significance consists in its being able to set a series of thoughts
going. \textsc{Goblot} therefore says that a concept is nothing but a possibility
of a series of judgments (une virtualité, une possibilité indéfinie de
jugements).\footnote{\textit{Traité de Logique.} 1918. p.\ 87.} He therefore
begins his exposition of logic with the judgments and goes from there over to
concepts and inferences. Earlier still, \textsc{Sigwart} had maintained the
priority of judgments in relation to the concept. But the matter is surely more
involved than the two excellent logicians present it. There is a constant
reciprocal action between concept and judgment. The formation of concepts itself
takes place through judgments, but the judgments are clear only when the
concepts connected in them are each determinate. The history of thought shows us
how the starting-point is taken now in the single concepts, now in the judgment.
In epistemology this comes forward in the relation between fundamental concepts
(categories) and fundamental propositions
(principles).\footnote{Cf.\ \textit{Den menneskelige Tanke.} § 51; 57; 106.}

There is something in our thinking (and in our consciousness in general) which
never comes forward for immediate observation, but which must be presupposed
or, if one will, constructed, because what immediate observation exhibits is
otherwise unintelligible. Thus synthesis, the fundamental form of conscious life
and the first fundamental concept of all cognition, is not itself an observable
element. Nor can our willing as such be immediately
observed.\footnote{\textit{Psykologi}\textsuperscript{6} p.\ 177 f. (Forkortet
Udgave p.\ 33 f.; 84 f.)} And so it goes with all inner and outer activity: we
observe its results and infer to the energy to which the work is due. The
activity of thought itself is not observed. Every single thought is a member in
a course of thought, and we infer to thought's energy from those alterations in
the course of thought which are independent of momentary observations (except in
so far as the alteration consists in such observations being brought into
determinate connection with the previously given content of consciousness).
Everywhere we observe only succession, not the activity which stretches through
all the members of the succession and which we often have laboriously to trace
out. Here Hume's criticism still holds good, in the domain of the human sciences
as in that of natural science. His violent isolation of the single members in
the course of thought led indirectly to the right apprehension of the concept of
activity in all domains. But it is Leibniz who has directly indicated the right
application of concepts such as force, energy and working, in maintaining that
before there can be talk of force a lawlike connection between changing states
must have been exhibited.\footnote{\textit{Den nyere Filosofis
Historie.}\textsuperscript{3} II p.\ 136 f.}

Thinking proper is distinguished from the ordinary course of representations
only by being directed towards the conditions for solving a task, towards means
and ways of reaching a goal. It consists in an analysis of the connection within
which the task or the goal is assumed to have its place, and then seeks to find
a law among the elements to which the analysis leads. A simple example is
presented by orientation when one has lost one's way. The different directions
are tried and compared. If, on the other hand, one wanders without a determinate
goal, one strikes into the first road that offers. Already \textsc{Hobbes} ---
who so often has to be the scapegoat in those criticisms of ``association
psychology'' of which people seem unable to tire --- saw this, when he declares
that it is through constantly holding fast to a goal (\textit{frequens ad finem
respectio}. De Corpore XXV, 9) that we learn to find order and connection.

\hnum{10} The concepts of analogy and of inference from analogy were first
brought out with clarity by Aristotle. Whereas for Plato, even after he had had
his last single combat with the Eleatics, it stood as the great goal to lead all
cognition back to one Idea, Aristotle is convinced that we end in a plurality of
concepts (categories), between which not identity but only analogy holds.
Between the different domains in which thinking works, analogies too can perhaps
be exhibited; thus Aristotle holds that the relation between matter
(possibility) and form (actuality) recurs everywhere --- but this relation
itself takes on a different character in the different domains (that is to say,
within the different sciences).

As regards the conducting of proof, Aristotle distinguishes between deduction
and induction. By deductive inference he understands the syllogistic bringing of
something under a genus-concept by the help of a species-concept as intermediate
member. His doctrine of inference is thus built on a presupposed classification.
By induction he understands a complete enumeration of all the subject-matters
that come under a concept (ἡ ἐπαγωγὴ διὰ πάντων. ed.\ Bekker p.\ 68 b). But
there is also, he adds (p.\ 69 a), an inference from a single example, or rather
from example to example. This inference thus goes neither, like deduction, from
a whole to a part, nor, like induction, from all the parts to a whole, but it
goes from one part to another part (ὡς μέρος πρὸς μέρος). It is distinguished
from induction in that not all the parts, but only a single part, lies before
us. Aristotle's designation for inference from analogy is Paradigm
(παράδειγμα), which means example, but can also mean model. The single
subject-matter can by inference from analogy be said to stand as a model, a
type, in whose light a whole circle of subject-matters is seen, in that they
present like relations between properties or parts among themselves. --- The
word ἀναλογία itself Aristotle uses mostly in the sense of mathematical
proposition, identity of relation (ἰσότης λόγων), while for qualitative likeness
of relation he frequently uses the indeterminate expression τὸ ἀνάλογον.

In more recent times it has been seen more and more clearly that there is no
difference in principle between induction and analogy, since on the ground of
experience one can never be sure of having all possible subject-matters of a
certain kind given.\footnote{Except of course in such cases as those where it
lies in the very definition of the concept that there must be a certain number
of subject-matters. Cf.\ \textit{Den menneskelige Tanke.} p.\ 172 f.} Incomplete
induction is then distinguished from inference from analogy only by having more
examples to start out from. Naturally it is of importance to have as many
examples as possible to start out from. But for modern methodology the main
thing is whether, by an analysis of the subject-matters before us (one or
several), a proposition can be found which can provisionally be set up so as
then to be tested by seeing whether the inferences drawn from it find
confirmation in new experiences. This method, which was first applied on a great
scale by Galilei and Newton, is characteristic of modern natural science.
Analogy and induction give, as \textsc{Kant} expresses it in his Logic (ed.\
1801, p.\ 210), merely ``logische Präsumtionen''. Strangely enough, he does not
add that it may all the same often be of importance to test these surmises.

\textsc{Kant} himself saw clearly that there is a kind of inferring which is
neither deduction, induction nor inference from analogy. He refers here to the
principle by which mathematical thinking differs from formal logical thinking.
In virtue of this principle an identically varying difference-series can be
constructed, that is, a series within which the difference between the members
following one another is always the same --- where, accordingly, there is
identity of difference. The simplest example of this is the number-series, which
is built up by what Kant calls ``Synthesis des Gleichartigen''. Besides in the
number-series, the same principle appears --- surely under the influence of the
number-series as a model, and so in analogy with it --- in the series of times,
degrees and places, when (as with number) all special qualitative properties are
purged away from these concepts. Judgments built on such series-formations
(which may be called identical quality-series) Kant calls synthetic judgments
\textit{a priori}. In analogy with these series built on identity of difference,
Kant has now constructed a concept of ``possible experience'' or ``nature'' as a
scheme within which every subject-matter to which we are to be able to ascribe
reality must find its place. We shall at a later point have occasion to discuss
the significance of this construction.

While we are on the prehistory of the concept of analogy, it must be pointed out
that in philosophical attempts to reach a world-view, a concluding cognition,
the dispute in modern times stands between those who would reach a conclusion
through inferences, that is, by the help of the principle of identity, and those
who maintain that analogy is the last means of thought when it is a question of
conclusion. It is the opposition between Plato and Aristotle repeating itself in
new forms. On Plato's side stand Spinoza and Hegel, on Aristotle's side Leibniz,
Schopenhauer, Lotze and, in the most recent time, Bergson. From \textsc{Søren
Kierkegaard}'s \textit{Papirer}\footnote{\textit{Papirer.} Ed.\ P.\ A.\ Heiberg
and V.\ Kuhr. V. p.\ 29.} one sees that he was aware of the opposition between
identity and analogy at the decisive points, in good agreement with the critical
tendency of his epistemology. He had become clear that a systematic conclusion
was impossible by the help of the principle of identity; and since induction and
analogy permit no strict inferring, he came to the result that one always
concludes (as one begins) with a leap --- and it was surely a satisfaction to
him to find an application here for one of his favourite concepts. But the
problem is still more involved than he saw, and it is precisely the leap that
poses it anew: the leap takes place in a curve, and so brings about a new
continuity, though on another plane.

\hnum{11} When a manifold of subject-matters lies before thinking, its first step
will naturally be to order them in series whose members are as closely connected
as possible.\footnote{In what follows I make use of the attempt to set up and
order the most important forms of series-formation which I have made in
\textit{Den menneskelige Tanke.} p.\ 162--173.} The most unfortunate relation
would be that between all the subject-matters there subsist equally great
differences, so that they can be put in any order whatever. As already shown
above, it would be wrong to believe that such chaotic difference-series were
peculiar to the primitive consciousness; this is characterized precisely by an
involuntary interpretation, which lies at the base of the reflection it is able
to institute. Only for a critical investigation, which is not content with any
interpretation that offers, can it look as if the given were a chaos. As we have
seen, the reality of a chaotic difference-series could not be proved. But the
provisional assumption of a chaos can contain an invitation to closer
investigation, and here analogy offers a starting-point. What would be common to
the members of such a series would be precisely the circumstance that they are
elements in a chaos; that holds for \textit{A}, for \textit{B}, for \textit{C},
the whole series through. Herein lies the possibility of an inference in the
Aristotelian second figure, one of the ways in which an analogy can be
formulated. No positive inference can be derived from this figure, but there may
be occasion to investigate whether after all there might not be other points of
likeness between \textit{A}, \textit{B} and \textit{C} than the one meagre point
of belonging to a chaos. The second figure has in general the significance of
setting a task. Already in this simple example, then, we see analogy's
significance for thinking, as spurring, provoking and awakening. Analogy asks;
but whether it answers is quite another question. And its very spurring effect
presupposes that the participation has been dissolved, for by participation
likeness at a single point is precisely made into likeness as regards the whole.

We do not get much further when the members lie fast and cannot without more ado
change places, as when the order of succession is determined purely
chronologically. It would have to be that the question were expressly raised why
the subject-matters were experienced or remembered in just this determinate
order. This too presupposes a wakeful reflection. A child who is taught the
Mosaic story of creation might, e.g., ask why plants were created before animals,
or how light could be created before sun and moon.

If we have a regularly varying difference-series, such as, e.g., the
colour-series, then not only do the members lie fast, but a certain continuity
also declares itself in the order of succession, in that the qualitative likeness
to the neighbouring colour is always greater than that to a colour lying further
off. And there can be analogous relations within the series, in that, e.g.,
violet may be said to stand to blue as orange to red. But above all the question
can come up how this determinate continuous order is to be explained. Goethe
believed he could find an explanation by keeping to the qualitative analogies.
But in spite of his protest, science has gone beyond this standpoint of quality
by exhibiting an analogy between the degrees of refraction of light and the
members of the quality-series. Instead of the analogy within the series, we get
analogy between two series, of which the one is quantitative, the other
qualitative. It will appear in what follows that an inference from analogy in
this form is of decisive significance for the modern age's apprehension of
existence.

If a series is symmetrical, that is, can be read in the reverse direction without
the single members being interchanged or any member eliminated, there is a new
possibility of analogy within the series. If (to use de Morgan's example)
\textit{A} is in conflict with \textit{B}, and \textit{B} with \textit{C}, there
is analogy between \textit{A} and \textit{C} with respect to their relation to
\textit{B}. The inference from analogy can be formulated in the second figure:
\textit{A} is an enemy of \textit{B}, \textit{C} is an enemy of \textit{B}. The
third figure can also be used: \textit{B} is an enemy of \textit{A} and is an
enemy of \textit{C}. But the analogy leads here to asking how the enmity between
\textit{A} and \textit{C} more precisely stands. It may be that the common
relation makes \textit{A} and \textit{C} friends. But it may also be that the
common enmity is not enough to overcome their mutual indifference, or the enmity
which for other reasons subsists between them.

With a progressive difference-series (that is, one in which each member stands to
its predecessor as its successor stands to it) and with partial identity-series
(that is, such as those in which each member is predicate to the preceding and
subject to the following, or consequent of the preceding and ground of the
following), not only are inferences from analogy like those already brought
forward possible, but inferences can also be drawn by the simple elimination of
intermediate members. A progressive difference-series is formed, e.g., by the
numbers, the times, the degrees and the places, which we came to discuss above
from another point of view. A partial identity-series is formed when predicates
themselves become subjects for new predicates, and the consequents grounds for
new consequents. The two kinds of series are characteristic of the formal
sciences (logic and mathematics). But, as we have already seen, such series also
acquire significance for purely qualitative difference-series. If it can be shown
that a progressive difference-series or a partial identity-series within the
circle of our subject-matters runs point for point parallel with a fixedly
ordered quality-series, so that to a determinate member in the one series there
answers a determinate member in the other, then from a transition within the one
series one can infer to a transition within the other. All exact empirical
science rests on such a setting together of quantitative and qualitative series.
Where such a setting together is possible, one can regard the quality-series as
if they were purely logical or mathematical series. One can, e.g., regard the
colour-series as if the difference between the colours consisted merely in the
difference between the degrees of refraction of light; and one can regard
unavoidable qualitative successions as if they were merely logical transitions
from ground to consequent (in a partial identity-series), or mathematical
transitions from one degree of refraction to another. But it is and remains
analogies that are built on here, even though there has here been made the great
advance, in relation to the inferences from analogy discussed above, that not
merely is occasion given for questions, but determinate expectations and surmises
are grounded --- hypotheses, that is, are formed, which can be sought confirmed
either by new members being looked for in the qualitative series, or by the
quantitative series being shaped with greater exactness than before, perhaps even
in such a way that a new quantitative series is constructed. Whereas certain
irregularities in the orbit of Uranus could be explained from Newton's theory of
gravitation, the irregularities in the orbit of Mercury cannot be explained from
that theory, and \textsc{Einstein} has then attempted an explanation from the
electromagnetic theory, from which the theory of gravitation itself would then be
derivable.\footnote{\textsc{Einstein:} \textit{Über die spezielle und die
allgemeine Relativitätstheorie}\textsuperscript{6}\textit{,} p.\ 70. ---
\textsc{Eddington:} \textit{Space, Time and Gravitation.} 1920. p.\ 123--125.}

Whether the two series will continue to run parallel cannot be known with full
certainty. However long a continuation of the quality-series may be, it does not
vouch for it that new subject-matters may not arise which do not order themselves
according to the rule stated in the quantitative series. Already \textsc{Leibniz}
saw that a complete verification is an impossibility. ``A hypothesis' success'',
he says,\footnote{\textit{Nouveaux Essais.} IV, 17, 5. (Op.\ phil.\ ed.\
Erdmann. p.\ 397.)} ``does not prove its truth''. To the complete confirmation
of a hypothesis there would belong what Leibniz calls a complete reversal (un
parfait retour) of the train of thought, that is, a demonstration that from the
quality-series one came back with necessity to this determinate quantity-series
and to no other. One would have to be able to reverse the inference and make its
earlier premisses the conclusion. An objectively correct inference can, after
all, be drawn from wrong premisses, and science's first (or last)
presuppositions are therefore difficult to establish.

A relation of identity is in any case not present here, though a dogmatic
tendency leads again and again to assuming one, in order to be able to rest in a
First (or a Last). Quality simply cannot be derived from quantity, even though
the study of the quality-series can lead to the setting up of proportionally
varying quantity-series. Quantity is itself a quality which (as with number,
time, degree, place) lends itself to the setting up of progressive
difference-series. Science aims everywhere, where it can be done, at forming such
series, in order then by their help to find lawfulness within other series whose
qualitative character cannot be reduced in that way. We shall later come to
discuss some particular examples of the great analogy on which all modern exact
empirical science rests, often without being conscious of it itself. We shall
also come to consider more closely Kant's great merit of clarity on this point,
as well as the misgiving that still remains in the face of his theory. Besides
Kant, it is especially the famous physicist Clerk Maxwell (in philosophical
respects a disciple of Kant at second hand) who has maintained the significance
of the concept of analogy in this connection. To him Ernst Mach has on essential
points attached himself. ---

In conclusion it is to be remarked that even where absolute identity-series can
be formed, we meet analogy as a presupposition. It was the Eleatic ideal to
reduce everything to identity. If that could be reached, all existence's
fundamental subject-matters could be presented in the series
\textit{A = B = C = D} and so on. But (as I have shown in detail in
\textit{Den menneskelige Tanke} p.\ 177--201) every absolute identity between
real subject-matters presupposes a determinate point of view in relation to which
they can be regarded as identical, and so substituted for one another. From the
point of view which the law for a group of subject-matters states, these
subject-matters can be regarded as identical, though they may have many other
properties than those the law concerns. It is the same relation as that between
examples of the same concept.


% ============================================================
%  III. ANALOGI MELLEM ERKENDELSESFUNKTIONER -- pp. 61--79
%       §§ 12--14 -- notes 30--40
% ============================================================

\chapter{Analogy between Functions of Cognition}
\label{ch:3}

\hnum{12} From a psychological point of view it becomes more and more impossible
to assume sharp oppositions between the different stages of mental life. There
shows itself a constant continuity between the different stages, and thereby an
analogy between the different stages, and thereby again between the different
mental functions, becomes possible --- even if we cannot, as there was occasion
to insist above, without more ado identify the functions of the higher stage of
development with those of the lower. There is here a connection which defies the
distinctions men have wished to maintain. If, e.g., the so often drawn
distinction between matter and form in mental life could be carried through,
there would be a limit both to analogy and to continuity. But sensation itself
already defies this distinction. That is seen already when we compare
sense-sensations with the corresponding physiological processes (so far as we
know them): what is physiologically a composite process comes forward in the
sense-sensation gathered together into a relatively simple element. We trace
here already the law of synthesis. But this element itself is determined from
the very outset partly by its relation to immediately preceding elements, partly
by its relation to elements appearing at the same time. And very quickly
reproduced elements are woven together with it. An absolutely simple and
independent sensation, which should be an absolutely passive (i.e.\ externally
determined) phenomenon, cannot be exhibited. Only approximately can such a
complete isolation be exhibited. \textsc{Salomon Maimon}\footnote{Cf.\
\textit{Den nyere Filosofis Historie}.\textsuperscript{3} III, p.\ 122.} already
objected against Kant's doctrine of sensations as mere ``matter'' that an
absolutely passive sensation was to be compared with $\sqrt{2}$, and so could be
exhibited only approximately. We thus find already in the simplest sensation an
analogy both to what at higher stages comes forward clearly as synthesis, and to
that distinguishing, that mutual determination of differences, which is peculiar
to reflection.

A transition between sen\-sa\-tion and re\-flec\-tion is formed by mem\-o\-ry- and
imag\-i\-na\-tion-intuition. In intuiting, not only relations of difference but also
relations of likeness make themselves felt; they condition the fusion or the
organization which every intuiting bears the stamp of.\footnote{The difference
between fusion and organization as two kinds of mental whole-formation I had
occasion to bring out in \textit{Den store Humor.} p.\ 15--34.} Memory and
imagination present, like sensation, an analogy to the way in which reflection
strives to bring connection between the subject-matters accessible to it. The
involuntary expectations which are the first form of memory and imagination form
an analogy to the hypotheses which reflection sets up by the way of analysis. ---
An intermediate form between intuiting and reflection comes forward in the
involuntary alteration of an intuitional image, by which it can become
articulated, can change direction and centre of gravity, without reflection
needing to play any part. There will, especially in some natures (those called
above the intuitive), be a tendency to hold fast to intuiting as long as
possible and to keep reflection's analysis away. It is an example of the frequent
confusion of identity with analogy when many psychologists read judgments into
the involuntary functions of sensation and intuition. The judgment becomes
necessary only when immediate intuiting cannot, not even through articulations
and displacements, hold all the subject-matters that draw attention to
themselves. And the judgment will always be an imperfect surrogate for the
intuition, since reflection can hardly separate out all the elements contained in
the intuition and bring them into the new form of connection which the judgment
is. The single judgment will therefore, for its supplementing, require new
judgments, so as thereby to get more decimals for the determination of the
irrational relation between intuiting and reflection. --- One and the same type
governs, psychologically regarded, the development of the life of cognition
through its different stages. More fully than there was room for here, I have
sought to demonstrate this in \textit{Den menneskelige Tanke.} p.\ 34--54.

It may have its importance to see that continuity, and an analogy conditioned by
it, make themselves felt as regards other sides of mental life too. Answering to
the opposition between sensation on the one side and memory, imagination and
reflection on the other, there comes forward in the domain of the life of
feeling the opposition between elementary and ideal feelings. Elementary
feelings hang together with vital sensations answering to alterations in the
circulation of the blood, respiration and digestion, or with sensations
answering to influences from without. Representations play no part, neither
recollections, imaginings nor reflections. Even with such feelings differences
of quality make themselves felt --- e.g.\ between a feeling of strength and a
feeling of lassitude, between a feeling of ease and oppression, between a feeling
of freedom and a feeling of anxiety. With ideal feelings, representations
(recollections, imaginings, reflections) play a greater or smaller part.
Thereby the states get a more determinate character, a more or less pronounced
direction. The indeterminate feeling of ease and freedom becomes joy in that
which rightly or wrongly stands as its cause; the indeterminate feeling of
anxiety becomes fear of something that stands as threatening, and so on. But
though we have here different stages in the development of the life of feeling
before us, there is a constant transition between them. Ideal feelings are never
independent of our organic states or of the sense-sensations that emerge at the
moment, but may precisely be predisposed by these. Common to everything we in
psychology call feeling is the difference between pleasure and displeasure; it
holds for the ideal as for the elementary feelings, and these are only symptoms
that feeling is present at all. No psychologist takes it into his head to derive
the special qualities that come forward in elementary and ideal feelings from
the abstract concepts of pleasure and displeasure. But it has an interest in this
connection to note the analogies between elementary and ideal feelings, as well
as the way in which the transition from the former to the latter takes place. To
apprehend as original the pronounced qualities which only the ideal feelings
present would conflict with what genetic psychology teaches, especially the
investigation of the life of feeling in the animal and in the infant.

In his interesting book \textit{Følelseslivets Grundelementer} (1920)
\textsc{Carl Jørgensen} has given a series of good descriptions of different
feelings. One notices the good observer all through. But he lays no weight at
all on the more indeterminate (and yet qualitative) feelings which I call
elementary, and he denies that representations play any part in the development
of the life of feeling towards more pronounced forms. And in his criticism of my
conception he overlooks (1) that I already assume, in the case of the elementary
feelings, certain differences of quality (which therefore are not due to the
influence of representations); (2) that in my terminology I understand by joy
\textit{joy in} something (a ``joy'' that was not joy in something I should
reckon under the feeling of power, of ease or of freedom); (3) that I do not
assume that it is due to representations that feelings acquire quality at all,
but only that under the influence of recollections, imaginings and reflections
there arise what I have expressly designated as ``\textit{new} qualities of
feeling''. (See the fifth edition of my Psykologi, used by the author, p.\ 304;
305; 340.) I cannot recognize that the author has taken pains to enter into my
conception; had he done so, he could not have christened it ``the pleasure- and
displeasure-theory'' and given it the appearance that I assume feelings of
pleasure and displeasure without any more special qualities whatever. ---

In the domain of the life of will (when we take the word will in its widest
sense) we also find a series of stages: need, impulse, deliberation, intention,
decision. The analogy between them rests on the active and concentrated
character of the mental states determined by them. (On this I refer to my
treatise \textit{Begrebet Villie} in the third series of my \textit{Mindre
Arbejder}.) ---

In all the domains of mental life, qualitative differences make themselves felt
between the different stages. The one stage cannot be reduced to the other; but
there are analogies between them which testify to the dominion of one and the
same type.

\hnum{13} It is one thing that within psychology it can be possible to find
continuity between the different stages of development in spite of their
qualitative difference, and that analogies can therefore be exhibited between
more elementary and more elaborated mental functions; it is another whether on
this point there is not a difference in principle between the psychological and
the epistemological consideration of thinking and cognition. In the opinion of
some, quite other principles must be applied when one asks about our cognition's
validity than when one asks about its actual constitution and development. I
shall not here come back to the general question of the relation between
psychology and epistemology, which I have discussed in \textit{Totalitet som
Kategori} (p.\ 3--20), but shall only make some remarks on the relation between
valid and actual thinking. It is clear, I think, that since it is men who think
the valid thoughts, the valid thoughts must also be actual thoughts, thoughts
that are really thought --- while conversely thoughts that are actually thought
are not always valid thoughts. Herewith the relation between psychology and
epistemology is given. However high scientific thinking may raise itself, however
many new means and forms of thought it may develop in order to solve its task, it
must all the same be psychologically possible, and psychology must, using the
history of the sciences, investigate its possibility and its actual development.
An anecdote which Kant somewhere tells on another occasion expresses the relation
in a good image. An Indian, on a visit to a European, saw the cork spring out of
a beer bottle and the froth foam out, and the Europeans were amused at his great
astonishment at it. But then he said: ``I am not astonished that it comes out,
but at how you got it in!'' Such is precisely psychology's position over against
science, as over against art and religion. Psychology can as little bring forth
science as it can bring forth art or religion. But it can and must set itself the
task of investigating how these forms of spiritual culture are possible. History
shows too that science arises out of the involuntary need for orientation in the
world. If we are in doubt whether we stand before something real or not, or as to
which road we are to strike into when we notice that we are astray, then we
already have use for formal science, that is, for inferences, measurements and
calculations, and for real science, that is, for the finding out of causal
relations, processes of development and whole-formations, as regards the
subject-matters before us. One of the most important differences between
prescientific and scientific thinking is, as we have already had occasion to
mention, that in the former involuntary expansion plays an essential part: one
asks (if one asks at all) not ``Why?'' but ``Why not?'' --- and lets things go as
long as they will, whereas scientific thinking is born of doubt and tests every
single step. Naturally a scientific genius applies the principles of strict
thinking involuntarily; in the genius, as in instinct, need and capacity are
inseparably joined.

Kant already drew attention to the fact that ``even the plain human
understanding'' thinks according to certain principles, and in the most recent
time it has been energetically maintained, especially by \textsc{Émile
Meyerson} with the support of the history of the sciences, that all science
develops out of sound human understanding (\textit{sens
commun}),\footnote{The Danish expression ``sund Menneskeforstand'' is better than
the German ``der gemeine Verstand'', the French ``sens commun'' and the English
``common sense'', inasmuch as in these last designations there lies something
current, common and traditional, which is not always ``sound''. \textsc{Léon
Brunschvicg} has aptly distinguished between ``bon sens'' as ``fonction de
discernement'' and ``sens commun'' as ``tradition de consentement universel'', and
finds this difference already drawn in Montaigne and Descartes.
(\textit{L'expérience humaine et la physique.} 1922, p.\ 574 f.).} correcting its
principles step by step where that is needed. There is not identity but analogy
between science and sound human understanding. As Meyerson puts it: ``The
concepts of sound human understanding have been brought forth by a process which
is indeed unconscious, but for the rest strictly analogous to the process by
which the scientific theories are shaped'' (Identité et
Réalité\textsuperscript{2} p.\ 393). The essential reason why science cannot stop
at sound human understanding Meyerson finds in this, that although the latter can
be led to seek and perhaps find lawfulness in what is lived through, it does not
satisfy the need to set something lasting, something constant, in place of the
more or less fleeting experiences; and he explains thereby the constant
re-emergence of the atomic doctrine in the history of science. (Op.\ cit.\ p.\
420--423.) But it would perhaps be just as correct to say that science does not
find the lawfulness at which sound human understanding stops special, exact and
comprehensive enough, and that the atomic theory keeps emerging anew precisely
because only this theory makes possible a quite exact representation of the most
special lawfulness. --- To this point of difference between the excellent French
epistemologist and myself I shall later have occasion to return. ---

By the following example the analogy between psychology and epistemology may
perhaps be marked. Psychology draws attention to the fact that for a
consciousness of time to develop, there must make itself felt in consciousness an
opposition between something continuously subsisting and something changing. If
one state were forgotten as soon as another state arose, consciousness of time
would not be possible. A simple sensation of change, what we may call a sensation
of time, presupposes a relatively fixed background against which the change can
be noticed. In the simplest case this constant background will be formed by the
feeling of life (the feeling answering to the given organic state), which can
maintain itself while different, more peripheral experiences succeed one another.
A representation of time proper (as distinct from the mere sensation of change)
presupposes that certain experiences recur after an interval, so that a rhythm
can be apprehended. Here a need or an impulse will be able to form the fixed
background, and a rhythm of privation, striving, attainment (or disappointment)
can be experienced. It is this practical triad that from the first carries the
apprehension of time. Practically, the interest in the rhythm of time goes before
the interest in the measurement of time, to which the history of religion and of
magic testifies.\footnote{\textsc{Hubert:} \textit{La représentation du temps
dans la religion et dans la magie} (Hubert et Mauss: Mélanges d'histoire des
religions), p.\ 195.} An objective consciousness of time develops through the
observation of regularly changing natural phenomena (day and night, the seasons
and so on) and forms, by the exact measurements astronomy makes possible, the
transition to the scientific apprehension of time. The alterations this
apprehension of time has undergone in modern times I have had occasion to discuss
in the treatise \textit{Relation som Kategori} (p.\ 62--77). It is especially the
question of the starting-point for the measurement of time that becomes decisive.
As a consequence of the Copernican theory, Bruno maintained that every star must
have its own time; and as a consequence of the electromagnetic theory of light,
Einstein has maintained that the measurement of time depends on whether the
system of co-ordinates laid at the base is relatively at rest or in motion. Even
with the most exact and special scientific theory of time and its measurement
carried through, the analogy to the primitive apprehension of time does not fall
away, just as it has developed out of it through the constant work of
centuries. ---

It is with the concepts of number, degree and place in a similar way as with the
apprehension of time. They too have developed, into the form in which science
now works with them, out of qualitative observations, and the analogy with the
initial stages has not fallen away. At all stages the principles for orientation
in the world of subject-matters are at work.

\hnum{14} The history of the sciences shows us remarkable examples of how certain
fundamental thoughts emerge again and again, though in such different forms and
with such different grounding that one is not seldom inclined to deny both the
historical connection --- which is nevertheless actually present --- and the
abiding significance they may have for continued research. There does after all
make itself felt a continuity in the course of scientific development, and,
connected with it, interesting analogies which point to a certain fundamental
type for human thinking and for the demands it is to satisfy.

a. The history of atomic theory shows us how through the ages two apparently
opposite tendencies make themselves felt in human thinking: on the one side the
need to find something that subsists in spite of all changes, and on the other
the need to presuppose separate starting-points for the changes that experience
demonstrates. Atomic theory arises through the co-operation of these two
tendencies. It is due to the experience that one does not reach real
understanding by stopping at the qualitative changes the senses show us.
Democritus therefore declared that the sense-qualities had only apparent, not
real existence. When one quality replaces another, something seems to arise and
something to be lost; but when we apply ``a finer organ of thought'' than
sensation, we reach a ``genuine cognition'' which shows us that the case stands
quite otherwise: nothing arises or perishes, but what happens is that certain
small parts group themselves in another way than before. These small parts are
homogeneous, for otherwise they could not act upon one another. It is in the name
of thought that Democritus maintains his atomic doctrine, and he was therefore
set alongside Plato by Sextus Empiricus, since both maintained that only what can
be thought (τὰ νοητά) can be true.\footnote{On an obscurity in Democritus'
doctrine of sensation and thought, see my treatise on Democritus and Plato.
(Mindre Arbejder. III. p.\ 98).}

Historically, Democritus' atomic doctrine has exercised a great influence. His
writings are, apart from some fragments, lost --- one of the greatest losses
within the literature of antiquity. But in the form Epicurus had given it, it was
presented in brilliant fashion in Lucretius' famous didactic poem; it influenced
Bruno in the Renaissance, and it was taken up again by Basso and Gassendi, from
the last of whom Newton took it over (in his Optics). From Newton, John Dalton in
turn took it over,\footnote{Cf.\ \textsc{Roscoe:} \textit{John Dalton,} p.\
127--145.} in showing, following Lavoisier's quantitative method, that the atom
of an element has a constant weight, and that composite substances arise through
combinations of different elements in determinate proportions by weight.

For the researchers of modern times there was this misgiving about Democritus'
and Newton's atoms, that they were absolutely hard and indivisible. Then force
would have to be lost at their collision. Huyghens and Leibniz therefore assumed,
as a consequence of their hypothesis of the subsistence of force, that the atoms
are elastic. Leibniz emphasized this very strongly against the Newtonian Clarke,
who, like others of Newton's disciples --- e.g.\ Maclaurin --- rejected the
subsistence of force and found precisely in the absolutely hard atoms a decisive
proof against it. Leibniz assumed that the force which the atom as a whole loses
at collision is preserved in its parts, which are set vibrating (agitation) by
the collision. Force is thus not lost, but is distributed over more points. It is
as when large money is changed into small.\footnote{\textit{Lettres entre Leibniz
et Clarke.} (Leibnitii opera philosophica ed.\ Erdmann, p.\ 775).} On this
utterance of Leibniz, \textsc{Henri Poincaré} has remarked that one could not in
any better way express the hypothesis that lies at the base of the mechanical
theory of heat.\footnote{See the quotation from Poincaré in \textsc{Léon
Brunschvicg:} \textit{L'expérience humaine et la causalité physique,} p.\ 227.}
Through this modification of atomic theory the transition is formed to the
apprehension of the most recent time, according to which atomicity is relative,
every atom being a whole little world-system whose elements are called electrons.
Democritus' ``finer thinking'' was thus not yet fine enough; but the science of
modern times has developed the fineness in high degree --- and why should
thinking not be able to become finer still?

\textsc{Émile Meyerson} has in his writings, and especially in an interesting
discussion in the French philosophical society (Bulletin de la Société française
de philosophie 3, mars 1910), exhibited with great expert knowledge the analogy
between ancient and modern atomism. He will not say that Democritus and Lucretius
have ``priority'' over modern chemists and physicists; even to call them
forerunners would, he thinks, be to misapprehend the true spirit of the history
of the sciences. And yet, he maintains, the setting together of ancient and
modern atomic theory is not arbitrary; a real analogy takes place, in that the
two theories spring from the same train of thought.

b. As it is with ancient atomism in relation to modern atomism, so it is on the
whole with ancient thinking in relation to modern thinking. The opposition is
clear, and often enough emphasized.\footnote{See, e.g., \textit{Den menneskelige
Tanke,} p.\ 123--134.} For Plato, the most pronouncedly Greek thinker, it was
science's task to form a system of concepts (``Ideas'') under which all the
subject-matters before us could be brought. In such a system one would have the
key to all existence, he thinks. He sought in his own way the identity behind
everything different and changing. The possibility of a classificatory system
within which all the different subject-matters found their place under ever
higher points of view of likeness was to him a testimony that the world of Ideas
is the true world. In modern times the concept of law has been put in the place
of the concept of Idea. Men do not flee from the differences and the changes, but
find in the laws for the existence of subject-matters and for the changes that
take place the identity which Plato found in the world of Ideas. Bacon and Bruno
already interpret the Platonic Ideas as the laws of subject-matters, and
Galilei's Platonic studies were of no small importance for his research. The
principle of identity rules in the classificatory systems of the ancient thinkers
as in the explanatory systems of modern times. In ancient science the qualities
are sought reduced to identities through concepts of species and of kinship. In
modern science qualities are apprehended as identities by the help of
determinations of quantity. Modern thinking too has its world of Ideas, namely in
the circle of points of view from which it becomes possible not only to classify
by the help of the principle of identity, but also --- by the help of the same
principle --- to reach rational explanations. The world of Ideas has thus, as it
were, sunk down into the world of experience; for it is no matter of course that
the subject-matters observation offers should be able to make classificatory and
explicative cognition possible. It might have been that existence was so
constituted that there was no possibility either of classification or of
explanation! But every victory, great or small, that science wins is a testimony
to the contrary.

The analogy between ancient and modern science also lays itself open to view in
this, that the difficulties which presented themselves for Plato's Ideas, and for
which his genius and his love of truth gave him so clear an eye (see the first
part of the dialogue \textit{Parmenides}), make themselves felt in new forms for
modern science's cognition of laws. For Plato the difficulty was this: when the
world of Ideas is closed in itself, the one Idea being derivable from the other,
and all Ideas finally from a highest Idea, how can there then subsist a world of
subject-matters whose whole nature is not exhausted by the elements that find
their place in the Ideas? In a corresponding way it stands as a constant problem
for modern science how the single subject-matters, in their concrete, individual
peculiarity, should be derivable from the general laws, if these are to be
science's last word or perhaps even to stand as the expression of existence's
innermost essence. Here is a problem that makes itself felt both in the domain of
natural science and in that of the human sciences. The problem sets itself in new
forms; but it always follows thought on this road.

c. Within the research of modern times a significant analogy makes itself felt
between the natural science of the seventeenth and of the nineteenth century. The
seventeenth century founds the thought-life of modern times. The principal
thoughts which later received their magnificent experimental confirmation come
forward here for the first time. The principle of identity was boldly applied to
the relation between cause and effect. Descartes and Spinoza declared that if
there is in the effect anything other and more than there is in the cause, then
full explanation has not been reached. Huyghens and Leibniz applied this thought
in their assertion of the subsistence of force: neither can new force arise at
the transition from rest to motion, nor can force disappear at the transition
from motion to rest. Against this thought the Newtonians \textsc{Clarke} and
\textsc{Maclaurin} asserted that it conflicted with
experience.\footnote{\textsc{Maclaurin:} \textit{Exposition des Découvertes de
Newton.} Trad.\ franç.\ 1749, p.\ 88.} But the thought was taken up again --- as
self-evident! --- towards the middle of the nineteenth century by the researchers
who at about the same time discovered the proposition of a relation of
equivalence between the different special forces of nature, according to which at
the transition from one form of energy to another energy neither arises nor
disappears. A study of these researchers' works shows that they were led by the
thought that there must be a relation of equivalence between cause and effect. It
stood for them as self-evident; but even the ``self-evident'' has its history, and
this determinate self-evident thing did not exist before the seventeenth century.
Fortunately the researchers in question \textsc{(Robert Mayer, Colding, Joule,
Helmholtz)} did not content themselves with a deduction from a ``self-evident''
proposition, but sought proofs from observations and experiments. A ``leading
thought'' is, after all, not a thought that can without more ado be used as a
major premise in a syllogism or as a presupposition for a deduction, but is a
thought that can serve as a model (``Paradigm''), as a foundation for an analogy
that can lead to the winning of cognition in new domains. The thought of the
subsistence of force has significance in and for itself, namely in the domain of
pure mechanics. But it was of course possible that its significance did not reach
beyond this domain, so that it could not help towards an understanding of the
mutual relations of the special forces of nature. This was the opinion of
positivism's founder Auguste Comte, for whom the different branches of science
were separated in principle from one another, so that a thought which can fit in
the one domain had no significance for the other domains. The different branches
of physics would, in his opinion, continue to stand separate from one another. He
did indeed assume that in the long run it would not be possible for the human
spirit to think strictly scientifically in one domain and unscientifically in
others; there will involuntarily be a striving after unity in method and theory.
But he ascribed no scientific significance to the analogy which here made itself
felt; he derived it rather from the influence of habit. The founders of the
proposition of the subsistence of energy, on the other hand, started out from the
possibility that there could objectively be exhibited relations between the
special forces of nature answering to the relation between motion and rest in
pure mechanics, and it was in reliance on this analogy that they went to their
work. --- We shall in what follows meet more examples of the significance of
analogy where the relation between different sciences is concerned; but the
analogy is not always so fruitful as in the case just discussed.

In my book \textit{Ledende Tanker i det 19.\ Aarhundrede} I thought I could
characterize this century as carrying through and verifying, in relation to the
thoughts of the preceding centuries, at whose appearance the science of modern
times came into being at all. In his fine and instructive presentation of the
history of physics and chemistry in the nineteenth century, \textsc{Niels
Bjerrum} thinks that I have been too modest on natural science's behalf. To this
I will say that it has by no means been my meaning to deny that the great
thought-work performed within natural science in the nineteenth century has an
independent and outstanding significance, and that much that is new and
surprising came forward in the course of that century. But for a comparative
historical consideration it is precisely new and surprising to see how thoughts
which hitherto had come forward only in the abstract forms of epistemology and
mechanics were now carried through in detail on the ground of experience and
experiment. By ``leading thoughts'' must first and foremost be understood the
directions in which thinking and research go. The direction is one thing, the
thought-work another. The very introducing of a new direction does, however,
require thought-work, and it may be a very difficult question to decide whether
the founding of a wholly new direction in research (as when the principles of the
new science were set up in the seventeenth century) requires less or more
thought-work than the carrying through of an already introduced direction in
detail. Like all work, thought-work too is the overcoming of a resistance, and
the resistance which thinking's first pioneers met was perhaps quite as great as
that which the nineteenth century's researchers, who carried through their
thoughts, met on their way.\footnote{The literary fate which the first works of
the four researchers mentioned above provisionally met testifies that here too
there was resistance enough.} The resistance has been of a somewhat different
kind, and a sure measure for it we do not have. It was in any case not my task to
investigate this point. I find the nineteenth century's greatness in the
energetic faithfulness with which leading thoughts that had already seen the
light beforehand were held fast and carried through.

d. Within the history of philosophy, \textsc{Hans Brøchner} has in his book
\textit{Bidrag til Opfattelse af Filosofiens historiske Udvikling} (1869) given
interesting contributions to the elucidation of analogy's significance for
philosophical thinking. He finds a connection between the great periods in the
history of philosophy, regarded as wholes, and the single attempts of thought
within each of them, in that these attempts of thought now anticipate what comes
forward in fully pronounced form only when the period has reached its high point,
now continue to bear the period's stamp after it has properly closed. Brøchner
assumes a mysterious influence of a period upon the single thinkers; but what
lies before us is plainly analogies between attempts of thought which take place
under approximately the same historical conditions, nearer to or further from
those conditions' culmination.

The great periods which the history of philosophy presents are often introduced
by the putting forward of views which are worked out more closely only later in
the period. Thus in the oldest Greek philosophy there are anticipated the
thoughts that first came to full clarity in Plato. At the beginning of the
Renaissance, attempts are made again and again at what Giordano Bruno first gives
in systematic connection. In many respects Kant anticipates the thoughts of
Romantic idealism. --- On the other hand, thoughts that have been fostered within
a period as a warranted expression of its special tasks can work hamperingly and
limitingly on thinkers whose personal striving would in and for itself have had
to carry them beyond the period's standpoint. The possibilities which presented
themselves in Aristotle and Plotinus for a standpoint other than the classical
Greek one were hampered by presuppositions common to all Greek philosophizing.
But such a hampering influence from the past can often call forth more deeply
going thoughts than would otherwise be possible. Thus Greek thinking has
exercised a beneficial influence in the Middle Ages, and the thinking of the
Middle Ages in modern times.

Brøchner regards the history of philosophy as a world closed in itself. The
relations become more involved when it is taken into consideration that the way
in which philosophizing is done is at any given time conditioned not only by the
philosophizing of preceding times, but in high degree also by new experiences and
by the work that takes place within other sciences. But even within its
determinate limitation, Brøchner's fine and spirited consideration has its
interest for the elucidation of analogy's appearance in human thinking.


% ============================================================
%  IV. ANALOGI MELLEM ERKENDELSESOMRAADER -- pp. 79--128
%      §§ 15--21 -- notes 41--77.  The longest chapter, ~82.5k characters, and
%      the only one whose divisions carry real headings (16--21).
%      122 of the book's 179 small-caps spans fall here.
% ============================================================

\chapter{Analogy between Domains of Cognition}
\label{ch:4}

\hnum{15} To the different domains of cognition, that is, to the different groups
of sciences, there answer different groups of categories. Common to all sciences
are the fundamental categories: synthesis and relation, continuity and
discontinuity, likeness and difference; and they are at the same time common to
prescientific and to scientific thinking. To logic and mathematics there answer
the formal categories: identity, negation, rationality. To the natural and the
human sciences there answer the real categories: causality, totality,
development. To the ethical sciences there answer the ideal categories: the
value-concepts.

The analogy between the prescientific and the scientific consciousness I have
already discussed. I now go over to exhibiting analogies between the formal
sciences among themselves, between formal and real sciences, between real
sciences among themselves, and finally between formal and real sciences on the
one side and the ethical sciences on the other.

\subsection*{16. Logic and Arithmetic}

Every science must itself, through analysis of its subject-matters and its
methods, fix its necessary fundamental concepts in as simple and as fruitful a
form as possible. But thereafter the question can be raised from the
philosophical side whether the fundamental concepts thus set up do not in turn
presuppose still simpler concepts, and, if they do, whether they can then without
more ado be derived from these or must be set up as new, more special concepts. A
concept may very well presuppose another concept without being derivable from it.
As regards the relation between logic and mathematics, the question becomes this:
whether they lie on the same line, resting on the same foundation, so that
mathematics signifies a straightforward continuation of logic; or whether
mathematics comes forward with new fundamental concepts, between which and the
purely logical fundamental concepts there can subsist only analogy, not identity.

Purely logical thinking would, as shown in the foregoing, stand powerless if all
the subject-matters before us formed either a chaotic difference-series or an
absolute identity-series. It is powerless both where there are no real likenesses
at all and where there are no real differences at all. Even under these, logically
regarded, desperate conditions, mathematical thinking can take one step more ---
and a step that can draw many steps after it. The differences or the identities
can, namely, be counted; that is, a series can be formed between whose members
there is identity of difference, a series that can always be continued according
to the law of identity of difference, the same law by which the series begins,
and it is indifferent whether the actually given series of differences or of
identities is limited. The number-series is due to a construction which consists
merely in the repetition of one and the same act of thought: every new member is
found by positing a difference identical with the difference between earlier
members in the series. The concept of number is built on a point of likeness
between real differences (namely that they are objects for thought or attention)
and a point of difference between real likenesses (namely that they are objects
for different acts of thought). --- The prescientific consciousness already
attempts to form such series. The child ``counts'' its skittles (``one more, one
more, one more'' and so on) until it grows tired of it; the trouble is only with
the summing up. The savage counts on his fingers and sums up by closing his hand.
Even if no more subject-matters are given, the operation can be carried out by the
help of self-created signs. Through such productive repetition the doctrine of
number becomes something other and more than formal logic. From the principle of
identity the possibility of always adding new members to those already given
cannot be derived. The principle of identity is applied logically to qualities as
to quantities; but through productive repetition no new qualities arise; there is
formed only a quantity-series or, as it can most simply be expressed, the pure
concept of magnitude comes into being. Arithmetic's fundamental concepts are thus
more special than pure logic's, and they can as little be derived from these as a
species-concept can be derived from a genus-concept at all. If one wished to say
that a number is a general concept or a class-concept, like the other concepts
logic operates with, there is all the same this important difference between
number-classes (e.g.\ five-classes, six-classes and so on) and other classes
(e.g.\ the class of horses): that whereas the definitions of the latter contain
qualitatively different elements, the definition of a number-class contains only
mutually identical elements.\footnote{Cf.\ already \textsc{Plato:}
\textit{Staten} VI, p.\ 526 A, where it is said of the arithmetical units:
ἴσον τε ἕκαστον πᾶν παντὶ καὶ οὐδὲ σμικρὸν διαφέρον, μόριον τε ἔχον ἐν ἑαυτῷ
οὐδέν.} The concept horse has for its content: ``odd-toed hoofed animal with a
specially developed middle toe and long metatarsus, with a complete set of teeth
and with a bony ring about the eye''. The concept 6 has its content given in the
series 1 + 1 + 1 + 1 + 1 + 1, which arises through the repetition of an identity
of difference --- a repetition we can ourselves carry out without awaiting
observation. For logic the number-series is only an example of a peculiar kind of
series-formation, an identically varying difference-series or also a progressive
difference-series (according as the talk is of ordinal or of cardinal numbers).
It is indifferent for logic how such a series comes into being; it takes it as a
fact that it can be formed. But what for logic is a single example is for
arithmetic of fundamental significance. The number-series is grounded on a
postulate, in that the possibility of a construction such as the one described is
asserted. It is through this postulate that we come from logic over into
arithmetic. And the same process which leads to forming the number-series leads to
continuing it through negative, fractional, irrational and imaginary numbers.

This postulate, which contains what may be called ``mathematical induction'' or
``the law of repetition'' (loi de récurrence), was applied in purely mathematical
form by Maurolyco (16th cent.) and by Pascal (17th cent.).\footnote{\textit{Œuvres
de Pascal,} éd.\ Brunschvicg. III, p.\ 456; VIII, p.\ 363 (with Pierre Boutroux's
notes).} \textsc{John Locke} applies it as a foundation for the concepts of
number, space and time (Essay II, 12, 8 and 13, 1--2). As characteristic of
mathematics in opposition to logic it has been apprehended by Leibniz, Kant,
Boole and Poincaré.\footnote{Cf.\ \textit{Totalitet som Kategori,} p.\ 33--39.}
Plato's doctrine of mathematics as lying between pure logic (``dialectic'') and
the world of observations (Staten p.\ 511 A. D; 525 A; 527 B) is thereby
confirmed.

In the most recent time \textsc{Hilbert} has declared himself not only against
those who, like Frege and others,\footnote{So B.\ \textsc{Russell:}
\textit{Introduction to mathematical philosophy.} Chap.\ 2.} apprehend the
doctrine of number as a part of logic, regarding numbers as species-concepts, but
also against those who, like Poincaré and others, build it on the capacity to
repeat an act of thought. This last conception Hilbert designates as mystical. In
supposed opposition to it he maintains\footnote{\textit{Neubegründung der
Mathematik.} (Abhandlungen des mathematischen Seminars der hamburgischen
Universität. I, 2 (1922), p.\ 161--163.} that one must begin with ``certain
discrete objects, given in an intuitable way prior to all thinking''. Such
objects we have in the signs, the doctrine of number's proper objects: ``Am Anfang
ist das Zeichen!'' The sign 1 is a number. And a sign which begins with 1 and ends
with 1, in such a way that in the interval after each 1 there always comes a +
and after each + a 1, is also a number. Thus we get the series 1 + 1 + 1 + 1 ....
In order to be able more easily to communicate with others, we then form the sign
2 for 1 + 1, 3 for 1 + 1 + 1 and so on.

I cannot see otherwise than that this arrangement leads back to the conception
which has been maintained by Poincaré and his mathematical and philosophical
predecessors. A sign and a series of signs come into being, after all, only
through acts of attention or of thought, and a series such as the one set up
presupposes the possibility that such an act can be repeated in one and the same
way (which naturally, psychologically regarded, is only approximately possible).
There is nothing mystical in the fact that we can repeat our thoughts or, if one
will, our signs in such a way that the series is constantly continued according to
the same rule which led to the setting up of the two first members.

In this, however (and it is in this connection the most important thing), Hilbert
agrees with the other group of researchers: that mathematics cannot be derived
from logic. A sign-language of the kind he introduces logic has no use for.

There rules, then, no identity between arithmetic and logic. But analogy does
rule. Number can stand to number as ground to consequent, and Hilbert accordingly
has use for the logical sign → as an expression for the relation between ground
and consequent. (It is the same sign that I used in \textit{Den menneskelige
Tanke} p.\ 168 to denote the relation between the members in a partial
identity-series.) The purely logical relations are found again in the mathematical
domain, and that in so uniquely clear and exact a way that the best and most
fruitful examples of logical inferring lie before us here. But they are all the
same only examples among other examples. And within the doctrine of number there
is analogy between algebra and arithmetic, in that in algebra one fastens
attention only on the way in which the magnitudes are combined, without knowing
the magnitudes themselves. Thus (\textit{a} + \textit{b})\,\textit{c}\textsuperscript{2}
expresses all cases in which a magnitude is determined by the product of the sum
of two magnitudes and the square of a third magnitude, indifferently as to what
values \textit{a}, \textit{b} and \textit{c} have in the special
cases.\footnote{\textsc{Pierre Boutroux:} \textit{L'idéal scientifique des
mathématiciens.} 1920, p.\ 84.} We find again here the logical relation between a
concept and the examples for which it holds. Every example must be substitutable
for another everywhere, if the concept is to hold. In the doctrine of number this
substitution is unconditional; with concepts of qualitative subject-matters it
holds in each single case only from one determinate point of view.

The relation between logic and arithmetic becomes analogy; on the foundation that
comes into being through the construction of the number-series by productive
repetition, a new and more perfect application of the logical relation between
ground and consequent is possible; but that construction cannot be derived from
pure logic as such. There is given a new contribution in cognition's development
with the arising of the concept of number.

\subsection*{17. Arithmetic and Geometry}

Not only the doctrine of number, but geometry too rests on the postulate of the
possibility of repeating one and the same act of thought. Whatever an act of
thought's content or object may be, I can always add a new act of thought to the
earlier ones. I can think of more horses than those I have seen, and more
``signs'' than those I have already written down, more line-segments than those I
have already drawn. But on closer investigation men have all the same been led to
draw an ever sharper distinction between the doctrine of number and the doctrine
of other subject-matters, whether these are geometrical figures or concrete
objects. What is peculiar to number is its discrete nature. It is determined
merely by being a unit alongside other units, or by consisting of units that have
this property. But in the subject-matters that are counted there are always other
properties besides that of being countable. That is already the case with
geometrical lines. In and for itself one could, instead of expressing units by
vertical strokes (1 + 1 + 1 + 1 ....), use horizontal strokes (--- --- --- ---
....); and when then the one immediately continued the other, a straight line
would be constructed. This arises then, to that extent, by the same ``Synthesis
des Gleichartigen'' as the number-series. \textsc{Kant} has thought of the matter
in precisely this way when he says: ``Wir können uns keine Linie denken, ohne sie
im Gedanken zu ziehen'', and when he maintains that we ``im Ziehen der geraden
Linie bloss auf die Synthesis des Mannigfaltigen Acht haben''. (``Kritik der
reinen Vernunft''\textsuperscript{2} p.\ 154.) There does after all take place an
addition at every moment while we draw the line, just as every time we add a unit
in the number-series. But the great question is whether we have the right to say
that the constant additions, while we draw the line, are equal. With the
number-series this question did not arise, and for the simple reason that each
single member is regarded merely under the point of view of an addition --- solely
as something that is added, without there being any talk of extension or of other
properties.

This difference in principle between the doctrine of number and geometry has
worked its way forward successively. Greek mathematics had its foundation in
geometry. Relations of number were made intuitable by relations of space --- the
product of two numbers, e.g., by the rectangle of two lines representing the
factors. With Descartes came a turning-point. Not only were figures now
conversely expressed by numbers and equations stating the relation between the
figures' elements, but geometrical problems were reduced to algebraic problems.
There was, as \textsc{Zeuthen}\footnote{\textit{Matematikens Historie.} II, p.\
285.} puts it, a complete reversal of the traditional relation between algebra
and geometry, and thereby is marked a reform in the whole of mathematics, which
hitherto had had its acknowledged foundation in geometry but from now on got it in
algebra. A step further went \textsc{Maxwell} (1855),\footnote{\textit{Über
Faraday's Kraftlinien.} Übers.\ von \textsc{Boltzmann} (1895). Einleitung p.\ 4.}
when he declared that all scientific applications of mathematics rest on relations
between the laws of physical magnitudes and the laws of the whole numbers. When
\textsc{Boltzmann}, in a note to his translation of this passage in
\textsc{Maxwell}, says that so far as he knows the standpoint that measurements of
magnitudes of space and time by numbers rest merely on analogy with the relations
subsisting between the numbers has not since been taken up, this no longer holds.
Several researchers have continued \textsc{Maxwell}'s train of thought. Thus
\textsc{Ernst Mach} (who in general has drawn attention to \textsc{Maxwell}'s
epistemological way of looking at things) lays great weight on the limits of
measurement's exactness. ``The coincidence of a measuring-rod and the
subject-matter that is to be measured cannot'', he says,\footnote{\textit{Erkenntnis
und Irrtum.} (1905), p.\ 328 f.} ``be established with unlimited exactness''. And
the step has been taken right out by \textsc{Hjelmslev} in his presentation of
analytic geometry.\footnote{\textit{Elementær Geometri.} Tredie Bog. 1921, p.\ 71;
91; 180 f.} \textsc{Hjelmslev} distinguishes sharply between the geometry of
reality, which deals with things, and abstract geometry, which deals with numbers.
In a chapter on the theory of geometrical measurement he shows that no measurement
of a line-segment presented in a drawing can be wholly exact; not one but several
results are obtained, between which a choice must be made if the line-segment is
to be ``fixed''. Drawings are constantly necessary for guidance and description,
but they can be apprehended only as symbols of the abstract lines which can be
expressed by determinate numbers.

\textsc{Zeuthen}\footnote{\textit{Hvorledes Matematiken i Tidsrummet mellem Platon
og Euklid blev rationel Videnskab} (1917), p.\ 9.} regarded, in the last writing
we had from his hand, the ideal or abstract figures --- whose concepts are formed
by detachment from what sensation shows --- as symbols, just as the letters and
operation-signs of modern mathematics may be said to be. Such ideal or abstract
figures can now, according to \textsc{Hjelmslev}, be expressed only in numbers. In
an interesting way the two Danish mathematicians have thus each seized his own
side of the correlation which the concept of symbol contains. When \textit{A} is a
symbol of \textit{B}, \textit{B} is namely also a symbol of \textit{A}.
\textsc{Hjelmslev}, whose conception of geometry is most nearly to be called
realistic, apprehends the real geometrical forms as symbols of the abstract forms,
which according to him are expressed only by numbers. He can say with
\textsc{Goethe:} ``Alles Vergängliche ist nur ein Gleichnis!'' But
\textsc{Zeuthen}, whose conception of geometry is most nearly to be called
idealistic, since he maintains the necessity of ideal definitions in geometry,
conversely regards these ideals as symbols of the real geometrical relations. ---

There lies here an extraordinarily interesting example of how analogy can be the
only way in which the relation between different domains of cognition can be
marked. And here, as at other points, the development of thinking has led to
men's finding analogy where they formerly found identity.

\subsection*{18. Formal and Real Science}

Analogy will now show itself to be the right expression not only when the
relation between the different formal sciences (logic, arithmetic, geometry)
among themselves is to be determined, but also when the relation between formal
science and real science is to be determined. --- We keep provisionally to
natural science as the representative of real science, and go over later to
discussing the relation between natural science and the human sciences.

It is the relation between the concepts of ground and cause that is here in
question. On this point lies \textsc{Kant}'s great deed. He has set analogy as
the expression for the relation between ground and cause, in place of the
presupposition of these concepts' identity from which dogmatic philosophers and
investigators of nature had started out. The relation between cause and effect
signifies, according to \textsc{Kant}, only that the succession in time of two
subject-matters (events) is analogous to the logical-mathematical relation
between ground and consequent. Real science works without interruption to
demonstrate that every change that occurs in the world, everything new that
appears for our observation, stands to preceding changes as an inference stands
to its premisses, that is, as consequent to ground. All mystery in the causal
relation thereby falls away. We know nothing of this relation except that it is
at once a relation of ground and a relation of time. Schematically we have here
an analogy between two series: a partial identity-series (\textit{A} →
\textit{B} → \textit{C} → \textit{D} . . .), in which each member stands to the
following as ground to consequent; and an indeterminately varying
difference-series (\textit{a}\textsubscript{μ} \textit{b}\textsubscript{γ}
\textit{c}\textsubscript{π} \textit{d} . . .), where the members' order of
succession is merely chronologically determined. From a member in the first
series one can infer to a determinate member in the second series, and thereby,
within the second series, from one member to a following one, even if one thereby
gets outside the members actually before us. A cause is a subject-matter which
appears for our observation at a certain point of time, and from which it can be
inferred that a certain other subject-matter will appear at a later point of
time. It was through his study of Newton's physics, and under the influence of
the sting Hume's criticism of the concept of cause had left behind, that
\textsc{Kant} arrived at this conception.\footnote{The clearest passages in
\textsc{Kant} are \textit{Kritik der reinen Vernunft}\textsuperscript{1} p.\ 112
and \textit{Prolegomena,} § 30.} By virtue of the analogy between a
logical-mathematical relation and a relation of succession, experience becomes
something other and more than a collecting of single observations. This new
concept of experience \textsc{Kant} sets in opposition to Hume's conception of
experience as a sum of isolated, mutually independent observations --- that is,
as a chaotic difference-series, Spinoza's \textit{experientia vaga}. But his
exposition stands in high degree in need of closer determination and
simplification\footnote{In the third edition of \textit{Den nyere Filosofis
Historie} I have already (III, p.\ 278 f.) given the grounds for the necessity of
a simplification of \textsc{Kant}'s exposition. But there I kept mostly to the
purely formal side.} before it can be taken up as a member in our own time's
epistemology.

a. The causal relation itself \textsc{Kant} discusses in the section of the
``Kritik der reinen Vernunft'' headed ``the Analogies of Experience''. But two
sections go before it, called ``the Axioms of Intuition'' and ``the Anticipations
of Perception'', and here \textsc{Kant} does not seem to think he has use for the
concept of analogy. --- ``The Axioms of Intuition'' lays it down that all
phenomena have extension in space or time or both. What \textsc{Kant} here
discusses is nothing but an application of what has above in this treatise been
called an identical quality-series, which can be constructed by the
logical-mathematical way, and whereby the concepts of number, place and time
acquire an ideal form, since there is identity of difference between the members
of the series --- between the numbers, the places and the moments each for
itself. But the expression analogy fits better here than the expression axiom.
The meaning is that every subject-matter must be referable to, or able to enter
as a determinate member in, one of these series. In spite of all special
differences of quality, these series hold for everything about which we are to
have strict experience. But between the series of the special differences of
quality and these series there is only analogy, not identity. --- And it stands
in a corresponding way with ``the Anticipations of Perception''. Here the talk is
of the series of degrees. Every subject-matter that is to be able to be an object
for strict experience must be able to enter as a member in a series within which
each member, as regards intensity, always has the same difference from its
predecessor --- that is, a progressive difference-series, like the series of
numbers, of places and of times. Since the actual subject-matters present many
other properties than intensity, the relation between the series of degrees and
the series of actual subject-matters can only be analogy. But it is apt when
\textsc{Kant} here speaks of the anticipations of perception; only he ought to
have used this expression of all his epistemological principles. Analogy can lead
to no more than a grounded expectation (i.e.\ a hypothesis), which then possibly
is confirmed by new observations and experiments.

b. \textsc{Kant} sets up three analogies of experience. --- The most important is
the second, which concerns the causal relation in its generality. Here it is
shown, as already discussed, that a causal relation is a relation of succession
between two subject-matters which presents an analogy with the relation between
ground and consequent. --- The first analogy maintains that in a causal relation
no increase or decrease occurs in what subsists. It is ``the Analogy of
Substance''. If it were to be set up at all, it ought to have been placed after
the general causal proposition, which it plainly enough presupposes. What it
demands is that the relation between ground and consequent which the causal
proposition presupposes shall be reversible, so that the consequent can become
ground and the ground consequent --- to which it answers that what was first
effect can be cause of what was before cause, and conversely, so that there is a
relation of equivalence between cause and effect. But such a proposition can be
set up only as anticipation and as ideal. The history of science shows that
everywhere it is possible, men seek to exhibit such a relation of equivalence.
Thereby one approaches a pure relation of identity to the degree that is possible
when one has to do with qualitatively different subject-matters; and one then at
the same time abstracts from the difference of time which in each single case is
present and may be of decisive importance for the direction in which the
following changes take place. But not all inferences can be reversed, any more
than all predicates can be made subjects for their earlier subjects as
predicates. If \textsc{Kant}'s ``analogy of substance'' holds, it is in any case
a particularly happy case of the causal analogy. And it is an unfortunate
expression that \textsc{Kant} uses when he here speaks of substance as a special
category. He does not take the word substance in the old, metaphysical sense, but
in a relative sense, as he himself expressly declares (Kritik der reinen
Vernunft\textsuperscript{1}, p.\ 525): ``Was in der Erscheinung Substanz heisst,
ist nicht absolutes Subjekt, sondern beharrliches Bild der Sinnlichkeit und
nichts als Anschauung, in der überall nichts Unbedingtes angetroffen wird''. The
concept of substance in the old, absolute sense is a dying
category.\footnote{Cf.\ \textit{Den menneskelige Tanke,} p.\ 217 f. --- On
Spinoza's concept of substance see \textit{Spinoza's Ethica,} p.\ 6; 19--20.
Cf.\ besides above, § 7 at the end.} --- As regards the third analogy, which
maintains that everything simultaneous must stand in reciprocal action, the
matter is quite simple, for reciprocal action presupposes causality. If we are to
understand, starting from a given subject-matter, the simultaneous existence of
another subject-matter, there must through more or fewer intermediate members be
a relation of causality between them or between their causes. We must, starting
from the one subject-matter, by the requisite orientation come upon the other.

c. All three analogies together assert, according to \textsc{Kant}, that all
subject-matters lie in one nature. By ``nature'' \textsc{Kant} understands
precisely ``the connection of subject-matters, as regards their existence,
according to necessary rules, that is, according to laws''. (Kritik der reinen
Vernunft\textsuperscript{1}, p.\ 216.) \textsc{Kant}'s concept of nature, which
coincides with the strict concept of experience, thus maintains analogy between
the series of successive and simultaneous subject-matters and series of ground
and consequent and of numbers, times, places and degrees. If experience is to be
possible, the subject-matters must be capable of being ordered in accordance with
such series.

d. While \textsc{Kant} carefully constructs the ``schematisms'' --- that is, the
progressive difference-series whose application to the subject-matters before us
is to make experience in the strict sense possible --- and while he carefully
distinguishes between the analogy which thereby plays a part and identity, he has
not seen that only an advancing, approximating application of this analogy is
possible. His concept of nature or of strict experience is rightly formed, but he
did not see that ``experience'' in this sense is an ideal, a great hypothesis
which properly can never be fully confirmed. It holds here too that a hypothesis'
success does not prove its truth. The great question (which Hume raised, and
which \textsc{Kant} thought he had answered by his definition of experience) is
whether all subject-matters that can be apprehended now and in the future permit
an application of those ideal series. \textsc{Kant}'s ``analogies of experience''
formulate in a way of genius the great hypothesis on which natural science has
worked for three hundred years, and which has again and again shown itself
fruitful. And even if he did not see the hypothetical character of his concept of
experience, it was all the same a great achievement that he saw and formulated so
clearly the principle that lies at the base here. Therefore he still marks a
principal point and a central point in the epistemology of modern times.

e. It was no wonder that a conception such as \textsc{Kant}'s could not at once
prevail. On the contrary, the connection between formal and real cognition which
he had sought to exhibit was dissolved. The tendency to reach absolute identity
as at once the highest form of thought and the highest result of thought, the
wish once for all to get beyond everything chaotic and manifold to rest in the
thought of something that does not again point beyond itself, was --- especially
under the influence of poetic and religious currents in the spiritual life of the
time --- too strong for men to be able to content themselves with \textsc{Kant}'s
apparently modest analogies. The realities were to be wholly taken up into the
``Idea'', into the pure form of thought. In German Romantic speculation this
tendency comes forward. \textsc{Goethe} and \textsc{Schelling} maintained that
the explanation of the qualities was to be found within the world of qualities
itself, certain qualities being laid at the base as ``primal phenomena''.
\textsc{Hegel} attempted to derive the world of qualities and of time from the
world of pure logical identities --- an attempt which at once, by its confidence
in the power of pure thought, bore the stamp of Romanticism, and pointed beyond
Romanticism, since \textsc{Hegel} acknowledged that full understanding cannot be
reached by stopping at the qualities.\footnote{In his spirited characterization
of Hegel's philosophy of nature, \textsc{Émile Meyerson} lays at the base a point
of view similar to the one applied here. Cf.\ \textit{De l'Explication dans les
Sciences,} II, p.\ 78: ``Hegel a réellement tenté une entreprise qui constitue un
postulat éternel de l'esprit humain''.} Every science that will give absolute
solutions must end in an attempt like \textsc{Hegel}'s. If those are right who
maintain natural science to be the absolute science, they must make the attempt
to derive the qualitative differences from the logical-mathematical principles by
whose help the world of qualities becomes intelligible to us. If one will not
acknowledge that strict real science reaches its goal through exact analogies
between quality-series and mathematical-logical series, one must make
\textsc{Hegel}'s attempt over again and crown modern natural science's striving to
explain all changes as motions whose laws can be formulated as equations, with an
attempt to derive the differences of quality from the identities lying at the
base in the equations. Then there would be given a proof that natural science had
solved the riddle of the world. The ``parfait retour'' which Leibniz demanded
would have been reached. Such an ideal surely stands in the background of many
investigators of nature's train of thought. If one asks what then becomes of the
qualities, which are after all the constant and immediately given, the answer runs
that they are only ``anthropomorphisms'' --- the same answer the theologians give
when it is asked how the psychological expressions used of the Godhead's essence
and working are properly to be understood. But in both cases the
``anthropomorphism'' sets a new problem. Qualities such as colour, sound and so
on, and such as thinking and willing, simply are subject-matters that lie before
our cognition, and the question is whether they can ever be absorbed, wholly taken
up into rational knowledge (whether this comes forward as a series of equations or
as a series of dialectical transitions). We can, as the history of natural science
shows in brilliant and unique fashion, build up a world of thoughts that can guide
us in the world of ``anthropomorphisms'' in which we live. And this guidance
natural science gives us, if the consideration set forth in the foregoing is
correct, by the help of exactly carried through analogies. --- Whether theology
can give us a corresponding guidance I leave here undecided. Attempts at a strict
carrying through of the analogies applied lead in quite another direction than
positive theology.\footnote{Cf.\ my \textit{Religionsfilosofi,} § 16--25.}

The impossibility of deriving differences and changes of quality from the
equations that mark the formal foundation of natural science's exact analogies
--- and in which there is a certain inclination to see the true reality --- is
seen already from this, that the formation of the identical quality-series (the
series of numbers, times, places and degrees) takes place by an analysing
purging-out of the empirical quality-series. From the concrete wholes before us
one takes out the formal elements which permit strict logical-mathematical
treatment. In strict science's exact results it is always only a part of reality
that is expressed, even if it is the part of reality which by the way of analogy
can lead to the understanding and the mastering of the other parts or sides of
reality. It is always nature that is explained by nature itself, in that the most
transparent relations are laid at the base for the interpretation of other
relations. Or, to use \textsc{Kant}'s definition of ``nature'': that which makes
nature ``nature'' is used to interpret the parts of nature which do not yet
satisfy the strict concept of nature. In contrast to the primitive mentality, the
basis of interpretation is now definitely brought out, and it is no longer any
foreign authority that is consciously or unconsciously invoked. But just as
little as cosmology (metaphysics), which always wishes to explain the whole of
existence from a part of it (from ``matter'', from ``spirit'' and so on), has ever
been able to reach a rational conclusion, just as little will any science be able
to do so.

In the most recent time the Marburg school has apprehended the elimination of the
qualitative elements as the goal of scientific cognition, regarding only
quantitative science as true science. Its attempt is a companion piece to
Hegel's, only that a system of equations takes the place of speculative
dialectic.\footnote{Cf.\ \textsc{Gustav Körber:} \textit{Der Marburger Logismus
und sein Verhältnis zu Hegel.} (Archiv für die Geschichte der Philosophie. 1917).}
In my treatise \textit{Totalitet som Kategori} (p.\ 46--48) I have had occasion
to state my position towards this tendency, which is one of the most energetic
and most significant that have developed since the taking up again of the problem
of cognition in the last half of the nineteenth century.

With other epis\-te\-mol\-o\-gists Kant's thought has come more into its own. Pos\-i\-tiv\-ism's
founder \textsc{Auguste Comte} laid great weight on the difference between the
different sciences, which he ordered in a series beginning with mathematics and
in whose following members special and complex qualities played a greater and
greater part. The last member in this series is sociology. Between the different
sciences there can be only analogy, not identity; and \textsc{Comte} even
misapprehended, as we have had occasion to see, such reductions as were
scientifically warranted. But his classification of the sciences, which advances
from the simple to the complex, follows an ordering similar to the one applied
shortly afterwards by \textsc{Renouvier} in the doctrine of
categories,\footnote{Cf.\ \textit{Relation som Kategori,} p.\ 13.} and which is
also followed in my own exposition. It holds for categories as for sciences that
the more complex cannot be derived from the more elementary, but that there may
all the same subsist instructive analogies between them. With full consciousness
and great clarity the point of view of analogy was maintained towards the middle
of the nineteenth century by \textsc{Clerk Maxwell}, as we have already had
occasion to see as regards the relation between arithmetic and geometry. Analogy
expresses in general, for \textsc{Maxwell}, the relation between formal and real
science. Attaching himself to him, another philosophizing investigator of nature
worked, \textsc{Ernst Mach}, whose epistemology is built above all on an
independent study of the history of natural science.\footnote{\textsc{Maxwell}
and \textsc{Mach} I have attempted to characterize in \textit{Moderne Filosofer,}
p.\ 82--92.} Of more recent researchers must especially be named \textsc{Émile
Meyerson} in his works already several times discussed. When \textsc{Pierre
Boutroux}\footnote{\textit{L'idéal scientifique des Mathématiciens,} p.\ 152;
243.} declares that mathematics is a method, not an objective science, a similar
conception of the relation between formal and real science seems to lie at the
base.

With outstanding investigators of nature, even where they do not, like
\textsc{Maxwell} and \textsc{Mach}, go more directly into epistemological
questions, one will all the same often find utterances that testify to full
clarity concerning the relation between data and their rational interpretation.
The word analogy is perhaps not named, but the expressions used point decidedly
in the direction of the concept of analogy. Already in the title of \textsc{Niels
Bohr}'s treatise \textit{Atomernes Bygning og Stoffernes fysiske og kemiske
Egenskaber} (1922) this concept is really intimated. The physical and chemical
properties of substances are the given. The qualitative differences they present
--- the atomic weights of the elements, their spectra and so on --- can be ordered
in series, to which there answer, in the theory of the atom's structure, positions
and transitions of the electrons within the atom. What interests epistemology,
apart from the results significant for objective science, is the way in which
these results are expressed. There are always intimated in the expressions
analogies, not relations of identity --- a ``signifies'', not an ``is'', between
the quality-series which forms the starting-point and the quantity-series which
leads to a determinate thesis about atomic structure. Thus there is talk (p.\ 1)
of ``elucidation'', (p.\ 33) of ``explanation, or rather understanding'', (p.\ 36)
of ``interpretation'', and on p.\ 45 it is said: ``as the spectrum shows, and as
the atomic model makes intelligible''. The construction of the atom's inner
structure, and of what goes on within it, is guided step by step by observation of
changes in the qualitative properties of substances. --- It is clear that if the
elements did not allow themselves to be ordered in series where the members
followed one another according to a certain law --- if they thus formed a chaotic
difference-series and not a progressive difference-series --- no rational analogy
could be set up.

f. If it is right that the relation between formal and real sciences is a relation
of analogy, then there lies in this that advances in the domain of the real
sciences cannot be won by deduction from the results of the formal sciences. As a
method an analogy can only be a leading thought, an anticipation, not a proof. It
can lead to the putting of questions, and a well put question will in a scientific
mind lead on to closer investigations.

The fundamental concepts of our cognition themselves, the categories, lay
themselves open to view most clearly in the questions that are
put.\footnote{See on this \textit{Den menneskelige Tanke,} p.\ 137 f.} In the
leading thoughts which lead to questions and investigations there shows itself
above all the continuity in the history of science. A new circle of such leading
thoughts was formed at the founding of modern science nearly 350 years ago, and
these thoughts have preserved their force to this day --- especially after the
nineteenth century had worked with so great energy to put experiences in the place
of analogies. As already discussed in another connection, this strikes the eye
particularly in the study of the treatises in which \textsc{Robert Mayer, Joule,
Colding} and \textsc{Helmholtz}, at about the same time but independently of one
another, founded the modern doctrine of the subsistence of energy. They all
declared that they at bottom regarded it as a matter of course that there must be
equivalence between cause and effect --- which fortunately did not keep them from
investigating the relation through observation and experiment. But the matter of
course they assumed has not been acknowledged at all times. It is no older than
the seventeenth century, whose philosophers and investigators of nature maintained
that in the effect nothing new could be added and nothing lost of what was in the
cause. \textsc{Huyghens} and \textsc{Leibniz} maintained this proposition in
mechanics, after \textsc{Descartes} and \textsc{Spinoza} had maintained it from
what we should call an epistemological point of view. The proposition was, in the
manner of the time, set up purely dogmatically. \textsc{Spinoza} does indeed
ground it by saying that if there were anything other and more in the effect than
in the cause, this more would have no cause and would thus stand unexplained. But
who says that everything must be explicable? What was to be shown was precisely
that explanation was here possible, and the great thing about the discoverers of
the energy proposition is that they prepared the way for it.

But even if analogy has significance only as a foundation for questioning and
seeking, or as a heuristic principle, it is naturally of importance that it be
grasped in its determinateness and its limitation. We have seen that this was
still wanting in \textsc{Kant} when he thought he could set up a proposition about
``the subsistence of substance'' even before he had set up the causal proposition.
His discovery of genius was that between the concepts of ground and cause there
was only analogy, not absolute identity. But as regards the concept of ground
itself (or the relation between ground and consequent), he did not distinguish
between inferences that can be reversed (so that there is equivalence between
ground and consequent) and inferences that cannot be reversed. All real science
stands and falls with the causal proposition, according to which one
subject-matter stands to another subject-matter as consequent to ground; but that
ground and consequent (cause and effect) should be interchangeable is a more
special assumption which cannot be derived from the general causal proposition,
even though it may be said to formulate the ideal for all real science. There can
be science even if science's highest ideal is perhaps never reached. --- With
\textsc{Robert Mayer} a similar lack of clarity makes itself felt. A friend of his
youth, who stood in lively communication of thought with him before his
epoch-making treatise was composed (1841), says of him: ``Though he was far from
being a schoolman philosopher, a quite independent reflection in the domain of
logic and metaphysics on the essence of causality had perhaps quite as great a
share in his discovery as exact investigation of nature''.\footnote{\textsc{Gustav
Rümelin:} \textit{Erinnerungen an Robert Mayer} (Reden und Aufsätze. Neue Folge.
1881, p.\ 383).} But \textsc{Mayer} thought, as is seen from his foundational
treatise \textit{Bemerkungen über die Kräfte der unbelebten
Natur},\footnote{Reprinted in \textit{Die Mechanik der Wärme.} Herausg.\ von
Weyrauch, p.\ 29.} that his new proposition followed ``with necessity'' from the
proposition \textit{causa æqvat effectum}, and that there was therefore
equivalence between motion and heat. \textsc{Mayer} did not know that the
corrective to his proposition had already been given in a treatise by
\textsc{Carnot} (1824), which showed that the transition from heat to motion does
not go on as easily as the transition from motion to heat. And he did not see
clearly that in the transition from logic and mechanics to special physics there
can be no talk of deduction, but only of an analogical reflection of the abstract
ideals in the concrete relations between qualities. --- Analogy's methodological
significance naturally does not fall away because it must in many ways be
reshaped and limited by the very work that is done under its guidance. ---

g. Of analogies with heuristic significance within one and the same domain, those
are especially of epistemological interest in which one starts out from an
ordering within a certain totality and investigates whether a similar ordering is
not found again within other totalities. For Galilei, thus, the discovery of
Jupiter's moons contained an analogy for the elucidation of the sun's relation to
the planets (the earth included), and so for the strengthening of the Copernican
conception. And in the most recent time the analogy between the inner structure
of the solar system and of the atom has worked to make things intuitable and to
carry them forward. Such examples lead us over from the concept of causality to
the concept of totality. There lies in the nature of thinking a need to form
wholes. Every inference is a logical whole, and at the base of the causal need
there lies likewise a need to abolish the single subject-matter's isolation by
bringing it into as firm a connection as possible with other subject-matters. But
from this no deduction leads to special results. The application of the concept of
totality rests, like that of the concept of cause, on an analogy which may be of
great importance as a leading thought but has no probative power. As I have had
occasion to say in another treatise (Totalitet som Kategori, p.\ 51), the concept
of totality can never be used as an explanation, as happens in vitalistic,
spiritualistic and mystical-sociological theories. Like all categories, this
concept is a form for thought-work and at the same time an ideal which is sought
to be found again and carried through in the particular. Where observation
immediately presents subject-matters with a character of wholeness --- an atom, a
heavenly body, an organism, a personality, a society --- there we are set just so
many tasks: first for description and classification, then for the investigation
of their arising and of the conditions of their
subsistence.\footnote{Cf.\ more closely \textit{Totalitet som Kategori,} p.\
52--66.} It is an analytic work on the foundation of what observation can present
that is here required. As \textsc{Leibniz} has said, it is the law for the
connection of the elements that constitutes a being's peculiar wholeness, and it
is such a law that the analytic work is to seek to find.

As it stands with the concept of totality, so it stands also with the concept of
development. This concept too is a form of thought which guides research both
involuntarily and with full consciousness. Its methodological significance shows
itself in this, that it contains an invitation to attempt to construct a series in
which each member conditions the following and is conditioned by the preceding,
and that in such a way that in later members a greater manifold of elements is
always connected in a more intimate way. It becomes a regularly varying
difference-series which cannot be reversed. (Cf.\ Den menneskelige Tanke, p.\ 233
f.) As regards the development of the single organism or the single organ, the
doctrine of epigenesis founded by \textsc{Caspar Friedrich Wolff} maintains that
there can be only analogy between the different stages, whereas earlier theories
found identity here, regarding the later stages as completely preformed in the
earlier. The same thought as in \textsc{Wolff} appeared at about the same time in
\textsc{Rousseau}, in his psychological-pedagogical theory, which maintained that
the difference between the child and the adult was not merely a difference in the
dimensions. In the nineteenth century the concept of development was also
attempted to be carried through for the organic species, and men sought to exhibit
a continuity in the processes of development, whereby the analogy between the
different stages acquired a real foundation. \textsc{Cuvier} was still content
that the analogy had its ground in the Creator's imagination, which after every
destruction of organic life caused by a geological catastrophe brought forth new,
analogous forms. In \textsc{Goethe} and the Romantic investigators of nature,
metamorphosis and development were only figurative expressions for the analogy of
the forms; they assumed displacements in nature's systematic imagination, but not
any real descent. Many of \textsc{Darwin}'s opponents, especially anatomists and
morphologists, stood at such a standpoint. When modern strict critics of Darwinism
sometimes declare that they do not understand the standpoint those morphologists
took, it is a sign that the analogy between the stages of development has to such
a degree shown itself in its real significance that researchers of the present
cannot enter into their own subject's past. Nothing better shows the importance of
the study of the history of science with special regard to the fundamental forms
within which research moves. Those who are led by certain thoughts do not always
discover the leading-string, and the consequence of that can be not only a lack of
historical understanding, but also a dogmatic confidence in the forms of their own
time.

Within biology a sharp distinction is drawn, especially since \textsc{Owen}'s
time, between analogy and homology,\footnote{Cf.\ \textsc{Yves Delage:}
\textit{Comparative Anatomy and the foundation of Morphology.} (Congress of Arts
and Sciences. 1904).} in such a way that the former rests on an agreement in
function and is thus of a physiological kind, while the latter rests on agreement
in structure and is thus of an anatomical kind. The lungs of mammals and the
swim-bladders of fishes are homologous, not analogous; the lungs of mammals and
the gills of fishes, on the other hand, are analogous, not homologous. There can
be a dispute between anatomists and physiologists whether analogy or homology is
of the greater importance. In part the famous dispute between \textsc{Cuvier} and
\textsc{St.\ Hilaire} was of this kind. \textsc{Cuvier} thought essentially of
unity in structure (unité de composition), while \textsc{St.\ Hilaire} understood
by unity analogy.\footnote{\textsc{Cuvier:} \textit{Histoire des progrès des
sciences naturelles.} Bruxelles. 1838. II, p.\ 455.} Yet it must be remarked that
epistemologically likeness of form, likeness in structure, is often to be called
analogy, since it is more the relation between the parts than qualitative likeness
at single points that is in question. According to the doctrine of epigenesis
there is, as already remarked, only analogy between different stages of
development of the same organ, even if the development is seen to have gone
continuously from the one form over to the other; and according to
\textsc{Cuvier}'s own account, \textsc{St.\ Hilaire} thought essentially of
analogies of structure in different species. Seen from this side it was thus a
dispute about likeness of form or likeness of relation. As I have remarked in my
Psykologi, likeness of form is a transitional form between qualitative likeness
and likeness of relation.

h. As regards the understanding of organic life, the possibility must, as already
touched on, be excluded that the concept of organism should itself contain an
explanation. Concepts such as life and organism stand themselves in need of
explanation. Since now the history of physiology shows that all scientific
explanations in these domains have been won through the leading back of organic
phenomena to physical and chemical relations, it might seem to be warranted to
apprehend the organism all through in analogy with a mechanism. But there always
remains here the difficulty that with a mechanical whole we can know beforehand
the elements by whose joining together it comes into being, just as we can see
bricks lying heaped up before they are joined into a building. So external a
relation does not obtain between the organism and its elements. One can find
again, by analysis of the organism and its processes, elements known beforehand,
but not explain these processes as the mere product of such elements. As a
biological researcher has aptly said: in biology the whole comes forward as given,
and we go back from the form of the whole to the elements it shows itself to
consist of; but in physics and chemistry we have conversely the elements given,
and the task is to understand their combination into the whole --- the atom is not
given, but is constructed from the material processes that can be
exhibited.\footnote{\textsc{Johan Hjort:} \textit{The Unity of Science,} p.\ 166.}
--- Against the analogy with mechanism there has often been set the analogy with
human production, now in a more external way, by using the comparison with a
master-builder's relation to the building he erects, now more spiritedly, by using
the comparison with an architect of genius and the plan that comes into being in
his creative imagination. But in both cases one then distinguishes between creator
and work, and here the analogy fails. --- \textsc{Kant}'s result in his
investigation of this problem stands fast to this very day, when he declares:
``Genau zu reden, hat die Organisation der Natur nichts Analogisches mit irgend
einer Causalität, die wir kennen'' (Kritik der Urtheilskraft § 65). This does not
now exclude that analogy's spurring effect can make itself felt here too in
research. Both the analogies adduced can, namely, be applied heuristically, and
are in fact applied by different researchers, in that now the question presents
itself what physical and chemical conditions an organic process presupposes, now
it is asked how far an organism is so equipped that it can maintain its
subsistence under determinate conditions, especially when these are subject to
change. The future must show whether this position of biology's is the final one.

i. In the question of analogy's significance in the scientific sense there makes
itself felt an interesting opposition between two excellent French epistemologists.
\textsc{Émile Meyerson} regards, as we have earlier seen, analogy as the last form
of scientific thinking about the subject-matters before us. On his conception,
man's rational need finds, by the carrying through of scientifically analysed
analogies, the satisfaction that is possible at all. But in the irreducible
qualities there always remain irrational
elements.\footnote{\textit{De l'Explication dans les Sciences.} I, p.\ 217--224;
II, p.\ 291--295 and \textit{Bulletin de la Société française de Philosophie.} 24.\
Febr.\ 1921, p.\ 59.} Against this conception \textsc{Léon Brunschvicg} objected
that it lays a false ideal of cognition at the base. If true science were to be the
carrying through of absolute identity as cognition's highest form, the attainment
of this ideal would be one with intellectual suicide. As soon as one gives up this
ideal, the supposed irrationality falls away, since it merely marks the opposite of
absolute identity. True science seeks points of view which make it possible to find
equivalence between complex subject-matters and their elements; but it cannot be
set up as its ideal to falsify the qualities before us by transforming them into
pure quantities.\footnote{\textit{Bulletin.} 1921, p.\ 59 f.
\textit{L'Expérience humaine,} p.\ 581.}

For my own part I have already, in a lecture I gave at Oxford in 1904, printed in
\textit{Mind} 1905 and in my \textit{Mindre Arbejder} (second series), professed
the point of view of analogy, and in \textit{Den menneskelige Tanke} (p.\ 181--201
and 309--318) I have taken up this consideration again. The clash between identity
and analogy runs through the whole history of philosophy. It has cost, and it
costs, resignation to see an analogy carried through as far as possible as the
highest that can be reached, and as the only way in which the different roads of
cognition can meet, and in which the relation between the different domains of
cognition can be determined. But since analogy is only likeness of relation, not
identity, there always remains an irrational element --- the more so as it cannot
be known beforehand whether the analogies now at our disposal can be carried
through in the case of new subject-matters, or whether it is possible to find new
analogies when unforeseen subject-matters demand them. Beyond exact analogies it
would be possible to get only if it could be shown that the inferences, in analogy
with which we apprehend the qualitative and temporal differences, could be
reversed --- so that one could be sure that no other premisses than those on which
our sciences build could lead to results agreeing with the observations.

\subsection*{19. Natural Science and the Human Sciences}

In the elucidation of the relation between formal and real science, natural
science has in the foregoing been used as the example of real science, because it
is incontestably the most perfect real science. The human sciences (philosophy,
the science of language, historical research) work constantly with qualities and
successions which it has not been possible to lead back to relations as clear as
those which natural science has for an essential part succeeded, by the way of
exact analogies, in exhibiting for its own subject-matters. But then it is a
matter of investigating the relation between natural science and the human
sciences, especially with regard to whether this relation is a relation of
identity or a relation of analogy.

One might ask whether there is at bottom any science other than the human
sciences. What we call natural science is after all the fruit of spiritual work,
and its methods, as well as the results reached by the help of those methods,
must be determined by the essence of spiritual work. Can natural science ever
become anything other and more than a great attempt to apprehend existence as
spiritual work is able to? Natural science is to that extent only a peculiar form
of human science, and the difference is only that the human sciences, besides
investigating the spiritual work which leads to science, also investigate other
forms of spiritual work --- for instance those which lead to art, morality and
religion. Of special interest in this connection, however, is that epistemology
itself is undoubtedly a human science, whose subject-matter is that capacity for
work, that energy, which lays itself open to view in the fact that our
representations undergo changes leading at once to greater connection among them
and to more exact determination of each of them. The epistemological way of
regarding things gives us a unifying point of view for all science --- perhaps
the only one there is. Far from covering over the differences between the
sciences, such a unifying point of view precisely sharpens the eye for those
differences. The different ways of working, and the different success with which
work is done in the different fields of work, gain precisely in interest from
such a point of view. And it now also becomes more easily possible to investigate
the relation between scientific work and spiritual work performed in other
domains of human mental life than the intellectual. Here too the unifying point
of view will let the peculiar differences come forward with greater clarity.

But even if a con\-sid\-er\-a\-tion such as the one now indicated should be the
fun\-da\-men\-tal and the de\-fin\-i\-tive one, the matter presents itself in many people's
opinion otherwise when we turn our attention to the different fields of
intellectual work. When it now shows itself that the work of natural science ---
whose field is the so-called material phenomena, among which is reckoned
everything that comes forward for us in the form of space --- proceeds by more
perfect methods and is attended by greater success than the work of the human
sciences, which does not get beyond differences of quality and of time, would it
not then be reasonable to make all intellectual work into work of natural science
--- to let the more perfect knowledge elucidate the less perfect knowledge and
take it up into itself, by apprehending all so-called spiritual subject-matters
as members in the great continuous connection between all so-called material
subject-matters which natural science more and more discloses to us? Everyone
concedes (however he may otherwise apprehend the relation between ``soul and
body'') that spiritual subject-matters, whatever constitution they may have, hang
together with certain physiological processes; and since these last decidedly
fall within natural science's domain, would it not then be most correct to take
one's starting-point here and regard the natural-scientific apprehension as the
only valid one, and then apprehend the human sciences (in so far as they can be
called science at all) as a part of natural science?

When I attempt to answer this question, I must first refer to what I have
attempted to show in earlier sections of this treatise: that all real science
builds on analogy with logical and mathematical thinking, so that the question
finally becomes one of that thinking's validity. Are inferences and equations
spiritual or material? If they must be said to be forms of, and results of,
intellectual work, it must be asked who draws the inferences and grounds the
equations. We thereby stand before the relation between epistemology and
psychology, and must dwell a little on it before we go over to the relation
between psychology and physiology, which more immediately concerns our problem.

a. As regards the first relation, I must refer to my treatise \textit{Totalitet
som Kategori} (p.\ 4--18; 46--48), where I have sought to show that the
circumstance that psychology's possibility is one of epistemology's
subject-matters does not abolish psychology's independence, any more than
physics' independence is abolished by its also being such a subject-matter. The
epistemological subject which is presupposed in all inferences and equations is
now one of psychology's most interesting subject-matters. By the help of the
history of science, psychology seeks to show how this subject has developed from
being merely observing and describing to so high a stage that it has been possible
to say that it is this subject that prescribes to nature its laws. It is in
constant development; at each time it is at a determinate stage of that
development. It does not exist in and for itself; that one can assume only by a
leap from epistemology to metaphysics. And at every one of its stages of
development it sets psychology the task of exhibiting how its development has been
possible at all --- and that for the good reason that at every stage it presents
specifically psychological, and not merely purely logical and mathematical,
properties. It is the principles of spiritual growth that psychology here applies
to the subject of cognition which all science presupposes. There are
epistemologists who would most gladly slip everything psychological off
themselves; they are hit by the irony with which \textsc{Kierkegaard} speaks of
the speculative philosophers of his time. The Marburg school, so meritorious in
other respects, is hit by such irony. Naturally psychology cannot build on other
principles than those acknowledged by epistemology; but that it has in common
with --- epistemology itself, which naturally cannot build on other principles
than such as it can acknowledge when it one day gets finished. (Cf.\ above, § 7.)

b. As regards the relation between psychology and physiology, I can also refer to
earlier utterances --- partly in my \textit{Psykologi} and my \textit{Psykologiske
Undersøgelser}, partly and especially to my university programme
(\textit{Filosofiske Problemer}) (1902), whose content is taken up in abbreviated
form in \textit{Den menneskelige Tanke} (1910) (p.\ 22--34; 329--338). I shall in
this connection dwell only on a few points. An observing (respectively
experimenting) and describing psychology is the presupposition for finding the
physiological connection which in some people's opinion\footnote{In this
direction goes, within our own literature, the acute treatise of \textsc{Jørgen
Fr.\ Jørgensen: }\textit{Problemet om Forholdet mellem Sjæl og Legeme.}
(The journal ``Vor Tid''. 1913.)} is to give us the only scientific explanation of
mental life. In order to ``explain'' something, I must first know what it is that
is to be explained; otherwise I cannot know whether the right explanation has been
found. So-called psychophysiology therefore always presupposes psychology proper.
The programme of leading all psychology back to physiology can be carried through
only when empirical-experimental psychology has first done its work and thereby
made precise the phenomena that are to be explained. Psychology has great
difficulties to carry through in its own domain, and it is of no use to proceed
precipitately. The discontinuity of mental phenomena, their qualitative
differences and their constant changing, as well as their manifold and
complicated interweavings, constantly set new tasks for descriptive, genetic and
analysing psychology.

If we now look at the physiological ``explanations'' which are to impress the
psychologist by their exactness, it shows itself that they have been won by the
help of analogies drawn from direct psychological observations. They are attempts
to construct a physiological mechanics answering more or less to what
psychological ob\-ser\-va\-tion shows. It is a psychological fact that the more
fre\-quent\-ly two rep\-re\-sen\-ta\-tions accompany one another, the more easily the
transition later takes place from the one to the other. This is to find its
explanation in the ``path-forming'' of nerve-cells and nerve-fibres. But in
reality one knows no more of what such nerve path-forming properly is than what
lies in the fact (which is both a psychological and a physiological fact) that a
function goes on more easily by repetition. One has got a mechanical image; that
is the whole of it. And such an image might very well have been applied directly
in psychology itself, since all psychological expressions show themselves
etymologically to be fetched from the outer world.

One of the first and most remarkable attempts to apprehend psychology as
physiological mechanics was made by \textsc{Thomas Hobbes}, who will maintain that
sensation and thinking are motion in space and nothing else. But he then stops
involuntarily at the fact that motion can be apprehended (become phenomenon), and
that is not explained by the motion itself. (See the appendix to this section.) In
\textsc{Richard Avenarius}' acute and instructive work \textit{Kritik der reinen
Erfahrung} the attempt is made to carry through the conception that a scientific
apprehension of experience and thinking is got only by regarding these functions
as processes in the central nervous system. A whole series of such processes is
then depicted; but when it is asked what they signify, it shows itself that they
are built on analogy with directly observed mental phenomena. The so-called
``objective vital series'', which is to be regarded as the independent variable in
relation to the ``subjective vital series'', is in reality dependent on the
latter. (Cf.\ my characterization of \textsc{Avenarius} in \textit{Moderne
Filosofer}.) While \textsc{Hobbes} does in flashes come to intimate the subjective
presuppositions of physiological psychology, and while \textsc{Avenarius}
expressly acknowledges that there are two ``vital series'' standing in a functional
relation to one another, \textsc{Bechterew} in his \textit{Objektive Psychologie}
(1913) has attempted to write a psychology avoiding all psychological
designations. But the consequence is that the reader must constantly guess why
just these determinate physiological processes are described in this determinate
connection. Thus great weight is laid on ``concentration-reflexes'', in opposition
to sporadically isolated reflexes, and the reader then guesses that it is
attention that is being thought of. The expression ``concentration'' can very well
be used in a purely psychological sense, of our gathering of the mind in a
determinate direction, just as we saw before that one speaks of path-forming as a
good image for a practised course of thought. The place given to the
``concentration-reflexes'' is due to an analogy with mental concentrations, such as
descriptive psychology can present them. \textsc{Bechterew} himself concedes that
to the difference which observing psychology draws between the unconscious and the
conscious, the involuntary and the voluntary, he can find nothing physiologically
corresponding; and therefore, he says, this difference has no significance for
objective psychology. If objective psychology must give up before differences of
the kind stated, then it has not much significance. I think now that
\textsc{Bechterew} gives up too early. Even if one does not think that
physiological psychology is the only scientific one, it can and must all the same
be set as an inescapable task (whether it can be solved or not) to find
physiological correlates to those differences so significant in the psychological
respect.

Physiological processes are directly measurable in a degree in which psychical
processes are not. When we set the two kinds of process together, we stand again
before the relation between a quantity-series and a quality-series. Here too
identity and analogy must not be confused. The members of one of the two series
can be used as symptoms of what goes on in the other. If one cannot connect any
meaning with a thought's being the product of the reciprocal action of material
elements, or with a thought's bringing forth an alteration in the body, then one
must stop at a relation of analogy\footnote{For the rest, the popular theory of
reciprocal action would also build on analogy, since it presupposes that the
relation between soul and body is analogous to the relation between two bodies
(two billiard balls, e.g., or --- when it is to be more ``scientific'' --- two
groups of atoms). \textsc{Descartes} lets body and soul ``collide'' with one
another, and \textsc{Lotze} has nothing against this theory.} which may perhaps
one day become a relation of proportion. The conception we come to by keeping to
what can be observed is that there are in any case some beings which have both
psychical and material properties answering to one another, without the one group
of properties being derivable from the other. It is a riddle that here remains ---
an example of the fact that where we must stop at an analogy, we stop at a
question. Such a relation may for that matter occur in purely psychological form.
There may be in a man's character properties that cannot be derived from one
another, and which we therefore do not understand, because we do not know the
character itself deeply enough to see how the two different properties hang
together in their inmost ground. As regards the relation between the spiritual and
the material, we are so placed that we can now infer from mental states to states
in the organism, now conversely. \textsc{Spinoza}, the founder of the conception
to which I here most nearly attach myself, does indeed already go both
ways.\footnote{\textit{Spinoza's Ethica.} Analyse og Karakteristik, p.\ 48 f.;
87 f.}

It is finally the relation between the concepts of personality and organism that
we here come to stand before. What \textsc{Kant} said of the organism, that at
bottom no analogy to it can be found, holds for the personality too. So intimate a
connection as that between the elements of mental life we know no parallel to. The
more interesting is it that those two concepts are mutually analogues.

\textit{Appendix to § 19.} \textsc{Hobbes} apprehended everything, sensation and
thinking too, as motion in space and \textit{nothing else}. But at one place
especially it dawns on him that the matter is after all not so simple. It is
finally through sensation that phenomena --- that is, motions --- come into being
for us, and that science becomes possible. And when it is then asked after the
causes of sensation itself, we cannot take our starting-point from any other
phenomenon than precisely that one (sensation) itself (ad causarum sensionis
investigationem ab alio phænomeno \textit{præter eam ipsam} initium sumi non
posse. De Corpore XXV, 1). It is on this account that \textsc{Hobbes} declares
that of all phenomena the most remarkable is the very fact that anything can
become a phenomenon (id ipsum τὸ φαίνεσθαι), \textit{admirabilissimum}.

I understand this remarkable passage in the following way. That anything can
become a phenomenon \textsc{Hobbes} finds so remarkable because a motion cannot
become a phenomenon for another motion (that is, be \textit{known} by it). The
one motion can continue the other, hamper it or stop it. Sensation must then be
different from \textit{mere} motion, though it is \textit{also} motion. Now since
it is through sensation that we know everything, it must also be through sensation
itself that we know sensation. But to sense one's own sensation is to remember;
and there is accordingly, \textsc{Hobbes} maintains, memory and comparison bound
up with all sensation. (Sentire se sensisse, memoria est. --- Sensioni necessario
adhæret memoria aliqva, qva priora cum posterioribus comparari et alterum ab
altero distingui potest. De Corpore XXV, 1; 5.) It is in consequence of this
conception of sensation that he sets up his famous proposition that in all
sensation there is contained a manifold (phantasmatum varietas), and that always
to sense the same and not to sense at all come to the same thing (sentire semper
idem et non sentire ad idem recidunt) (5). From this purely psychological
observation he then infers that the motion to which sensation is bound must be
capable of being held fast and recalled, and that there must be conditions for
this in the sense-organ. More clearly the difference between the subject of
cognition and the motion which according to \textsc{Hobbes} is the only thing
objectively existing cannot be expressed. But \textsc{Hobbes} was too much of a
dogmatist for this flash of thought to make him give up his metaphysics of motion.
It went with him as it later went with \textsc{Schelling}, who (as I have shown in
\textit{Den nyere Filosofis Historie}\textsuperscript{3} III, p.\ 158) had an
analogous flash of thought, in setting up very clearly the problem: how can there
exist not only nature, but also a science of nature? This flash hindered
\textsc{Schelling} as little from going further in his Romantic philosophy of
nature as \textsc{Hobbes}' flash hindered him from going further in his
mechanistic philosophy of nature.

In \textsc{Frithiof Brandt}'s excellent monograph \textit{Den mekaniske
Naturopfattelse hos Thomas Hobbes} (1921), in which light is thrown on so many
hitherto overlooked or misunderstood points in \textsc{Hobbes}' philosophy, the
passage adduced above is discussed at length (p.\ 365--373). The author seeks to
show that there is in reality no inconsistency here in \textsc{Hobbes}. But no
account is taken of the proposition which in my eyes is decisive, that the
understanding of sensation must be won from sensation itself, nor of the
consequent circumstance that from the purely psychological understanding of
sensation's nature one must \textit{infer} to the constitution of the sense-organ
and to what must go on in it.

My interpretation of the utterance of \textsc{Hobbes} in question as a flash of
thought leading beyond his absolute metaphysics of motion is confirmed by his
statement (De Corpore VI, 7) that one can come to ethics and political science
without going the road of natural science, but merely by building on the
psychological experiences of need and impulse which everyone can make directly in
himself (per uniuscujusqve proprium animum examinantis experientiam). Not only
need and impulse, but sense-sensation too, we know directly by this road.

\subsection*{20. Ethics and Theoretical Science}

When ethics is put in a certain opposition to theoretical science, this does not
rest on one's having the right within ethics to think and infer in another way
than within all other science. There is no special ethical thinking as regards
grounding and the conducting of proof. But ethics has presuppositions which other
sciences need not set up. Every thinking, however theoretical, rests on the
presupposition that truth has value, that it is valuable to seek after it. To such
a degree is this presupposition necessary that one does not even consider it
necessary to name it. When in the present a philosophical tendency, whose most
considerable representative is \textsc{Heinrich Rickert}, has all the same laid
weight on this presupposition,\footnote{\textit{Den nyere Filosofis
Historie}\textsuperscript{3} IV, p.\ 219.} the warrant lies in this, that the
problem of value is thereby drawn out of the isolation which can easily follow
when it is seen only in its significance for aesthetics, ethics (economics
included) and the philosophy of religion. When such isolation is averted, there
comes forward a fundamental connection in human spiritual life, a connection which
lies deeper than the difference between the different problems, but which at the
same time grounds the difference that there is. When science is preferred to
artistic or religious striving, then it is value that stands against value, and a
clash takes place between different spiritual needs. The problems come into being
by each of these spiritual needs setting thoughts in motion. It then depends on
whether these thoughts can be brought into agreement with one another. But here it
is that the great significance of the value of truth, all science's tacit
presupposition, comes forward, precisely in relation to all other values. The
thoughts which the different values set in motion must --- for the values' own
sake --- be subjected to a purely theoretical testing, a testing whose sole purpose
is to find the thoughts' grounding and their mutual relation: a testing, that is,
from the value of truth as the ruling one. Of all the problems, therefore, the
problem of cognition takes precedence. On the result to which the discussion of
the problem of cognition leads depends the validity that can be ascribed to
thoughts set in motion by other motives than the purely intellectual. In the
treatment of the problem of cognition it is presupposed that we have no other
interest than that of finding the truth, however this may be constituted. It then
becomes the question how other interests that stir in us can be united with this
interest. In my book \textit{Den menneskelige Tanke} (see especially p.\ 22) I
have therefore proceeded by presupposing that an intellectual interest has
developed, a joy in thought-work itself, which leads to thought's being able to
develop according to its own laws and according to the demands of the
subject-matters. The other interests, for which thought-work is merely a means and
not an end in and for itself, will finally find themselves best served by the
intellectual interest's ruling for a time, and by a theoretical revision being
undertaken from time to time of the thoughts they have especially called to life.

In the ethical domain it cannot now be avoided that interests other than the
purely intellectual make themselves felt. All ethical thinking presupposes a
fundamental value and aims at finding means and ways to maintain it in the
judgments that are passed and the social orderings that are formed. In an apt way
\textsc{Henri Poincaré} has in his treatise \textit{La Morale et la
Science}\footnote{\textit{Dernières Pensées.} 1913, p.\ 229.} expressed this by
saying that ethical thinking has a premise which is imperative, since it springs
from a need demanding satisfaction, while the other premisses are indicative,
since they state the means and ways by which the satisfaction can under the given
conditions be reached. Since there is an imperative premise, the conclusion must
necessarily also be imperative. The consistency that can here be in question
becomes a consistency of the relation between means and ends. If the grounding of
the need or value which gives rise to the imperative is asked for, the answer
must --- since value can be measured only by value --- be fetched from the
fundamental value which the special value presupposes. And if the grounding of
this fundamental value is asked for in turn, we must keep to the fact that every
determination of value presupposes a whole. Values do not float in the air, but
are the expression of a whole's need for subsistence. It may be a single moment,
a period of life, a single interest, that one makes into one's all in all; it may
be one's own personality, or the class and stock it belongs to, or finally the
subsistence and development of the whole human race, that constitutes the last
presupposition of imperative ethics. The dispute of theories rests finally on the
fact that different wholes can be laid at the base, and that thereby the
fundamental values too become different.\footnote{The connection between the
concepts of totality and value I have exhibited especially in \textit{Den
menneskelige Tanke} 145--146; 151, and in \textit{Totalitet som Kategori,}
Kap.\ V.} An attempt at scientific ethics can here proceed only by taking the
starting-point from a determinate whole, drawing all inferences from the
conditions of its subsistence and development, and thereafter seeking to show how
the other wholes are to be valued in accordance with the results so won. In this
last investigation there is special ground for analogy. Even if the other wholes,
in virtue of the starting-point, must be treated as mere means, such treatment
becomes purposeful only when there is an understanding of their nature. And if
they are more than means --- if a relation of solidarity or even of sympathy
follows from the given starting-point --- a still more deeply going understanding
of them is required; fortunately in such cases the capacity for understanding can
grow with the need. Between peculiar wholes there is a relation of analogy, not a
relation of identity. Ethical thinking's validity (from the given starting-point)
rests on the strictness with which the analogies are carried through.

One might think that the most comprehensive whole must in and for itself furnish
the right starting-point for ethical valuation, so that its subsistence became the
self-evident fundamental value. But besides extent, inwardness and strength also
come into consideration with the concept of value. Within a less comprehensive
whole there may be an inwardness and a strength in the union which do not so
easily arise within a greater whole; and the greater whole's value will for an
essential part rest on the independence it is able to acknowledge and to further
in the smaller wholes it comprehends.

By the analogies between the different wholes it becomes intelligible that there
are certain ethical concepts which are formed, and must be formed, at all
standpoints, because they express the formal relation between a whole, regarded
as fundamental value, and the conditions of its subsistence and development. Such
formal ethical concepts are right and duty. The single whole is to maintain itself
against other wholes and against outer conditions generally; that is its right,
which follows from its position as fundamental value. But its task is besides to
guard the means to its subsistence, and also to work for the subsistence of those
wholes which are means for it, or to which it stands in a relation of solidarity
or of sympathy. Under the concept of duty belongs prudence, the understanding of
the nature of the means and of the necessity of procuring them. Especially with
regard to the significance and application of these formal concepts, the one
ethical standpoint will be able to learn from the other. It is impressed on us, as
regards the duty of prudence, by the parable of the unjust steward (Luke XVI),
which shows that the children of darkness often have at their disposal, within
their own sphere, a prudence which is wanting in the children of light in their
work. It has been said in our own day by an experienced statesman that the task of
clever men often consists in remedying the misfortunes good men have caused. ---
\textsc{Leibniz} has defined the concept of justice in such a way that right and
duty, humanity and prudence are united in it: \textit{caritas sapientis}!

Not only between different ethical standpoints among themselves can analogy make
itself felt in a fruitful way, but also between ethical consciousness and
scientific consciousness. The scientific consciousness, which seeks clarity and
understanding for their own sake, can become a model (a Paradigm, to use
Aristotle's expression) for ethical consciousness, both as regards disposition and
as regards consistency. The same can be said of the artistic consciousness in its
striving to achieve the perfect. Nothing, says \textsc{Michel Angelo}, makes the
soul so pure and noble as striving to accomplish a perfect work. In all spiritual
domains it is a matter of uniting unity and closedness with the free unfolding of
a fullness of force. Justice (in the sense of the word adduced) may be called the
ethical synthesis, as truth is the intellectual synthesis. (Den menneskelige
Tanke, p.\ 365--367.) In the treatise on morality and science adduced above,
\textsc{Henri Poincaré} says: ``Intellectual habits have a moral resonance. When
one accustoms oneself to disregard particulars and accidents, because they are of
no importance for the understanding, there will arise a natural inclination to
subordinate particular interests to general ones altogether''.

\subsection*{21. World-view and Science}

It is a true word of \textsc{Wundt}'s that it is easier to reject all metaphysics
than to act in accordance with that rejection. By many hidden ways combinations
and constructions steal into the background of consciousness, and now and then
they even press forward into the foreground. Again and again a critical
investigation of such constructions will be necessary. And first and foremost it
is of importance that one become clear about the constitution of the problem one
here stands before. As I have attempted to show in another place (Den
menneskelige Tanke, p.\ 306 f.), there is an analogy between the problem of
existence and the problem of cognition. Cognition becomes, as regards its
validity, a problem through the fact that human thought is indeed a fact, an
actually given subject-matter, but that it is not the only subject-matter and yet
is to hold both itself and all other subject-matters. It is to develop according
to its own laws and yet to build a world in which all subject-matters, itself
included, can find a place answering to their peculiarity. Thought is to rule and
is yet an element in a world greater than itself. The world is great and thought
is small, however energetically it stirs. The problem of existence --- that is,
the question of the possibility of a scientific world-view --- is likewise due to
a relation between part and whole. But here the question is whether any single
subject-matter, any part of existence, can stand as the type for the whole of
existence, as that primal phenomenon which gives the characterization of the whole
of existence in brief. All historically given world-views, religious or
speculative, are crystallizations, involuntary or constructed with conscious
reflection, about a primal phenomenon. In the series of these crystallizations one
can read off which subject-matters have at different times drawn interest to
themselves and become unconscious or conscious central points in men's minds. ---
It is clear that in the coming into being of a world-view the thinking is by
analogy, ἐν παραδείγματι as Aristotle says. One subject-matter is used as a model
for all subject-matters.

We have here operated without more ado with the concepts of part and whole. But
have we the right to start out from the assumption that existence constitutes a
whole? --- \textsc{Bertrand Russell} has declared that the representation of a
universe as a whole is a remnant of pre-Copernican astronomy which not even
\textsc{Kant}'s criticism has been able to get rid of.\footnote{\textit{Mysticism
and Logic,} p.\ 987.} It is doubtless also so that it is the sensible picture of
the world that determines the representations in this domain, even though science
has long since given up its absolute validity. But there is all the same a deeper
foundation for the tendency that here makes itself felt. The fundamental law of
all categories, as of the life of consciousness generally, is the law of
synthesis. There works here, to use \textsc{Kant}'s expression, a hidden art in
man's interior which weaves sense-sensations into a sense-image, sense-images into
concepts and laws, and finally, by the help of the lawlike connection, a great
image of the whole, more or less carried through and elaborated. Criticism always
comes afterwards. And the most important critical objection against an attempt to
fix such images of the whole is that it would be the death of thought. Since all
thinking is finally a gathering-together, there is presupposed in all thought-work
a manifold, a discontinuity, which is to be overcome. And with a complete and
concluded satisfaction of the need for wholeness, that very ``hidden art'' would
die a straw death. It would no longer do what in its finest and noblest forms it
has been able to do: shape great ideals, set goals pointing beyond what has
hitherto been won. Discontinuity must have a place among the fundamental
categories, to remind us that there is always work to be done. Every whole that is
found can mark only a provisional station. Our cognition works constantly on a
scale whose highest step would be the unattainable absolute identity-series, and
whose lowest step would be the undemonstrable chaotic difference-series.

In the most pronounced cosmological systems --- \textsc{Parmenides', Spinoza's,
Hegel's} --- identity has been the last and highest point of view. The ideal was
to apprehend everything as in the last instance one. It was at bottom an oriental
ideal that here came to rule in European philosophizing. A turning-point came with
\textsc{Leibniz}, for whom identity marks a methodological principle in our
grounding and proving course of thought and can acquire cosmological significance
only by the way of analogy. \textsc{Leibniz} maintains precisely, with great
emphasis, that all cosmology (metaphysics), and therewith all theology too, builds
on analogy --- finally on analogy with the human spirit and the harmony and
wholeness that can be found in it, or that it at any rate strives after. Still
more searchingly \textsc{Kant} sought to draw the boundary between the science
that works with the principle of identity and the faith that works with analogies.
At the beginning of the nineteenth century new attempts were made at a cosmological
doctrine of identity. But it seems as though the point of view of analogy is
winning more and more adherence, even among thinkers of a predominantly
speculative tendency.\footnote{Cf.\ e.g.\ \textsc{Bosanquet:} \textit{The meeting
of extremes in contemporary philosophy} (1921), p.\ 17.} The figurative character
of the representations which are formed, involuntarily or with full consciousness,
at thought's boundary and beyond it, is more and more making its way. This
figurative character indicates that a choice has been made between the mutually
analogous series --- the formal (logical-quantitative) and the real
(qualitative-successive) series, and within the real series the natural-scientific
and the human-scientific. But the analogy between these series is and remains
science's last result. A scientific conclusion would require the identity of the
series, not mere analogy. It simply is so that our exact concept of reality comes
into being through the formal series' making the real series intelligible, and
that the formal series are themselves a side of reality --- the side we become
aware of as soon as we try to orient ourselves in existence.

How one takes one's stand at the very boundary of scientific thinking will depend
in part on which epistemology one adheres to. There is here a somewhat different
situation for pragmatism and for criticism. Both conceptions maintain that the
thoughts which can be shaped when we have reached that boundary can only be
analogies. But pragmatism, which on the whole is inclined to ascribe validity to
all thoughts springing from human need, and which is not interested in a critical
investigation of our fundamental concepts, will also have no interest in
investigating how far one can still trace, in the limiting concepts which are
human cognition's last word, the fundamental forms characteristic of it. Criticism,
on the other hand, will be interested in finding again the structure of the human
spirit beyond the domain in which scientific concepts can be formed; in the very
kind of the figurative representations that are still possible beyond that
domain's boundaries, it will seek for peculiarities of human spiritual life. It
lays weight on thought's homologies, on the kinship which a common structure can
condition, not merely on the analogies that arise from thought's having to be
adapted to new conditions.

Cosmology cannot become science. It stands between science and art. It forms the
transition to poetry and religion, the domains in which the need to find concrete
expressions for the experiences of life shows itself freely and involuntarily.
Even there philosophy will seek not only analogy but also homology, since the
nature of the human spirit will lay itself open to view there too.


% ============================================================
%  V. POETISK OG RELIGIØS SYMBOLIK -- pp. 129--137 -- § 22 -- notes 78--79
% ============================================================

\chapter{Poetic and Religious Symbolism}
\label{ch:5}

\hnum{22} Out of that merging into involuntary analogies which is peculiar to
human spiritual life at the most primitive stage we know, there has developed,
through many transitional forms, the scientific consciousness which demands
identity as the strict presupposition of every transition of thought, and which,
where this presupposition cannot be exhibited, works with rational analogies.
Science's development is a member in the great spiritual division of labour that
comes with advancing culture. But when one element is released from the primitive
connection in which all spiritual forces are joined undifferentiated, other
elements are released too, or at any rate get other conditions to work under.
Alongside, and more or less under the influence of, the independence of
intellectual culture, poetry and religion too have in greater or lesser degree
released themselves from the original boundness. In opposition to the rational
analogies there then gradually come forward the emotional analogies characteristic
of poetry and religion.

The immediate absorption, or rather enthralment, by a subject-matter can give rise
to an equally immediate need to shape it in an image or in words, and thereby to
make known its value for oneself and for others. We live or feel ourselves into
the subject-matter, strive to follow it in its fate and to get a share in its
properties. The individual and immediately intuitable, which no science can
exhaust in its concepts or in its laws, poetry gives us in its images, sprung from
lived experience and living-into --- and perhaps even with so complete an
individualization as the subject-matters given in experience did not have. And
what thought can express in pale forms --- the opposition between the low and the
high, the bitter and the sweet, the fleeting and the abiding --- poetry, especially
as drama, leads us immediately into. Whether what is given us in this way is the
rule or the exception, poetry cannot itself decide. Every example is a symbol, but
not every symbol is an example. (Cf.\ above, § 10.) Poetry can set problems only
in so far as the very fact that a single, unique, individual subject-matter belongs
to existence can serve to characterize existence, or can at least give occasion for
serious reflection. But often the emotion awakened by the subject-matter will be so
strong that no single symbol satisfies, and a whole series of symbols is developed.
It is a series-formation of another kind than the intellectual one. (See § 5.) Not
the inner, logical relation between the members of the series, but the emotion
underlying all the members, here conditions the continuity and binds the members
together into a series.

It marks, as \textsc{Schiller} already saw, an advance of culture when there can
be joy in the image, apart from whether the subject-matter is in our possession or
not, and apart from whether it has any significance as a sign or as a means. So
long as ``participation'' (``realism of the whole'') rules, this joy cannot arise.
But even when the image is freed from practical need of one kind or another, the
joy in it does not thereby become of a purely intellectual kind. It is a
displacement that has taken place, and with the displacement there follows an
ennobling; the interest has moved from the subject-matter's reality to the image's
ideality, and has thereby itself undergone a metamorphosis.

A distinction must here be drawn between the mind that gives birth to the image
and the mind that is affected by the image once it has come into the world.

In the poet, feeling is the primary thing. The poetic images or similitudes
proceed from a feeling of the subject-matter's value, which awakens a need to
dwell on it and to see it from all sides. Perhaps the mind gets air only through a
whole series of similitudes. These then present analogies among themselves, but
these analogies are due to the inwardness with which the poet's mind holds fast to
the subject-matter and its value. Thus there is analogy between Jesus' parables
among themselves, because they are borne by one and the same fundamental mood ---
that awakened by the idea of the kingdom of heaven.

The reader or hearer is here otherwise placed than the poet himself. The image or
the similitude shows him, perhaps with striking clarity and in an expression of
the whole, the relation that is in question. And in relation to this intellectual
effect, the feeling becomes derived. The poet need not have set himself the express
purpose of bringing forth this effect. In the moment of inspiration he in any case
does not. He sings as the bird sings.

In his interesting treatise \textit{Metafor och Fiktion. Ett bidrag til Poetikens
Teori} (Lund 1922), Dr.\ \textsc{Nyman} relies on the statements of experimental
subjects to whom different poetic images were presented. Here it follows of itself
that the feeling becomes secondary, and that the image's effect in calling forth an
intellectual bird's-eye view comes to stand in the first rank. But for a full
aesthetic appreciation of the significance of images and similitudes, one must ask
both after the origin and after the effect. And the poet one cannot make into an
experimental subject, so as by that road to find out how the image comes into
being. He must be spied on. And \textsc{Julius Lange} has shown himself to have
spied well, when in his definition of artistic value he set as the first member
``the value a subject-matter has for the artist''. In return, he laid too little
weight on the ways and means by which that value could be called forth for hearer
or reader. (Cf.\ above, § 4 b.)

All poetry stands in connection with life and its fates, and it may therefore
become difficult to distinguish between poetry and religion, to decide where the
one ends and the other begins, or conversely. Religion too speaks in images and
symbols. The distinction that can be drawn rests essentially on this, that the
shaping of the images \textit{as} images plays a greater part for poetry than for
religion. The fact is, no doubt, that the underlying experiences are here deeper
and more strongly enthralling, and address themselves, with the demands that spring
from them, more to the will than to the imagination. The need to live immediately
in the symbols, and thereby to get a share in the experiences from which the
symbols sprang when they were formed at first hand, is characteristically
religious; and on that rests the great part tradition often plays in religion. It
is not only because tradition in and for itself constantly exercises its power over
men's minds that dogmas and forms of worship maintain themselves so long. It is
also because there the individual stands before symbols in whose coming into being
he knows he has had no part, so that image and thing come to stand as one, and
immediate living-into becomes more easily possible. As little as the individual is
conscious of having had a part in the symbol's formation, so little does he
distinguish between interpretation and lived experience, between image and reality.
The poet complains that ``life hovers between vision and possession''; but this
difference does not exist at pronouncedly religious stages. There will here in any
case be a point where vision and possession coincide. On this it rests that
religion bears the stamp of the primitive standpoint of participation, and of its
involuntary analogies, longer than science and poetry do.

The divorce between lived experience and interpretation, between possession and
vision, scarcely goes on without pain. The religious problem, as it presents itself
in our own day, concerns essentially the question whether this divorce can be
carried through without harm to the energy and concentration of spiritual life.

For in lived experience we are receptive; life streams in over us with its mighty
waves. In interpretation we are more self-active; it is a damming in and a
limiting, a self-defence against the strong influences. The old Frisians said that
God made the sea, but the Frisians the dykes. (\textit{Deus mare, Frisones littora
fecerunt}.) The damming in may be necessary to preserve the fertility of the soil
on which life is to be lived. But here individual forces intervene, and there may
even come a tension between lived experience and interpretation. So long as
tradition gives us both at once --- or rather, so long as it is not seen that
tradition itself contains an interpretation --- no properly religious problem arises.
There can at most arise the question how what has been handed down is to be
appropriated and understood.

But the real problem religion has in common with cosmology. In every symbol a part
serves to elucidate the whole. In the religious symbols too, a part or side of
existence is to elucidate, indeed to reveal, the whole of existence's inmost
essence. Men cling, in all forms of religion or religiosity, to a single
subject-matter in the confidence that it marks a breaking through of the highest
thing of value. The choice stands between this determinate interpretation and the
possibility of a chaos of valuelessness and meaninglessness. In \textsc{Pascal} one
can study this Either---Or particularly, which forms an analogy to the logical
opposition between a chaotic difference-series and a partial or absolute
identity-series. There was for \textsc{Pascal} only one single gleam of light in
the darkness of existence. The profound \textsc{Jacob Böhme} too, deeply seized by
the great riddles, halted for a while at this Either---Or. But he seized the way out
of facing the oppositions and apprehending existence as a great battlefield where
the strife stood between light and darkness, purity and lust, love and hate; and he
found here the foundation for a bold ethical conduct of life, while \textsc{Pascal}'s
dilemma led him to strict and solitary asceticism as the only way out.

When a part is to stand as the expression of the whole, there will involuntarily
make itself felt a striving to prove that the part in question really is typical
--- perhaps even that all other subject-matters that experience presents can be
derived from it. Thereby religion becomes theology and speculation, and sophistries
lie near. As we have seen, an example is always a symbol, but a symbol not always an
example. What becomes a symbol may be a \textit{unicum}. The religious poet says of
his God that he ``saves the trembling dove''. But are there not many poor doves that
are not saved? And if they were all saved, how would it then go with the goshawks,
which are after all subject-matters too? Nor is there any lack of attempts to take
the hawks for symbols. --- Religious symbolism at its height will regard the
subject-matter it originally started out from as derived in relation to the
generalization or absolutizing of that subject-matter to which it involuntarily
leads. The father-symbol is originally fetched from the fatherly love that can be
experienced in human relations. But then the relation is turned round, and the fact
that fatherly love can be practised in human relations is derived from God's being
Father (Eph.\ III, 15). (Cf.\ above, § 5.) --- Or it is a single historical figure
which, by purity of mind and by its love for men, makes vivid the wish to see in it
a revelation of the world's inmost ground. But not only do there stand over against
such a figure types of a quite different kind. There also arises the question
whether any single figure at all is able to unite in itself all the valuable
properties which, through the changing ages and cultures, must be presupposed in one
who is to be the constant model for striving men.

\textsc{Stanley Hall} has in a remarkable and richly contentful book (Jesus Christ
in the light of Psychology. 1921) made an attempt to show that the figure of Jesus,
when apprehended with historical criticism and with a thorough psychological sense,
is and remains the great model for human life, under all life's fates and in all
life's striving. \textsc{Stanley Hall} thinks that by psychology's help he has got
beyond the opposition between orthodoxy and criticism, and in particular that he has
overcome the difficulty which arises when the model is to give guidance in the face
of interests, tasks and experiences which did not, and could not, lie before the
historical figure that has been made into the model. The Jesus of the ``new age''
whom \textsc{Stanley Hall} would present, and for whose presentation he uses all the
resources of the newer psychology, is to unite in himself all great ideals, no
matter where they may come from. He is clear that between these ideals, proceeding
from different regions of life and history, there will be oppositions so great that
a complete harmony will be impossible; but he maintains that it is not possible in
any personality at all.\footnote{He will come to seem a baffling paradox, until it
is understood that personality means a synthesis of elements too manifold and
diverse ever to be completely harmonized (p.\ 238).} But the question is whether a
new image of personality is really brought forth by combining the ideals of many
different ages and cultures. And even if it were possible, the new model would all
the same become so essentially different from the one that has worked on men's minds
through centuries, that continuity with the historical figure which is still to lend
its name would be broken.\footnote{\textsc{Stanley Hall}'s predecessors in
religious-psychological research within American literature --- especially
\textsc{Starbuck, Coe} and \textsc{Leuba} --- strongly emphasized the one-sidedness
of the traditional shaping of the image of Jesus. And \textsc{Coe} in particular also
lays much weight on what is questionable in regarding a single type as alone
warranted. See my treatise \textit{Om nogle religionsfilosofiske Arbejder fra den
nyeste Tid.} 1903. (Mindre Arbejder. Anden Række).}

Neither in models nor in dogmas can the continuity of human spiritual life be
sought, but in the deep need which, in reciprocal relation with life's course, its
gifts and its tasks, its heights and its hollows, gives rise to the changing symbols
--- between which, therefore, analogy but not identity can take place. The analogy
makes it possible to learn from models and dogmas other than our own. But
identification is impossible, since it conflicts with life's nature and conditions.

The more religion becomes a poetry of life without dogmatic tendency, the less will
it assert that the sources of value which it hopes will continue to flow, however
much resistance they may meet, are at the same time the sources from which all the
subject-matters given in our experience spring. It will be enough for it that a
stream of value constantly goes through the world, and it will awaken a striving not
only to get a share in what this stream carries with it, but also to clear away the
hindrances to its course, so far as it stands in men's power. ---

In my \textit{Religionsfilosofi} these thoughts are more closely developed. The
purpose of this concluding part of the treatise has only been to intimate that the
emotional analogies which make themselves felt in practical and religious symbolism
are peculiar counterparts --- though in part also opposites --- to the rational
analogies which are of such great significance in science. Both kinds of analogy
stand more and more in opposition to the analogy that lies hidden in primitive
``participation''.

\end{document}
